% BAGEL 在 Intel Arc Pro B70 上的 M* 与 vLLM-Omni 对比：框架支持范围、模型接入工作量、本机实测（16:9 幻灯片）。
% 版式、配色与宏沿用 share/omni_runtime_report_v2（common.tex / deck.tex）。
% 数据：scripts/make_tex_tables.py 从 results/*/timings.csv 生成 generated/*.tex。编译：bash build.sh
\PassOptionsToPackage{table}{xcolor}
\documentclass[aspectratio=169,10pt]{beamer}

% 两份文档共用的部分：字体、配色、TikZ 样式、表格与代码块、标签与内容宏。
% 不含页面版式：幻灯片外观在 deck.tex，两份文档的标题与内容在 main.tex、main_framework.tex。
% 使用前需要先在调用文件里写 \PassOptionsToPackage{table}{xcolor} 和 \documentclass。

% ---------- 字体：Lato（拉丁）+ Noto Sans CJK SC（中文），等宽用 DejaVu Sans Mono ----------
\usepackage[UTF8,scheme=plain]{ctex}
% 注意：fontconfig 会把名称 "Lato" 解析成 Hairline 字重，所以直接按文件加载；缺文件时退回 Noto Sans CJK SC。
\IfFileExists{/usr/share/fonts/truetype/lato/Lato-Regular.ttf}{%
  \setmainfont{Lato-Regular.ttf}[Path=/usr/share/fonts/truetype/lato/,BoldFont=Lato-Bold.ttf,ItalicFont=Lato-Italic.ttf,BoldItalicFont=Lato-BoldItalic.ttf]%
  \setsansfont{Lato-Regular.ttf}[Path=/usr/share/fonts/truetype/lato/,BoldFont=Lato-Bold.ttf,ItalicFont=Lato-Italic.ttf,BoldItalicFont=Lato-BoldItalic.ttf]%
}{%
  \setmainfont{Noto Sans CJK SC}\setsansfont{Noto Sans CJK SC}%
}
\setmonofont{DejaVu Sans Mono}[Scale=0.84]
\setCJKmainfont{Noto Sans CJK SC}
\setCJKsansfont{Noto Sans CJK SC}
\setCJKmonofont{Noto Sans CJK SC}
\renewcommand{\familydefault}{\sfdefault}

\usepackage{booktabs,array,tabularx,multirow,graphicx,xcolor}
\usepackage{tikz}
\usetikzlibrary{arrows.meta,positioning,calc,fit,shapes.geometric,backgrounds,decorations.pathreplacing}
\usepackage{fancyvrb}
\usepackage[skins]{tcolorbox}
\usepackage{hyperref}

% ---------- 配色：取自公司模板（深蓝 001E50、蓝 0054AE / 0046C9 / 00377C、青 00C7FD、紫 6515AD / AF99FF、浅底 F5F8FF）----------
% 语义：Teal = M*（主蓝）、Orange = vLLM-Omni（紫）、Blue = SGLang-Omni（青）、Gold = 状态 / 交接（深蓝）、Red = 警示。
% 变量名沿用旧版（Teal / Orange / ...），只改取值，页面内容里的颜色引用不用动。
\definecolor{Ink}{HTML}{0A1A33}
\definecolor{Navy}{HTML}{001E50}
\definecolor{Teal}{HTML}{0054AE}
\definecolor{Blue}{HTML}{0092C2}
\definecolor{Orange}{HTML}{6515AD}
\definecolor{Red}{HTML}{C2334A}
\definecolor{Gold}{HTML}{00377C}
\definecolor{Cyan}{HTML}{00C7FD}
\definecolor{Lav}{HTML}{AF99FF}
\definecolor{Muted}{HTML}{5B6B82}
\definecolor{Line}{HTML}{C9D3E3}
\definecolor{Mist}{HTML}{F5F8FF}
\definecolor{TealBg}{HTML}{E6EFF9}
\definecolor{BlueBg}{HTML}{E0F6FD}
\definecolor{OrangeBg}{HTML}{EFE7F8}
\definecolor{RedBg}{HTML}{FBE9EC}
\definecolor{GoldBg}{HTML}{E4EAF4}
\hypersetup{colorlinks=true,linkcolor=Ink,urlcolor=Blue,citecolor=Blue}

\newcolumntype{Y}{>{\raggedright\arraybackslash}X}
% Pull top-aligned columns closer to the title rule without changing table-only slides.
\newcommand{\minipagetop}{\vspace{-.3cm}}

% ---------- 通用宏 ----------
\newcommand{\memo}[1]{{\color{Muted}\footnotesize #1}}
\newcommand{\pill}[2]{{\setlength{\fboxsep}{1.6pt}\raisebox{.05ex}{\colorbox{#1}{\textcolor{white}{\fontsize{6.5}{7}\selectfont\bfseries #2}}}}}
\newcommand{\tagE}{\pill{Teal}{实验}}
\newcommand{\tagC}{\pill{Blue}{源码}}
\newcommand{\tagO}{\pill{Gold}{官方}}
\newcommand{\tagH}{\pill{Orange}{假设}}
\newcommand{\tagG}{\pill{Red}{未测}}
\newcommand{\labSame}{\pill{Teal}{方向一致}}
\newcommand{\labMag}{\pill{Gold}{量级不同}}
\newcommand{\labBase}{\pill{Orange}{口径不同}}
\newcommand{\labRev}{\pill{Red}{方向相反}}
% 伪代码块：Verbatim 中用 \kw{} \cm{} \id{} 着色（花括号需写成 \{ \}）
\newcommand{\kw}[1]{\textbf{\textcolor{Teal}{#1}}}
\newcommand{\cm}[1]{\textcolor{Muted}{#1}}
\newcommand{\id}[1]{\textcolor{Orange}{#1}}
\newcommand{\hi}[1]{\textbf{\textcolor{Red}{#1}}}
\newtcolorbox{codebox}[2][]{enhanced,colback=Mist,colframe=Line,boxrule=.4pt,arc=2pt,left=4pt,right=3pt,top=2pt,bottom=2pt,
  title={#2},coltitle=white,colbacktitle=Teal,fonttitle=\fontsize{6.8}{8}\selectfont\bfseries,boxsep=1pt,toptitle=1.2pt,bottomtitle=1.2pt,lefttitle=4pt,before skip=0pt,#1}
% 右侧说明框：高度自适应，可连续堆叠
\newtcolorbox{nbox}[3][]{colback=#2!6,colframe=#2,boxrule=.7pt,arc=2.5pt,left=4pt,right=4pt,top=3pt,bottom=3pt,
  before upper={\def\codesize{\relax}\noindent\textbf{\textcolor{#2}{#3}}\par\vspace{1pt}},fontupper=\fontsize{7.7}{9.5}\selectfont\raggedright,before skip=0pt,after skip=4pt,#1}
\newcommand{\hl}[1]{\textbf{\textcolor{Teal}{#1}}}
\newcommand{\warn}[1]{\textbf{\textcolor{Red}{#1}}}
\newcommand{\codesize}{\small}
\newcommand{\code}[1]{\texttt{\codesize #1}}
\newcommand{\tb}[1]{{\bfseries #1}}
% Walk 名字：小字号，长名（如 image_gen_guided3）能放进 2.25 cm 的框
\newcommand{\wname}[1]{{\fontsize{6.6}{8}\selectfont\bfseries #1}}

% 图片来源注记
\newcommand{\figsource}[1]{{\color{Muted}\fontsize{6.5}{7.8}\selectfont #1}}

% ---------- TikZ 图元样式（固定高度、字号，避免重叠） ----------
\tikzset{
  >={Latex[length=1.8mm]},
  box/.style={draw=#1,fill=#1!8,rounded corners=2.5pt,line width=.8pt,align=center,font=\footnotesize,minimum height=.62cm,inner xsep=4pt,inner ysep=2.5pt},
  box/.default=Teal,
  sbox/.style={box=#1,font=\scriptsize,minimum height=.45cm,inner xsep=3pt,inner ysep=1.5pt},
  sbox/.default=Teal,
  plain/.style={align=center,font=\footnotesize,inner sep=1pt},
  tiny/.style={align=center,font=\scriptsize,inner sep=1pt,text=Muted},
  flow/.style={->,draw=Muted,line width=.9pt},
  ctrl/.style={->,draw=Blue,line width=.9pt,densely dashed},
  strm/.style={->,draw=Orange,line width=1.2pt},
  back/.style={->,draw=Teal,line width=1pt,rounded corners=3pt},
  wk/.style={box=#1,minimum height=.62cm,font=\scriptsize,inner xsep=3pt},
  panel/.style={draw=Line,fill=Mist,rounded corners=3pt,line width=.6pt},
}

% 对数轴上的点与连线（第 18 页）
\newcommand{\two}[3]{\pgfmathsetmacro{\xa}{ln(#2/0.5)/ln(2)*1.5}\pgfmathsetmacro{\xb}{ln(#3/0.5)/ln(2)*1.5}\draw[Muted!60,line width=.8pt] (\xa,#1)--(\xb,#1);}
\newcommand{\dotat}[3]{\pgfmathsetmacro{\xx}{ln(#2/0.5)/ln(2)*1.5}\node[#3] at (\xx,#1) {};}

% 幻灯片版式（两份汇报共用）：16:9 Beamer，白底，左对齐标题，底部固定“要点条”，右下页码。
% 调用前需要：\documentclass[aspectratio=169,10pt]{beamer}，然后 % 两份文档共用的部分：字体、配色、TikZ 样式、表格与代码块、标签与内容宏。
% 不含页面版式：幻灯片外观在 deck.tex，两份文档的标题与内容在 main.tex、main_framework.tex。
% 使用前需要先在调用文件里写 \PassOptionsToPackage{table}{xcolor} 和 \documentclass。

% ---------- 字体：Lato（拉丁）+ Noto Sans CJK SC（中文），等宽用 DejaVu Sans Mono ----------
\usepackage[UTF8,scheme=plain]{ctex}
% 注意：fontconfig 会把名称 "Lato" 解析成 Hairline 字重，所以直接按文件加载；缺文件时退回 Noto Sans CJK SC。
\IfFileExists{/usr/share/fonts/truetype/lato/Lato-Regular.ttf}{%
  \setmainfont{Lato-Regular.ttf}[Path=/usr/share/fonts/truetype/lato/,BoldFont=Lato-Bold.ttf,ItalicFont=Lato-Italic.ttf,BoldItalicFont=Lato-BoldItalic.ttf]%
  \setsansfont{Lato-Regular.ttf}[Path=/usr/share/fonts/truetype/lato/,BoldFont=Lato-Bold.ttf,ItalicFont=Lato-Italic.ttf,BoldItalicFont=Lato-BoldItalic.ttf]%
}{%
  \setmainfont{Noto Sans CJK SC}\setsansfont{Noto Sans CJK SC}%
}
\setmonofont{DejaVu Sans Mono}[Scale=0.84]
\setCJKmainfont{Noto Sans CJK SC}
\setCJKsansfont{Noto Sans CJK SC}
\setCJKmonofont{Noto Sans CJK SC}
\renewcommand{\familydefault}{\sfdefault}

\usepackage{booktabs,array,tabularx,multirow,graphicx,xcolor}
\usepackage{tikz}
\usetikzlibrary{arrows.meta,positioning,calc,fit,shapes.geometric,backgrounds,decorations.pathreplacing}
\usepackage{fancyvrb}
\usepackage[skins]{tcolorbox}
\usepackage{hyperref}

% ---------- 配色：取自公司模板（深蓝 001E50、蓝 0054AE / 0046C9 / 00377C、青 00C7FD、紫 6515AD / AF99FF、浅底 F5F8FF）----------
% 语义：Teal = M*（主蓝）、Orange = vLLM-Omni（紫）、Blue = SGLang-Omni（青）、Gold = 状态 / 交接（深蓝）、Red = 警示。
% 变量名沿用旧版（Teal / Orange / ...），只改取值，页面内容里的颜色引用不用动。
\definecolor{Ink}{HTML}{0A1A33}
\definecolor{Navy}{HTML}{001E50}
\definecolor{Teal}{HTML}{0054AE}
\definecolor{Blue}{HTML}{0092C2}
\definecolor{Orange}{HTML}{6515AD}
\definecolor{Red}{HTML}{C2334A}
\definecolor{Gold}{HTML}{00377C}
\definecolor{Cyan}{HTML}{00C7FD}
\definecolor{Lav}{HTML}{AF99FF}
\definecolor{Muted}{HTML}{5B6B82}
\definecolor{Line}{HTML}{C9D3E3}
\definecolor{Mist}{HTML}{F5F8FF}
\definecolor{TealBg}{HTML}{E6EFF9}
\definecolor{BlueBg}{HTML}{E0F6FD}
\definecolor{OrangeBg}{HTML}{EFE7F8}
\definecolor{RedBg}{HTML}{FBE9EC}
\definecolor{GoldBg}{HTML}{E4EAF4}
\hypersetup{colorlinks=true,linkcolor=Ink,urlcolor=Blue,citecolor=Blue}

\newcolumntype{Y}{>{\raggedright\arraybackslash}X}
% Pull top-aligned columns closer to the title rule without changing table-only slides.
\newcommand{\minipagetop}{\vspace{-.3cm}}

% ---------- 通用宏 ----------
\newcommand{\memo}[1]{{\color{Muted}\footnotesize #1}}
\newcommand{\pill}[2]{{\setlength{\fboxsep}{1.6pt}\raisebox{.05ex}{\colorbox{#1}{\textcolor{white}{\fontsize{6.5}{7}\selectfont\bfseries #2}}}}}
\newcommand{\tagE}{\pill{Teal}{实验}}
\newcommand{\tagC}{\pill{Blue}{源码}}
\newcommand{\tagO}{\pill{Gold}{官方}}
\newcommand{\tagH}{\pill{Orange}{假设}}
\newcommand{\tagG}{\pill{Red}{未测}}
\newcommand{\labSame}{\pill{Teal}{方向一致}}
\newcommand{\labMag}{\pill{Gold}{量级不同}}
\newcommand{\labBase}{\pill{Orange}{口径不同}}
\newcommand{\labRev}{\pill{Red}{方向相反}}
% 伪代码块：Verbatim 中用 \kw{} \cm{} \id{} 着色（花括号需写成 \{ \}）
\newcommand{\kw}[1]{\textbf{\textcolor{Teal}{#1}}}
\newcommand{\cm}[1]{\textcolor{Muted}{#1}}
\newcommand{\id}[1]{\textcolor{Orange}{#1}}
\newcommand{\hi}[1]{\textbf{\textcolor{Red}{#1}}}
\newtcolorbox{codebox}[2][]{enhanced,colback=Mist,colframe=Line,boxrule=.4pt,arc=2pt,left=4pt,right=3pt,top=2pt,bottom=2pt,
  title={#2},coltitle=white,colbacktitle=Teal,fonttitle=\fontsize{6.8}{8}\selectfont\bfseries,boxsep=1pt,toptitle=1.2pt,bottomtitle=1.2pt,lefttitle=4pt,before skip=0pt,#1}
% 右侧说明框：高度自适应，可连续堆叠
\newtcolorbox{nbox}[3][]{colback=#2!6,colframe=#2,boxrule=.7pt,arc=2.5pt,left=4pt,right=4pt,top=3pt,bottom=3pt,
  before upper={\def\codesize{\relax}\noindent\textbf{\textcolor{#2}{#3}}\par\vspace{1pt}},fontupper=\fontsize{7.7}{9.5}\selectfont\raggedright,before skip=0pt,after skip=4pt,#1}
\newcommand{\hl}[1]{\textbf{\textcolor{Teal}{#1}}}
\newcommand{\warn}[1]{\textbf{\textcolor{Red}{#1}}}
\newcommand{\codesize}{\small}
\newcommand{\code}[1]{\texttt{\codesize #1}}
\newcommand{\tb}[1]{{\bfseries #1}}
% Walk 名字：小字号，长名（如 image_gen_guided3）能放进 2.25 cm 的框
\newcommand{\wname}[1]{{\fontsize{6.6}{8}\selectfont\bfseries #1}}

% 图片来源注记
\newcommand{\figsource}[1]{{\color{Muted}\fontsize{6.5}{7.8}\selectfont #1}}

% ---------- TikZ 图元样式（固定高度、字号，避免重叠） ----------
\tikzset{
  >={Latex[length=1.8mm]},
  box/.style={draw=#1,fill=#1!8,rounded corners=2.5pt,line width=.8pt,align=center,font=\footnotesize,minimum height=.62cm,inner xsep=4pt,inner ysep=2.5pt},
  box/.default=Teal,
  sbox/.style={box=#1,font=\scriptsize,minimum height=.45cm,inner xsep=3pt,inner ysep=1.5pt},
  sbox/.default=Teal,
  plain/.style={align=center,font=\footnotesize,inner sep=1pt},
  tiny/.style={align=center,font=\scriptsize,inner sep=1pt,text=Muted},
  flow/.style={->,draw=Muted,line width=.9pt},
  ctrl/.style={->,draw=Blue,line width=.9pt,densely dashed},
  strm/.style={->,draw=Orange,line width=1.2pt},
  back/.style={->,draw=Teal,line width=1pt,rounded corners=3pt},
  wk/.style={box=#1,minimum height=.62cm,font=\scriptsize,inner xsep=3pt},
  panel/.style={draw=Line,fill=Mist,rounded corners=3pt,line width=.6pt},
}

% 对数轴上的点与连线（第 18 页）
\newcommand{\two}[3]{\pgfmathsetmacro{\xa}{ln(#2/0.5)/ln(2)*1.5}\pgfmathsetmacro{\xb}{ln(#3/0.5)/ln(2)*1.5}\draw[Muted!60,line width=.8pt] (\xa,#1)--(\xb,#1);}
\newcommand{\dotat}[3]{\pgfmathsetmacro{\xx}{ln(#2/0.5)/ln(2)*1.5}\node[#3] at (\xx,#1) {};}
。

\setbeamercolor{normal text}{fg=Ink,bg=white}
\setbeamercolor{structure}{fg=Teal}
\setbeamercolor{alerted text}{fg=Red}
\setbeamertemplate{navigation symbols}{}
\setbeamertemplate{itemize item}{\raisebox{.15ex}{\color{Teal}\scriptsize$\blacksquare$}}
\setbeamertemplate{itemize subitem}{\raisebox{.15ex}{\color{Teal}\tiny$\blacksquare$}}
\setbeamertemplate{enumerate item}{\color{Teal}\bfseries\insertenumlabel.}
\setbeamersize{text margin left=1.0cm,text margin right=1.0cm}
\setbeamerfont{itemize/enumerate body}{size=\small}
\setbeamerfont{itemize/enumerate subbody}{size=\footnotesize}
\setlength{\leftmargini}{1.1em}
\setlength{\parskip}{0pt}
\setlength{\itemsep}{2pt}
\renewcommand{\arraystretch}{1.18}
\setlength{\tabcolsep}{4pt}

\newcommand{\deckpart}{}
% PDF 书签：每个 \setpart 是一级书签，每页标题是二级书签（打开 PDF 时可从侧边栏跳转）。
\hypersetup{bookmarks=true,bookmarksopen=true,bookmarksopenlevel=1,bookmarksnumbered=false,unicode=true}
\newcommand{\setpart}[1]{\renewcommand{\deckpart}{#1}\pdfbookmark[1]{#1}{part.\thepage.\the\numexpr\value{page}\relax}}
\newcommand{\framebookmark}[1]{\pdfbookmark[2]{\thepage\ \ #1}{frame.\thepage}}

\setbeamertemplate{frametitle}{%
  \vspace{-.05cm}%
  \framebookmark{\insertframetitle}%
  \begin{tikzpicture}[remember picture,overlay]
    \fill[Cyan] ([xshift=1.0cm,yshift=-.43cm]current page.north west) rectangle ++(.13cm,-.13cm);
    \node[anchor=west,inner sep=0,font=\fontsize{7}{8}\selectfont\bfseries,text=Muted] at ([xshift=1.25cm,yshift=-.495cm]current page.north west) {\deckpart};
    \node[anchor=north west,inner sep=0,font=\fontsize{14.5}{17}\selectfont\bfseries,text=Ink,text width=14.4cm] at ([xshift=1.0cm,yshift=-.7cm]current page.north west) {\insertframetitle};
    \fill[Cyan] ([xshift=1.0cm,yshift=-1.42cm]current page.north west) rectangle ++(1.0cm,-.045cm);
    \fill[Orange] ([xshift=2.0cm,yshift=-1.42cm]current page.north west) rectangle ++(1.0cm,-.045cm);
    \fill[Teal] ([xshift=3.0cm,yshift=-1.42cm]current page.north west) rectangle ++(1.0cm,-.045cm);
  \end{tikzpicture}%
  \vspace{.85cm}%
}
% 页脚：只保留页码，右下对齐。
\setbeamertemplate{footline}{%
  \begin{tikzpicture}[remember picture,overlay]
    \node[anchor=east,inner sep=0,font=\fontsize{7.5}{9}\selectfont\bfseries,text=Muted] at ([xshift=-1.0cm,yshift=.36cm]current page.south east) {\insertframenumber};
  \end{tikzpicture}%
}

% 页尾要点条：没有标签、不是每页都放；位置贴近页脚，内容区因此更大。
\newcommand{\takeaway}[1]{%
  \begin{tikzpicture}[remember picture,overlay]
    \node[anchor=south west,inner xsep=.38cm,inner ysep=.1cm,minimum height=.74cm,minimum width=14cm,text width=13.3cm,align=left,fill=Mist,font=\fontsize{8.6}{11}\selectfont] (sw) at ([xshift=1.0cm,yshift=.62cm]current page.south west) {#1};
    \fill[Teal] (sw.south west) rectangle ([xshift=.09cm]sw.north west);
  \end{tikzpicture}}

% 封面：白底、大标题、左下蓝色标签、细横线。调用时传入标题、副标题、标签行。
\newcommand{\deckcover}[4]{%
  \begin{frame}[plain,noframenumbering]
    \pdfbookmark[1]{封面}{title.page}%
    \begin{tikzpicture}[remember picture,overlay]
      \fill[Cyan] ([xshift=1.4cm,yshift=-1.35cm]current page.north west) rectangle ++(.2cm,-.2cm);
      \fill[Teal] ([xshift=1.4cm,yshift=-1.55cm]current page.north west) rectangle ++(2.4cm,-.16cm);
      \node[anchor=west,inner sep=0,text width=13.6cm,align=left] at ([xshift=1.4cm,yshift=.35cm]current page.west) {%
        {\fontsize{23}{30}\selectfont\bfseries\textcolor{Navy}{#1}}\\[.5cm]
        {\fontsize{12}{16}\selectfont\textcolor{Teal}{\bfseries #2}}\\[.1cm]
        {\fontsize{10.5}{14}\selectfont\textcolor{Muted}{#3}}};
      \draw[Line,line width=.5pt] ([yshift=1.05cm]current page.south west) -- ([yshift=1.05cm]current page.south east);
      \fill[Teal] ([yshift=.0cm]current page.south west) rectangle ++(2.6cm,1.05cm);
      \node[anchor=center,font=\footnotesize\bfseries,text=white] at ([xshift=1.3cm,yshift=.52cm]current page.south west) {2026-11-22};
      \node[anchor=west,inner sep=0,font=\fontsize{7}{9}\selectfont,text=Muted,align=left] at ([xshift=3.2cm,yshift=.52cm]current page.south west) {#4};
    \end{tikzpicture}
  \end{frame}}

% generated by scripts/make_tex_tables.py -- do not edit
\makeatletter
\newcommand{\val}[1]{\ifcsname val@#1\endcsname\csname val@#1\endcsname\else--\fi}
\newcommand{\raw}[1]{\ifcsname raw@#1\endcsname\csname raw@#1\endcsname\else--\fi}
\newcommand{\anyval}[1]{\ifcsname any@#1\endcsname\csname any@#1\endcsname\else--\fi}
\newcommand{\warm}[1]{\ifcsname warm@#1\endcsname\csname warm@#1\endcsname\else--\fi}
\newcommand{\pct}[1]{\ifcsname pct@#1\endcsname\csname pct@#1\endcsname\else n/a\fi}
\expandafter\def\csname any@final/mstar-cfgpar/t2i\endcsname{7.678}
\expandafter\def\csname val@final/mstar-cfgpar/t2i\endcsname{7.68}
\expandafter\def\csname raw@final/mstar-cfgpar/t2i\endcsname{7.678}
\expandafter\def\csname warm@final/mstar-cfgpar/t2i\endcsname{17.5}
\expandafter\def\csname any@final/mstar-cfgpar/i2t1\endcsname{0.085}
\expandafter\def\csname val@final/mstar-cfgpar/i2t1\endcsname{0.08}
\expandafter\def\csname raw@final/mstar-cfgpar/i2t1\endcsname{0.085}
\expandafter\def\csname warm@final/mstar-cfgpar/i2t1\endcsname{8.4}
\expandafter\def\csname any@final/mstar-cfgpar/i2t64\endcsname{2.458}
\expandafter\def\csname val@final/mstar-cfgpar/i2t64\endcsname{2.46}
\expandafter\def\csname raw@final/mstar-cfgpar/i2t64\endcsname{2.458}
\expandafter\def\csname warm@final/mstar-cfgpar/i2t64\endcsname{3.5}
\expandafter\def\csname any@final/mstar-cfgpar/i2i\endcsname{16.792}
\expandafter\def\csname val@final/mstar-cfgpar/i2i\endcsname{16.79}
\expandafter\def\csname raw@final/mstar-cfgpar/i2i\endcsname{16.792}
\expandafter\def\csname warm@final/mstar-cfgpar/i2i\endcsname{34.9}
\expandafter\def\csname any@final/mstar-cfgpar-nocompile/t2i\endcsname{8.476}
\expandafter\def\csname val@final/mstar-cfgpar-nocompile/t2i\endcsname{8.48}
\expandafter\def\csname raw@final/mstar-cfgpar-nocompile/t2i\endcsname{8.476}
\expandafter\def\csname warm@final/mstar-cfgpar-nocompile/t2i\endcsname{10.9}
\expandafter\def\csname any@final/mstar-cfgpar-nocompile/i2t1\endcsname{0.083}
\expandafter\def\csname val@final/mstar-cfgpar-nocompile/i2t1\endcsname{0.08}
\expandafter\def\csname raw@final/mstar-cfgpar-nocompile/i2t1\endcsname{0.083}
\expandafter\def\csname warm@final/mstar-cfgpar-nocompile/i2t1\endcsname{1.6}
\expandafter\def\csname any@final/mstar-cfgpar-nocompile/i2t64\endcsname{2.169}
\expandafter\def\csname val@final/mstar-cfgpar-nocompile/i2t64\endcsname{2.17}
\expandafter\def\csname raw@final/mstar-cfgpar-nocompile/i2t64\endcsname{2.169}
\expandafter\def\csname warm@final/mstar-cfgpar-nocompile/i2t64\endcsname{2.2}
\expandafter\def\csname any@final/mstar-cfgpar-nocompile/i2i\endcsname{18.251}
\expandafter\def\csname val@final/mstar-cfgpar-nocompile/i2i\endcsname{18.25}
\expandafter\def\csname raw@final/mstar-cfgpar-nocompile/i2i\endcsname{18.251}
\expandafter\def\csname warm@final/mstar-cfgpar-nocompile/i2i\endcsname{22.0}
\expandafter\def\csname any@final/mstar-tp2/t2i\endcsname{12.955}
\expandafter\def\csname val@final/mstar-tp2/t2i\endcsname{12.96}
\expandafter\def\csname raw@final/mstar-tp2/t2i\endcsname{12.955}
\expandafter\def\csname warm@final/mstar-tp2/t2i\endcsname{24.6}
\expandafter\def\csname any@final/mstar-tp2/i2t1\endcsname{0.330}
\expandafter\def\csname val@final/mstar-tp2/i2t1\endcsname{0.33}
\expandafter\def\csname raw@final/mstar-tp2/i2t1\endcsname{0.330}
\expandafter\def\csname warm@final/mstar-tp2/i2t1\endcsname{16.7}
\expandafter\def\csname any@final/mstar-tp2/i2t64\endcsname{3.156}
\expandafter\def\csname val@final/mstar-tp2/i2t64\endcsname{3.16}
\expandafter\def\csname raw@final/mstar-tp2/i2t64\endcsname{3.156}
\expandafter\def\csname warm@final/mstar-tp2/i2t64\endcsname{4.3}
\expandafter\def\csname any@final/mstar-tp2/i2i\endcsname{21.208}
\expandafter\def\csname val@final/mstar-tp2/i2i\endcsname{21.21}
\expandafter\def\csname raw@final/mstar-tp2/i2i\endcsname{21.208}
\expandafter\def\csname warm@final/mstar-tp2/i2i\endcsname{39.7}
\expandafter\def\csname any@final/omni-single-cfgpar2/t2i\endcsname{9.978}
\expandafter\def\csname val@final/omni-single-cfgpar2/t2i\endcsname{9.98}
\expandafter\def\csname raw@final/omni-single-cfgpar2/t2i\endcsname{9.978}
\expandafter\def\csname warm@final/omni-single-cfgpar2/t2i\endcsname{10.9}
\expandafter\def\csname any@final/omni-single-tp2/t2i\endcsname{17.040}
\expandafter\def\csname val@final/omni-single-tp2/t2i\endcsname{17.04}
\expandafter\def\csname raw@final/omni-single-tp2/t2i\endcsname{17.040}
\expandafter\def\csname warm@final/omni-single-tp2/t2i\endcsname{17.7}
\expandafter\def\csname any@final/omni-single-tp2/i2i\endcsname{28.076}
\expandafter\def\csname val@final/omni-single-tp2/i2i\endcsname{28.08}
\expandafter\def\csname raw@final/omni-single-tp2/i2i\endcsname{28.076}
\expandafter\def\csname warm@final/omni-single-tp2/i2i\endcsname{29.5}
\expandafter\def\csname any@final/omni-two-stage/t2i\endcsname{17.987}
\expandafter\def\csname val@final/omni-two-stage/t2i\endcsname{17.99}
\expandafter\def\csname raw@final/omni-two-stage/t2i\endcsname{17.987}
\expandafter\def\csname warm@final/omni-two-stage/t2i\endcsname{18.7}
\expandafter\def\csname any@final/omni-two-stage/i2t1\endcsname{0.193}
\expandafter\def\csname val@final/omni-two-stage/i2t1\endcsname{0.19}
\expandafter\def\csname raw@final/omni-two-stage/i2t1\endcsname{0.193}
\expandafter\def\csname warm@final/omni-two-stage/i2t1\endcsname{0.6}
\expandafter\def\csname any@final/omni-two-stage/i2t64\endcsname{1.882}
\expandafter\def\csname val@final/omni-two-stage/i2t64\endcsname{1.88}
\expandafter\def\csname raw@final/omni-two-stage/i2t64\endcsname{1.882}
\expandafter\def\csname warm@final/omni-two-stage/i2t64\endcsname{1.9}
\expandafter\def\csname any@final/omni-two-stage/i2i\endcsname{30.760}
\expandafter\def\csname any@final/omni-two-stage-graph/i2t1\endcsname{0.191}
\expandafter\def\csname val@final/omni-two-stage-graph/i2t1\endcsname{0.19}
\expandafter\def\csname raw@final/omni-two-stage-graph/i2t1\endcsname{0.191}
\expandafter\def\csname warm@final/omni-two-stage-graph/i2t1\endcsname{0.4}
\expandafter\def\csname any@final/omni-two-stage-graph/i2t64\endcsname{1.861}
\expandafter\def\csname val@final/omni-two-stage-graph/i2t64\endcsname{1.86}
\expandafter\def\csname raw@final/omni-two-stage-graph/i2t64\endcsname{1.861}
\expandafter\def\csname warm@final/omni-two-stage-graph/i2t64\endcsname{1.9}
\expandafter\def\csname any@screen/mstar-cfgpar/t2i\endcsname{7.702}
\expandafter\def\csname val@screen/mstar-cfgpar/t2i\endcsname{7.70}
\expandafter\def\csname raw@screen/mstar-cfgpar/t2i\endcsname{7.702}
\expandafter\def\csname warm@screen/mstar-cfgpar/t2i\endcsname{37.6}
\expandafter\def\csname any@screen/mstar-cfgpar/i2t64\endcsname{2.582}
\expandafter\def\csname any@screen/mstar-tp2/t2i\endcsname{12.946}
\expandafter\def\csname val@screen/mstar-tp2/t2i\endcsname{12.95}
\expandafter\def\csname raw@screen/mstar-tp2/t2i\endcsname{12.946}
\expandafter\def\csname warm@screen/mstar-tp2/t2i\endcsname{24.7}
\expandafter\def\csname any@screen/mstar-tp2/i2t64\endcsname{3.746}
\expandafter\def\csname any@screen/mstar-tp2-graph/t2i\endcsname{12.964}
\expandafter\def\csname val@screen/mstar-tp2-graph/t2i\endcsname{12.96}
\expandafter\def\csname raw@screen/mstar-tp2-graph/t2i\endcsname{12.964}
\expandafter\def\csname warm@screen/mstar-tp2-graph/t2i\endcsname{21.7}
\expandafter\def\csname any@screen/mstar-tp2-graph/i2t64\endcsname{3.790}
\expandafter\def\csname any@screen/omni-single-cfgpar2/t2i\endcsname{9.970}
\expandafter\def\csname val@screen/omni-single-cfgpar2/t2i\endcsname{9.97}
\expandafter\def\csname raw@screen/omni-single-cfgpar2/t2i\endcsname{9.970}
\expandafter\def\csname warm@screen/omni-single-cfgpar2/t2i\endcsname{11.0}
\expandafter\def\csname any@screen/omni-single-tp2/t2i\endcsname{17.060}
\expandafter\def\csname val@screen/omni-single-tp2/t2i\endcsname{17.06}
\expandafter\def\csname raw@screen/omni-single-tp2/t2i\endcsname{17.060}
\expandafter\def\csname warm@screen/omni-single-tp2/t2i\endcsname{17.8}
\expandafter\def\csname any@screen/omni-two-stage/t2i\endcsname{17.941}
\expandafter\def\csname val@screen/omni-two-stage/t2i\endcsname{17.94}
\expandafter\def\csname raw@screen/omni-two-stage/t2i\endcsname{17.941}
\expandafter\def\csname warm@screen/omni-two-stage/t2i\endcsname{30.8}
\expandafter\def\csname any@screen/omni-two-stage/i2t64\endcsname{1.881}
\expandafter\def\csname val@screen/omni-two-stage/i2t64\endcsname{1.88}
\expandafter\def\csname raw@screen/omni-two-stage/i2t64\endcsname{1.881}
\expandafter\def\csname warm@screen/omni-two-stage/i2t64\endcsname{2.2}
\expandafter\def\csname any@screen/omni-two-stage-graph/t2i\endcsname{17.970}
\expandafter\def\csname val@screen/omni-two-stage-graph/t2i\endcsname{17.97}
\expandafter\def\csname raw@screen/omni-two-stage-graph/t2i\endcsname{17.970}
\expandafter\def\csname warm@screen/omni-two-stage-graph/t2i\endcsname{18.4}
\expandafter\def\csname any@screen/omni-two-stage-graph/i2t64\endcsname{1.860}
\expandafter\def\csname val@screen/omni-two-stage-graph/i2t64\endcsname{1.86}
\expandafter\def\csname raw@screen/omni-two-stage-graph/i2t64\endcsname{1.860}
\expandafter\def\csname warm@screen/omni-two-stage-graph/i2t64\endcsname{2.1}
\expandafter\def\csname any@v032/diag-mstar-eager-mathsdpa/t2i\endcsname{13.901}
\expandafter\def\csname val@v032/diag-mstar-eager-mathsdpa/t2i\endcsname{13.90}
\expandafter\def\csname raw@v032/diag-mstar-eager-mathsdpa/t2i\endcsname{13.901}
\expandafter\def\csname warm@v032/diag-mstar-eager-mathsdpa/t2i\endcsname{17.2}
\expandafter\def\csname any@v032/diag-mstar-mathsdpa/t2i\endcsname{12.957}
\expandafter\def\csname val@v032/diag-mstar-mathsdpa/t2i\endcsname{12.96}
\expandafter\def\csname raw@v032/diag-mstar-mathsdpa/t2i\endcsname{12.957}
\expandafter\def\csname warm@v032/diag-mstar-mathsdpa/t2i\endcsname{41.7}
\expandafter\def\csname any@v032/diag-mstar-mathsdpa/i2t1\endcsname{1.445}
\expandafter\def\csname any@v032/diag-omni-single-mathsdpa/t2i\endcsname{17.660}
\expandafter\def\csname val@v032/diag-omni-single-mathsdpa/t2i\endcsname{17.66}
\expandafter\def\csname raw@v032/diag-omni-single-mathsdpa/t2i\endcsname{17.660}
\expandafter\def\csname warm@v032/diag-omni-single-mathsdpa/t2i\endcsname{18.2}
\expandafter\def\csname any@v032/diag-omni-single-mathsdpa/i2t1\endcsname{2.041}
\expandafter\def\csname val@v032/diag-omni-single-mathsdpa/i2t1\endcsname{2.04}
\expandafter\def\csname raw@v032/diag-omni-single-mathsdpa/i2t1\endcsname{2.041}
\expandafter\def\csname warm@v032/diag-omni-single-mathsdpa/i2t1\endcsname{3.1}
\expandafter\def\csname any@v032/diag-omni-single-no-ocl-fix/t2i\endcsname{17.651}
\expandafter\def\csname val@v032/diag-omni-single-no-ocl-fix/t2i\endcsname{17.65}
\expandafter\def\csname raw@v032/diag-omni-single-no-ocl-fix/t2i\endcsname{17.651}
\expandafter\def\csname warm@v032/diag-omni-single-no-ocl-fix/t2i\endcsname{18.2}
\expandafter\def\csname any@v032/diag-omni-single-no-ocl-fix/i2t1\endcsname{1.376}
\expandafter\def\csname val@v032/diag-omni-single-no-ocl-fix/i2t1\endcsname{1.38}
\expandafter\def\csname raw@v032/diag-omni-single-no-ocl-fix/i2t1\endcsname{1.376}
\expandafter\def\csname warm@v032/diag-omni-single-no-ocl-fix/i2t1\endcsname{2.3}
\expandafter\def\csname any@v032/mstar-tp2/t2i\endcsname{12.954}
\expandafter\def\csname val@v032/mstar-tp2/t2i\endcsname{12.95}
\expandafter\def\csname raw@v032/mstar-tp2/t2i\endcsname{12.954}
\expandafter\def\csname warm@v032/mstar-tp2/t2i\endcsname{24.7}
\expandafter\def\csname any@v032/mstar-tp2/i2t1\endcsname{0.962}
\expandafter\def\csname any@v032/mstar-tp2/i2t64\endcsname{3.754}
\expandafter\def\csname any@v032/mstar-tp2/i2i\endcsname{22.817}
\expandafter\def\csname val@v032/mstar-tp2/i2i\endcsname{22.82}
\expandafter\def\csname raw@v032/mstar-tp2/i2i\endcsname{22.817}
\expandafter\def\csname warm@v032/mstar-tp2/i2i\endcsname{52.5}
\expandafter\def\csname any@v032/mstar-tp2-nocompile/t2i\endcsname{13.905}
\expandafter\def\csname val@v032/mstar-tp2-nocompile/t2i\endcsname{13.91}
\expandafter\def\csname raw@v032/mstar-tp2-nocompile/t2i\endcsname{13.905}
\expandafter\def\csname warm@v032/mstar-tp2-nocompile/t2i\endcsname{17.1}
\expandafter\def\csname any@v032/mstar-tp2-nocompile/i2t1\endcsname{0.947}
\expandafter\def\csname any@v032/mstar-tp2-nocompile/i2t64\endcsname{3.040}
\expandafter\def\csname any@v032/mstar-tp2-nocompile/i2i\endcsname{24.666}
\expandafter\def\csname val@v032/mstar-tp2-nocompile/i2i\endcsname{24.67}
\expandafter\def\csname raw@v032/mstar-tp2-nocompile/i2i\endcsname{24.666}
\expandafter\def\csname warm@v032/mstar-tp2-nocompile/i2i\endcsname{25.3}
\expandafter\def\csname any@v032/omni-single-tp2/t2i\endcsname{17.656}
\expandafter\def\csname val@v032/omni-single-tp2/t2i\endcsname{17.66}
\expandafter\def\csname raw@v032/omni-single-tp2/t2i\endcsname{17.656}
\expandafter\def\csname warm@v032/omni-single-tp2/t2i\endcsname{18.2}
\expandafter\def\csname any@v032/omni-single-tp2/i2t1\endcsname{1.376}
\expandafter\def\csname any@v032/omni-single-tp2/i2t64\endcsname{3.423}
\expandafter\def\csname any@v032/omni-single-tp2/i2i\endcsname{30.740}
\expandafter\def\csname val@v032/omni-single-tp2/i2i\endcsname{30.74}
\expandafter\def\csname raw@v032/omni-single-tp2/i2i\endcsname{30.740}
\expandafter\def\csname warm@v032/omni-single-tp2/i2i\endcsname{30.9}
\expandafter\def\csname any@v032/omni-two-stage-2gpu/t2i\endcsname{18.528}
\expandafter\def\csname val@v032/omni-two-stage-2gpu/t2i\endcsname{18.53}
\expandafter\def\csname raw@v032/omni-two-stage-2gpu/t2i\endcsname{18.528}
\expandafter\def\csname warm@v032/omni-two-stage-2gpu/t2i\endcsname{30.4}
\expandafter\def\csname any@v032/omni-two-stage-2gpu/i2t1\endcsname{1.523}
\expandafter\def\csname val@v032/omni-two-stage-2gpu/i2t1\endcsname{1.52}
\expandafter\def\csname raw@v032/omni-two-stage-2gpu/i2t1\endcsname{1.523}
\expandafter\def\csname warm@v032/omni-two-stage-2gpu/i2t1\endcsname{1.8}
\expandafter\def\csname any@v032/omni-two-stage-2gpu/i2t64\endcsname{3.347}
\expandafter\def\csname val@v032/omni-two-stage-2gpu/i2t64\endcsname{3.35}
\expandafter\def\csname raw@v032/omni-two-stage-2gpu/i2t64\endcsname{3.347}
\expandafter\def\csname warm@v032/omni-two-stage-2gpu/i2t64\endcsname{3.3}
\expandafter\def\csname any@v032/omni-two-stage-2gpu/i2i\endcsname{40.125}
\expandafter\def\csname val@v032/omni-two-stage-2gpu/i2i\endcsname{40.13}
\expandafter\def\csname raw@v032/omni-two-stage-2gpu/i2i\endcsname{40.125}
\expandafter\def\csname warm@v032/omni-two-stage-2gpu/i2i\endcsname{52.5}
\expandafter\def\csname any@fact/mstar-cfgpar-nocache/i2t1\endcsname{0.293}
\expandafter\def\csname val@fact/mstar-cfgpar-nocache/i2t1\endcsname{0.29}
\expandafter\def\csname raw@fact/mstar-cfgpar-nocache/i2t1\endcsname{0.293}
\expandafter\def\csname warm@fact/mstar-cfgpar-nocache/i2t1\endcsname{6.7}
\expandafter\def\csname any@fact/mstar-cfgpar-nocache/i2t64\endcsname{2.656}
\expandafter\def\csname val@fact/mstar-cfgpar-nocache/i2t64\endcsname{2.66}
\expandafter\def\csname raw@fact/mstar-cfgpar-nocache/i2t64\endcsname{2.656}
\expandafter\def\csname warm@fact/mstar-cfgpar-nocache/i2t64\endcsname{3.6}
\expandafter\def\csname any@fact/mstar-cfgpar-nocache/i2i\endcsname{16.815}
\expandafter\def\csname val@fact/mstar-cfgpar-nocache/i2i\endcsname{16.82}
\expandafter\def\csname raw@fact/mstar-cfgpar-nocache/i2i\endcsname{16.815}
\expandafter\def\csname warm@fact/mstar-cfgpar-nocache/i2i\endcsname{32.3}
\expandafter\def\csname any@fact/omni-two-stage-mooncake-tcp/t2i\endcsname{17.983}
\expandafter\def\csname val@fact/omni-two-stage-mooncake-tcp/t2i\endcsname{17.98}
\expandafter\def\csname raw@fact/omni-two-stage-mooncake-tcp/t2i\endcsname{17.983}
\expandafter\def\csname warm@fact/omni-two-stage-mooncake-tcp/t2i\endcsname{20.7}
\expandafter\def\csname any@fact/omni-two-stage-mooncake-tcp/i2i\endcsname{29.922}
\expandafter\def\csname any@fact/omni-two-stage-nixl/t2i\endcsname{48.318}
\expandafter\def\csname any@fact/omni-two-stage-nixl/i2i\endcsname{121.569}
\expandafter\def\csname any@fact/omni-two-stage-prefixcache/i2t1\endcsname{0.053}
\expandafter\def\csname val@fact/omni-two-stage-prefixcache/i2t1\endcsname{0.05}
\expandafter\def\csname raw@fact/omni-two-stage-prefixcache/i2t1\endcsname{0.053}
\expandafter\def\csname warm@fact/omni-two-stage-prefixcache/i2t1\endcsname{0.6}
\expandafter\def\csname any@fact/omni-two-stage-prefixcache/i2t64\endcsname{1.744}
\expandafter\def\csname val@fact/omni-two-stage-prefixcache/i2t64\endcsname{1.74}
\expandafter\def\csname raw@fact/omni-two-stage-prefixcache/i2t64\endcsname{1.744}
\expandafter\def\csname warm@fact/omni-two-stage-prefixcache/i2t64\endcsname{1.7}
\expandafter\def\csname any@fact/omni-two-stage-shm/t2i\endcsname{17.962}
\expandafter\def\csname val@fact/omni-two-stage-shm/t2i\endcsname{17.96}
\expandafter\def\csname raw@fact/omni-two-stage-shm/t2i\endcsname{17.962}
\expandafter\def\csname warm@fact/omni-two-stage-shm/t2i\endcsname{18.6}
\expandafter\def\csname any@fact/omni-two-stage-shm/i2i\endcsname{30.744}
\expandafter\def\csname any@g3/mstar-cfgpar3-eager/t2i\endcsname{7.732}
\expandafter\def\csname val@g3/mstar-cfgpar3-eager/t2i\endcsname{7.73}
\expandafter\def\csname raw@g3/mstar-cfgpar3-eager/t2i\endcsname{7.732}
\expandafter\def\csname warm@g3/mstar-cfgpar3-eager/t2i\endcsname{17.4}
\expandafter\def\csname any@g3/mstar-cfgpar3-eager/i2t1\endcsname{0.089}
\expandafter\def\csname val@g3/mstar-cfgpar3-eager/i2t1\endcsname{0.09}
\expandafter\def\csname raw@g3/mstar-cfgpar3-eager/i2t1\endcsname{0.089}
\expandafter\def\csname warm@g3/mstar-cfgpar3-eager/i2t1\endcsname{4.5}
\expandafter\def\csname any@g3/mstar-cfgpar3-eager/i2t64\endcsname{2.424}
\expandafter\def\csname val@g3/mstar-cfgpar3-eager/i2t64\endcsname{2.42}
\expandafter\def\csname raw@g3/mstar-cfgpar3-eager/i2t64\endcsname{2.424}
\expandafter\def\csname warm@g3/mstar-cfgpar3-eager/i2t64\endcsname{3.3}
\expandafter\def\csname any@g3/mstar-cfgpar3-eager/i2i\endcsname{9.282}
\expandafter\def\csname val@g3/mstar-cfgpar3-eager/i2i\endcsname{9.28}
\expandafter\def\csname raw@g3/mstar-cfgpar3-eager/i2i\endcsname{9.282}
\expandafter\def\csname warm@g3/mstar-cfgpar3-eager/i2i\endcsname{24.5}
\expandafter\def\csname any@g3/mstar-cfgpar3-eager-nocache/i2t1\endcsname{0.304}
\expandafter\def\csname val@g3/mstar-cfgpar3-eager-nocache/i2t1\endcsname{0.30}
\expandafter\def\csname raw@g3/mstar-cfgpar3-eager-nocache/i2t1\endcsname{0.304}
\expandafter\def\csname warm@g3/mstar-cfgpar3-eager-nocache/i2t1\endcsname{5.9}
\expandafter\def\csname any@g3/mstar-cfgpar3-eager-nocache/i2t64\endcsname{2.685}
\expandafter\def\csname val@g3/mstar-cfgpar3-eager-nocache/i2t64\endcsname{2.68}
\expandafter\def\csname raw@g3/mstar-cfgpar3-eager-nocache/i2t64\endcsname{2.685}
\expandafter\def\csname warm@g3/mstar-cfgpar3-eager-nocache/i2t64\endcsname{3.4}
\expandafter\def\csname any@g3/mstar-cfgpar3-graph/t2i\endcsname{7.750}
\expandafter\def\csname val@g3/mstar-cfgpar3-graph/t2i\endcsname{7.75}
\expandafter\def\csname raw@g3/mstar-cfgpar3-graph/t2i\endcsname{7.750}
\expandafter\def\csname warm@g3/mstar-cfgpar3-graph/t2i\endcsname{15.2}
\expandafter\def\csname any@g3/mstar-cfgpar3-graph/i2t1\endcsname{0.085}
\expandafter\def\csname val@g3/mstar-cfgpar3-graph/i2t1\endcsname{0.09}
\expandafter\def\csname raw@g3/mstar-cfgpar3-graph/i2t1\endcsname{0.085}
\expandafter\def\csname warm@g3/mstar-cfgpar3-graph/i2t1\endcsname{3.0}
\expandafter\def\csname any@g3/mstar-cfgpar3-graph/i2t64\endcsname{1.951}
\expandafter\def\csname val@g3/mstar-cfgpar3-graph/i2t64\endcsname{1.95}
\expandafter\def\csname raw@g3/mstar-cfgpar3-graph/i2t64\endcsname{1.951}
\expandafter\def\csname warm@g3/mstar-cfgpar3-graph/i2t64\endcsname{1.9}
\expandafter\def\csname any@g3/mstar-cfgpar3-graph/i2i\endcsname{9.264}
\expandafter\def\csname val@g3/mstar-cfgpar3-graph/i2i\endcsname{9.26}
\expandafter\def\csname raw@g3/mstar-cfgpar3-graph/i2i\endcsname{9.264}
\expandafter\def\csname warm@g3/mstar-cfgpar3-graph/i2i\endcsname{14.9}
\expandafter\def\csname any@g3/mstar-cfgpar3-graph-nocache/i2t1\endcsname{0.300}
\expandafter\def\csname val@g3/mstar-cfgpar3-graph-nocache/i2t1\endcsname{0.30}
\expandafter\def\csname raw@g3/mstar-cfgpar3-graph-nocache/i2t1\endcsname{0.300}
\expandafter\def\csname warm@g3/mstar-cfgpar3-graph-nocache/i2t1\endcsname{2.1}
\expandafter\def\csname any@g3/mstar-cfgpar3-graph-nocache/i2t64\endcsname{2.152}
\expandafter\def\csname val@g3/mstar-cfgpar3-graph-nocache/i2t64\endcsname{2.15}
\expandafter\def\csname raw@g3/mstar-cfgpar3-graph-nocache/i2t64\endcsname{2.152}
\expandafter\def\csname warm@g3/mstar-cfgpar3-graph-nocache/i2t64\endcsname{2.2}
\expandafter\def\csname any@g3/mstar-cfgpar3vit-eager/t2i\endcsname{7.798}
\expandafter\def\csname val@g3/mstar-cfgpar3vit-eager/t2i\endcsname{7.80}
\expandafter\def\csname raw@g3/mstar-cfgpar3vit-eager/t2i\endcsname{7.798}
\expandafter\def\csname warm@g3/mstar-cfgpar3vit-eager/t2i\endcsname{31.5}
\expandafter\def\csname any@g3/mstar-cfgpar3vit-eager/i2t1\endcsname{0.107}
\expandafter\def\csname val@g3/mstar-cfgpar3vit-eager/i2t1\endcsname{0.11}
\expandafter\def\csname raw@g3/mstar-cfgpar3vit-eager/i2t1\endcsname{0.107}
\expandafter\def\csname warm@g3/mstar-cfgpar3vit-eager/i2t1\endcsname{6.6}
\expandafter\def\csname any@g3/mstar-cfgpar3vit-eager/i2t64\endcsname{2.410}
\expandafter\def\csname val@g3/mstar-cfgpar3vit-eager/i2t64\endcsname{2.41}
\expandafter\def\csname raw@g3/mstar-cfgpar3vit-eager/i2t64\endcsname{2.410}
\expandafter\def\csname warm@g3/mstar-cfgpar3vit-eager/i2t64\endcsname{3.1}
\expandafter\def\csname any@g3/mstar-cfgpar3vit-eager/i2i\endcsname{16.967}
\expandafter\def\csname val@g3/mstar-cfgpar3vit-eager/i2i\endcsname{16.97}
\expandafter\def\csname raw@g3/mstar-cfgpar3vit-eager/i2i\endcsname{16.967}
\expandafter\def\csname warm@g3/mstar-cfgpar3vit-eager/i2i\endcsname{31.7}
\expandafter\def\csname any@g3/mstar-cfgpar3vit-graph/t2i\endcsname{7.801}
\expandafter\def\csname val@g3/mstar-cfgpar3vit-graph/t2i\endcsname{7.80}
\expandafter\def\csname raw@g3/mstar-cfgpar3vit-graph/t2i\endcsname{7.801}
\expandafter\def\csname warm@g3/mstar-cfgpar3vit-graph/t2i\endcsname{13.9}
\expandafter\def\csname any@g3/mstar-cfgpar3vit-graph/i2t1\endcsname{0.107}
\expandafter\def\csname val@g3/mstar-cfgpar3vit-graph/i2t1\endcsname{0.11}
\expandafter\def\csname raw@g3/mstar-cfgpar3vit-graph/i2t1\endcsname{0.107}
\expandafter\def\csname warm@g3/mstar-cfgpar3vit-graph/i2t1\endcsname{3.3}
\expandafter\def\csname any@g3/mstar-cfgpar3vit-graph/i2t64\endcsname{1.971}
\expandafter\def\csname val@g3/mstar-cfgpar3vit-graph/i2t64\endcsname{1.97}
\expandafter\def\csname raw@g3/mstar-cfgpar3vit-graph/i2t64\endcsname{1.971}
\expandafter\def\csname warm@g3/mstar-cfgpar3vit-graph/i2t64\endcsname{1.9}
\expandafter\def\csname any@g3/mstar-cfgpar3vit-graph/i2i\endcsname{16.971}
\expandafter\def\csname val@g3/mstar-cfgpar3vit-graph/i2i\endcsname{16.97}
\expandafter\def\csname raw@g3/mstar-cfgpar3vit-graph/i2i\endcsname{16.971}
\expandafter\def\csname warm@g3/mstar-cfgpar3vit-graph/i2i\endcsname{19.7}
\expandafter\def\csname any@g3/mstar-phase-diff/t2i\endcsname{7.835}
\expandafter\def\csname val@g3/mstar-phase-diff/t2i\endcsname{7.83}
\expandafter\def\csname raw@g3/mstar-phase-diff/t2i\endcsname{7.835}
\expandafter\def\csname warm@g3/mstar-phase-diff/t2i\endcsname{15.7}
\expandafter\def\csname any@g3/mstar-phase-diff/i2i\endcsname{9.336}
\expandafter\def\csname val@g3/mstar-phase-diff/i2i\endcsname{9.34}
\expandafter\def\csname raw@g3/mstar-phase-diff/i2i\endcsname{9.336}
\expandafter\def\csname warm@g3/mstar-phase-diff/i2i\endcsname{18.4}
\expandafter\def\csname any@g3/mstar-phase-eager/i2t1\endcsname{0.092}
\expandafter\def\csname val@g3/mstar-phase-eager/i2t1\endcsname{0.09}
\expandafter\def\csname raw@g3/mstar-phase-eager/i2t1\endcsname{0.092}
\expandafter\def\csname warm@g3/mstar-phase-eager/i2t1\endcsname{6.8}
\expandafter\def\csname any@g3/mstar-phase-eager/i2t64\endcsname{2.476}
\expandafter\def\csname val@g3/mstar-phase-eager/i2t64\endcsname{2.48}
\expandafter\def\csname raw@g3/mstar-phase-eager/i2t64\endcsname{2.476}
\expandafter\def\csname warm@g3/mstar-phase-eager/i2t64\endcsname{3.3}
\expandafter\def\csname any@g3/mstar-phase-graph/i2t1\endcsname{0.087}
\expandafter\def\csname val@g3/mstar-phase-graph/i2t1\endcsname{0.09}
\expandafter\def\csname raw@g3/mstar-phase-graph/i2t1\endcsname{0.087}
\expandafter\def\csname warm@g3/mstar-phase-graph/i2t1\endcsname{3.4}
\expandafter\def\csname any@g3/mstar-phase-graph/i2t64\endcsname{2.013}
\expandafter\def\csname val@g3/mstar-phase-graph/i2t64\endcsname{2.01}
\expandafter\def\csname raw@g3/mstar-phase-graph/i2t64\endcsname{2.013}
\expandafter\def\csname warm@g3/mstar-phase-graph/i2t64\endcsname{2.0}
\expandafter\def\csname any@g3/mstar-tp2-3gpu-eager/t2i\endcsname{7.676}
\expandafter\def\csname val@g3/mstar-tp2-3gpu-eager/t2i\endcsname{7.68}
\expandafter\def\csname raw@g3/mstar-tp2-3gpu-eager/t2i\endcsname{7.676}
\expandafter\def\csname warm@g3/mstar-tp2-3gpu-eager/t2i\endcsname{18.8}
\expandafter\def\csname any@g3/mstar-tp2-3gpu-eager/i2t1\endcsname{0.259}
\expandafter\def\csname val@g3/mstar-tp2-3gpu-eager/i2t1\endcsname{0.26}
\expandafter\def\csname raw@g3/mstar-tp2-3gpu-eager/i2t1\endcsname{0.259}
\expandafter\def\csname warm@g3/mstar-tp2-3gpu-eager/i2t1\endcsname{4.7}
\expandafter\def\csname any@g3/mstar-tp2-3gpu-eager/i2t64\endcsname{2.764}
\expandafter\def\csname val@g3/mstar-tp2-3gpu-eager/i2t64\endcsname{2.76}
\expandafter\def\csname raw@g3/mstar-tp2-3gpu-eager/i2t64\endcsname{2.764}
\expandafter\def\csname warm@g3/mstar-tp2-3gpu-eager/i2t64\endcsname{4.0}
\expandafter\def\csname any@g3/omni-prof-i2i-cfgpar3/i2i\endcsname{11.264}
\expandafter\def\csname val@g3/omni-prof-i2i-cfgpar3/i2i\endcsname{11.26}
\expandafter\def\csname raw@g3/omni-prof-i2i-cfgpar3/i2i\endcsname{11.264}
\expandafter\def\csname warm@g3/omni-prof-i2i-cfgpar3/i2i\endcsname{12.8}
\expandafter\def\csname any@g3/omni-prof-t2i-cfgpar2/t2i\endcsname{10.020}
\expandafter\def\csname val@g3/omni-prof-t2i-cfgpar2/t2i\endcsname{10.02}
\expandafter\def\csname raw@g3/omni-prof-t2i-cfgpar2/t2i\endcsname{10.020}
\expandafter\def\csname warm@g3/omni-prof-t2i-cfgpar2/t2i\endcsname{10.9}
\expandafter\def\csname any@g3/omni-single-cfgpar3-i2i/i2i\endcsname{11.190}
\expandafter\def\csname val@g3/omni-single-cfgpar3-i2i/i2i\endcsname{11.19}
\expandafter\def\csname raw@g3/omni-single-cfgpar3-i2i/i2i\endcsname{11.190}
\expandafter\def\csname warm@g3/omni-single-cfgpar3-i2i/i2i\endcsname{12.9}
\expandafter\def\csname any@g3/omni-tp2-eager-cache/t2i\endcsname{17.951}
\expandafter\def\csname val@g3/omni-tp2-eager-cache/t2i\endcsname{17.95}
\expandafter\def\csname raw@g3/omni-tp2-eager-cache/t2i\endcsname{17.951}
\expandafter\def\csname warm@g3/omni-tp2-eager-cache/t2i\endcsname{18.6}
\expandafter\def\csname any@g3/omni-tp2-eager-cache/i2t1\endcsname{0.055}
\expandafter\def\csname val@g3/omni-tp2-eager-cache/i2t1\endcsname{0.05}
\expandafter\def\csname raw@g3/omni-tp2-eager-cache/i2t1\endcsname{0.055}
\expandafter\def\csname warm@g3/omni-tp2-eager-cache/i2t1\endcsname{0.3}
\expandafter\def\csname any@g3/omni-tp2-eager-cache/i2t64\endcsname{1.280}
\expandafter\def\csname val@g3/omni-tp2-eager-cache/i2t64\endcsname{1.28}
\expandafter\def\csname raw@g3/omni-tp2-eager-cache/i2t64\endcsname{1.280}
\expandafter\def\csname warm@g3/omni-tp2-eager-cache/i2t64\endcsname{1.3}
\expandafter\def\csname any@g3/omni-tp2-eager-cache/i2i\endcsname{30.027}
\expandafter\def\csname val@g3/omni-tp2-eager-cache/i2i\endcsname{30.03}
\expandafter\def\csname raw@g3/omni-tp2-eager-cache/i2i\endcsname{30.027}
\expandafter\def\csname warm@g3/omni-tp2-eager-cache/i2i\endcsname{31.0}
\expandafter\def\csname any@g3/omni-tp2-eager-chat/i2t1\endcsname{0.127}
\expandafter\def\csname val@g3/omni-tp2-eager-chat/i2t1\endcsname{0.13}
\expandafter\def\csname raw@g3/omni-tp2-eager-chat/i2t1\endcsname{0.127}
\expandafter\def\csname warm@g3/omni-tp2-eager-chat/i2t1\endcsname{0.6}
\expandafter\def\csname any@g3/omni-tp2-eager-chat/i2t64\endcsname{1.361}
\expandafter\def\csname val@g3/omni-tp2-eager-chat/i2t64\endcsname{1.36}
\expandafter\def\csname raw@g3/omni-tp2-eager-chat/i2t64\endcsname{1.361}
\expandafter\def\csname warm@g3/omni-tp2-eager-chat/i2t64\endcsname{1.4}
\expandafter\def\csname any@g3/omni-tp2-eager-chat-cache/i2t1\endcsname{0.054}
\expandafter\def\csname val@g3/omni-tp2-eager-chat-cache/i2t1\endcsname{0.05}
\expandafter\def\csname raw@g3/omni-tp2-eager-chat-cache/i2t1\endcsname{0.054}
\expandafter\def\csname warm@g3/omni-tp2-eager-chat-cache/i2t1\endcsname{0.6}
\expandafter\def\csname any@g3/omni-tp2-eager-chat-cache/i2t64\endcsname{1.254}
\expandafter\def\csname val@g3/omni-tp2-eager-chat-cache/i2t64\endcsname{1.25}
\expandafter\def\csname raw@g3/omni-tp2-eager-chat-cache/i2t64\endcsname{1.254}
\expandafter\def\csname warm@g3/omni-tp2-eager-chat-cache/i2t64\endcsname{1.5}
\expandafter\def\csname any@g3/omni-tp2-eager-i2tonly/i2t1\endcsname{0.122}
\expandafter\def\csname val@g3/omni-tp2-eager-i2tonly/i2t1\endcsname{0.12}
\expandafter\def\csname raw@g3/omni-tp2-eager-i2tonly/i2t1\endcsname{0.122}
\expandafter\def\csname warm@g3/omni-tp2-eager-i2tonly/i2t1\endcsname{0.5}
\expandafter\def\csname any@g3/omni-tp2-eager-i2tonly/i2t64\endcsname{1.373}
\expandafter\def\csname val@g3/omni-tp2-eager-i2tonly/i2t64\endcsname{1.37}
\expandafter\def\csname raw@g3/omni-tp2-eager-i2tonly/i2t64\endcsname{1.373}
\expandafter\def\csname warm@g3/omni-tp2-eager-i2tonly/i2t64\endcsname{1.4}
\expandafter\def\csname any@g3/omni-tp2-graph-cache/t2i\endcsname{18.010}
\expandafter\def\csname val@g3/omni-tp2-graph-cache/t2i\endcsname{18.01}
\expandafter\def\csname raw@g3/omni-tp2-graph-cache/t2i\endcsname{18.010}
\expandafter\def\csname warm@g3/omni-tp2-graph-cache/t2i\endcsname{18.4}
\expandafter\def\csname any@g3/omni-tp2-graph-cache/i2t1\endcsname{0.051}
\expandafter\def\csname val@g3/omni-tp2-graph-cache/i2t1\endcsname{0.05}
\expandafter\def\csname raw@g3/omni-tp2-graph-cache/i2t1\endcsname{0.051}
\expandafter\def\csname warm@g3/omni-tp2-graph-cache/i2t1\endcsname{0.3}
\expandafter\def\csname any@g3/omni-tp2-graph-cache/i2t64\endcsname{1.030}
\expandafter\def\csname val@g3/omni-tp2-graph-cache/i2t64\endcsname{1.03}
\expandafter\def\csname raw@g3/omni-tp2-graph-cache/i2t64\endcsname{1.030}
\expandafter\def\csname warm@g3/omni-tp2-graph-cache/i2t64\endcsname{1.0}
\expandafter\def\csname any@g3/omni-tp2-graph-cache/i2i\endcsname{30.087}
\expandafter\def\csname val@g3/omni-tp2-graph-cache/i2i\endcsname{30.09}
\expandafter\def\csname raw@g3/omni-tp2-graph-cache/i2i\endcsname{30.087}
\expandafter\def\csname warm@g3/omni-tp2-graph-cache/i2i\endcsname{31.5}
\expandafter\def\csname any@g3/omni-tp2-graph-chat/i2t1\endcsname{0.125}
\expandafter\def\csname val@g3/omni-tp2-graph-chat/i2t1\endcsname{0.13}
\expandafter\def\csname raw@g3/omni-tp2-graph-chat/i2t1\endcsname{0.125}
\expandafter\def\csname warm@g3/omni-tp2-graph-chat/i2t1\endcsname{0.4}
\expandafter\def\csname any@g3/omni-tp2-graph-chat/i2t64\endcsname{1.103}
\expandafter\def\csname val@g3/omni-tp2-graph-chat/i2t64\endcsname{1.10}
\expandafter\def\csname raw@g3/omni-tp2-graph-chat/i2t64\endcsname{1.103}
\expandafter\def\csname warm@g3/omni-tp2-graph-chat/i2t64\endcsname{1.1}
\expandafter\def\csname any@g3/omni-tp2-graph-chat-cache/i2t1\endcsname{0.051}
\expandafter\def\csname val@g3/omni-tp2-graph-chat-cache/i2t1\endcsname{0.05}
\expandafter\def\csname raw@g3/omni-tp2-graph-chat-cache/i2t1\endcsname{0.051}
\expandafter\def\csname warm@g3/omni-tp2-graph-chat-cache/i2t1\endcsname{0.4}
\expandafter\def\csname any@g3/omni-tp2-graph-chat-cache/i2t64\endcsname{1.022}
\expandafter\def\csname val@g3/omni-tp2-graph-chat-cache/i2t64\endcsname{1.02}
\expandafter\def\csname raw@g3/omni-tp2-graph-chat-cache/i2t64\endcsname{1.022}
\expandafter\def\csname warm@g3/omni-tp2-graph-chat-cache/i2t64\endcsname{1.0}
\expandafter\def\csname any@g3/omni-two-stage-3gpu-tp2-eager/t2i\endcsname{17.912}
\expandafter\def\csname val@g3/omni-two-stage-3gpu-tp2-eager/t2i\endcsname{17.91}
\expandafter\def\csname raw@g3/omni-two-stage-3gpu-tp2-eager/t2i\endcsname{17.912}
\expandafter\def\csname warm@g3/omni-two-stage-3gpu-tp2-eager/t2i\endcsname{19.4}
\expandafter\def\csname any@g3/omni-two-stage-3gpu-tp2-eager/i2t1\endcsname{0.123}
\expandafter\def\csname val@g3/omni-two-stage-3gpu-tp2-eager/i2t1\endcsname{0.12}
\expandafter\def\csname raw@g3/omni-two-stage-3gpu-tp2-eager/i2t1\endcsname{0.123}
\expandafter\def\csname warm@g3/omni-two-stage-3gpu-tp2-eager/i2t1\endcsname{0.4}
\expandafter\def\csname any@g3/omni-two-stage-3gpu-tp2-eager/i2t64\endcsname{1.481}
\expandafter\def\csname val@g3/omni-two-stage-3gpu-tp2-eager/i2t64\endcsname{1.48}
\expandafter\def\csname raw@g3/omni-two-stage-3gpu-tp2-eager/i2t64\endcsname{1.481}
\expandafter\def\csname warm@g3/omni-two-stage-3gpu-tp2-eager/i2t64\endcsname{1.4}
\expandafter\def\csname any@g3/omni-two-stage-3gpu-tp2-eager/i2i\endcsname{30.371}
\expandafter\def\csname val@g3/omni-two-stage-3gpu-tp2-eager/i2i\endcsname{30.37}
\expandafter\def\csname raw@g3/omni-two-stage-3gpu-tp2-eager/i2i\endcsname{30.371}
\expandafter\def\csname warm@g3/omni-two-stage-3gpu-tp2-eager/i2i\endcsname{31.1}
\expandafter\def\csname any@g3/omni-two-stage-3gpu-tp2-graph/t2i\endcsname{18.014}
\expandafter\def\csname val@g3/omni-two-stage-3gpu-tp2-graph/t2i\endcsname{18.01}
\expandafter\def\csname raw@g3/omni-two-stage-3gpu-tp2-graph/t2i\endcsname{18.014}
\expandafter\def\csname warm@g3/omni-two-stage-3gpu-tp2-graph/t2i\endcsname{18.4}
\expandafter\def\csname any@g3/omni-two-stage-3gpu-tp2-graph/i2t1\endcsname{0.123}
\expandafter\def\csname val@g3/omni-two-stage-3gpu-tp2-graph/i2t1\endcsname{0.12}
\expandafter\def\csname raw@g3/omni-two-stage-3gpu-tp2-graph/i2t1\endcsname{0.123}
\expandafter\def\csname warm@g3/omni-two-stage-3gpu-tp2-graph/i2t1\endcsname{0.3}
\expandafter\def\csname any@g3/omni-two-stage-3gpu-tp2-graph/i2t64\endcsname{1.100}
\expandafter\def\csname val@g3/omni-two-stage-3gpu-tp2-graph/i2t64\endcsname{1.10}
\expandafter\def\csname raw@g3/omni-two-stage-3gpu-tp2-graph/i2t64\endcsname{1.100}
\expandafter\def\csname warm@g3/omni-two-stage-3gpu-tp2-graph/i2t64\endcsname{1.1}
\expandafter\def\csname any@g3/omni-two-stage-3gpu-tp2-graph/i2i\endcsname{30.363}
\expandafter\def\csname val@g3/omni-two-stage-3gpu-tp2-graph/i2i\endcsname{30.36}
\expandafter\def\csname raw@g3/omni-two-stage-3gpu-tp2-graph/i2i\endcsname{30.363}
\expandafter\def\csname warm@g3/omni-two-stage-3gpu-tp2-graph/i2i\endcsname{31.8}
\expandafter\def\csname any@tp2g/mstar-tp2-eager-rerun/t2i\endcsname{7.680}
\expandafter\def\csname val@tp2g/mstar-tp2-eager-rerun/t2i\endcsname{7.68}
\expandafter\def\csname raw@tp2g/mstar-tp2-eager-rerun/t2i\endcsname{7.680}
\expandafter\def\csname warm@tp2g/mstar-tp2-eager-rerun/t2i\endcsname{14.0}
\expandafter\def\csname any@tp2g/mstar-tp2-eager-rerun/i2t1\endcsname{0.258}
\expandafter\def\csname val@tp2g/mstar-tp2-eager-rerun/i2t1\endcsname{0.26}
\expandafter\def\csname raw@tp2g/mstar-tp2-eager-rerun/i2t1\endcsname{0.258}
\expandafter\def\csname warm@tp2g/mstar-tp2-eager-rerun/i2t1\endcsname{11.0}
\expandafter\def\csname any@tp2g/mstar-tp2-eager-rerun/i2t64\endcsname{3.049}
\expandafter\def\csname val@tp2g/mstar-tp2-eager-rerun/i2t64\endcsname{3.05}
\expandafter\def\csname raw@tp2g/mstar-tp2-eager-rerun/i2t64\endcsname{3.049}
\expandafter\def\csname warm@tp2g/mstar-tp2-eager-rerun/i2t64\endcsname{3.6}
\expandafter\def\csname any@tp2g/mstar-tp2-graph-fixed/t2i\endcsname{7.640}
\expandafter\def\csname val@tp2g/mstar-tp2-graph-fixed/t2i\endcsname{7.64}
\expandafter\def\csname raw@tp2g/mstar-tp2-graph-fixed/t2i\endcsname{7.640}
\expandafter\def\csname warm@tp2g/mstar-tp2-graph-fixed/t2i\endcsname{15.1}
\expandafter\def\csname any@tp2g/mstar-tp2-graph-fixed/i2t1\endcsname{0.226}
\expandafter\def\csname val@tp2g/mstar-tp2-graph-fixed/i2t1\endcsname{0.23}
\expandafter\def\csname raw@tp2g/mstar-tp2-graph-fixed/i2t1\endcsname{0.226}
\expandafter\def\csname warm@tp2g/mstar-tp2-graph-fixed/i2t1\endcsname{5.0}
\expandafter\def\csname any@tp2g/mstar-tp2-graph-fixed/i2t64\endcsname{1.525}
\expandafter\def\csname val@tp2g/mstar-tp2-graph-fixed/i2t64\endcsname{1.52}
\expandafter\def\csname raw@tp2g/mstar-tp2-graph-fixed/i2t64\endcsname{1.525}
\expandafter\def\csname warm@tp2g/mstar-tp2-graph-fixed/i2t64\endcsname{1.5}
\expandafter\def\csname pct@topoMstarTti\endcsname{$-$41\%}
\expandafter\def\csname pct@topoMstarIti\endcsname{$-$21\%}
\expandafter\def\csname pct@topoMstarDec\endcsname{$-$22\%}
\expandafter\def\csname pct@topoMstarPre\endcsname{$-$74\%}
\expandafter\def\csname pct@topoOmniTti\endcsname{$-$41\%}
\expandafter\def\csname pct@cmpCfgTti\endcsname{$-$9\%}
\expandafter\def\csname pct@cmpCfgIti\endcsname{$-$8\%}
\expandafter\def\csname pct@cmpCfgPre\endcsname{$+$3\%}
\expandafter\def\csname pct@cmpCfgDec\endcsname{$+$13\%}
\expandafter\def\csname pct@bestTti\endcsname{1.30$\times$}
\expandafter\def\csname pct@bestIti\endcsname{1.67$\times$}
\expandafter\def\csname pct@bestDecG\endcsname{0.76$\times$}
\expandafter\def\csname pct@bestDecN\endcsname{0.86$\times$}
\expandafter\def\csname pct@omniGraphDec\endcsname{$-$1\%}
\expandafter\def\csname pct@verOmniTti\endcsname{$-$3\%}
\expandafter\def\csname pct@verOmniIti\endcsname{$-$9\%}
\expandafter\def\csname pct@verOmni2Dec\endcsname{$-$44\%}
\expandafter\def\csname pct@connMoonTti\endcsname{$+$0\%}
\expandafter\def\csname pct@connMoonIti\endcsname{$-$3\%}
\expandafter\def\csname pct@connNixlTti\endcsname{失败\textsuperscript{‡}}
\expandafter\def\csname pct@connNixlIti\endcsname{失败\textsuperscript{‡}}
\expandafter\def\csname pct@cacheOmni1\endcsname{$-$73\%}
\expandafter\def\csname pct@cacheOmni64\endcsname{$-$7\%}
\expandafter\def\csname pct@cacheMstar1\endcsname{$-$71\%}
\expandafter\def\csname pct@cacheMstar64\endcsname{$-$7\%}
\expandafter\def\csname pct@cacheMstarIti\endcsname{$-$0\%}
\expandafter\def\csname pct@i2tOff1\endcsname{$-$34\%}
\expandafter\def\csname pct@i2tOff64\endcsname{$-$29\%}
\expandafter\def\csname pct@i2tOn1\endcsname{$-$38\%}
\expandafter\def\csname pct@i2tOn64\endcsname{$-$29\%}
\expandafter\def\csname pct@i2tOff1x\endcsname{1.52$\times$}
\expandafter\def\csname pct@i2tOn1x\endcsname{1.60$\times$}
\expandafter\def\csname pct@i2tOff64x\endcsname{1.41$\times$}
\expandafter\def\csname pct@i2tOn64x\endcsname{1.41$\times$}
\expandafter\def\csname pct@msTokM\endcsname{37.7}
\expandafter\def\csname pct@msTokMnc\endcsname{33.1}
\expandafter\def\csname pct@msTokO\endcsname{26.8}
\expandafter\def\csname pct@msTokMoff\endcsname{37.5}
\expandafter\def\csname pct@msTokOg\endcsname{26.8}
\expandafter\def\csname pct@sumT2I\endcsname{$+$30\%}
\expandafter\def\csname pct@sumI2I\endcsname{$+$67\%}
\expandafter\def\csname pct@sumI1off\endcsname{$-$34\%}
\expandafter\def\csname pct@sumI64off\endcsname{$-$29\%}
\expandafter\def\csname pct@sumI1on\endcsname{$-$38\%}
\expandafter\def\csname pct@sumI64on\endcsname{$-$29\%}
\expandafter\def\csname pct@mFastT2I\endcsname{23\%}
\expandafter\def\csname pct@mFastI2I\endcsname{40\%}
\expandafter\def\csname pct@oFastI1off\endcsname{34\%}
\expandafter\def\csname pct@oFastI64off\endcsname{29\%}
\expandafter\def\csname pct@oFastI1on\endcsname{38\%}
\expandafter\def\csname pct@oFastI64on\endcsname{29\%}
\expandafter\def\csname pct@gCfgOn64\endcsname{20\%}
\expandafter\def\csname pct@gTp2G64\endcsname{50\%}
\expandafter\def\csname pct@gTp2Gap\endcsname{38\%}
\expandafter\def\csname pct@gDiffT\endcsname{$+$0\%}
\expandafter\def\csname pct@gDiffI\endcsname{$-$0\%}
\expandafter\def\csname pct@par_mstar-cfgpar_i2t1\endcsname{$-$56\%}
\expandafter\def\csname pct@par_mstar-cfgpar_i2t64\endcsname{$-$4\%}
\expandafter\def\csname pct@par_mstar-cfgpar-nocompile_i2t1\endcsname{$-$59\%}
\expandafter\def\csname pct@par_mstar-cfgpar-nocompile_i2t64\endcsname{$-$5\%}
\expandafter\def\csname pct@par_mstar-tp2_i2t1\endcsname{$-$66\%}
\expandafter\def\csname pct@par_mstar-tp2_i2t64\endcsname{$-$16\%}
\expandafter\def\csname pct@par_mstar-cfgpar-nocache_i2t1\endcsname{$-$65\%}
\expandafter\def\csname pct@par_mstar-cfgpar-nocache_i2t64\endcsname{$-$17\%}
\makeatother


\title{M* 与 vLLM-Omni：支持范围、接入工作量与 Intel B70 实测}
\date{2026-11-25}
\hypersetup{pdftitle={M* 与 vLLM-Omni：支持范围、接入工作量与 Intel B70 实测}}

\begin{document}

\deckcover{M* 与 vLLM-Omni\\支持范围、接入工作量，以及 BAGEL 在 Intel Arc Pro B70 上的实测}%
  {各自最优配置、等工作量、全部实测结果列出并标出最佳}%
  {框架支持范围 · 模型接入工作量 · BAGEL-7B-MoT · 并行方式 / torch.compile / XPU graph / 两阶段 connector / 前缀缓存}%
  {结果汇报\quad|\quad Intel Arc Pro B70 $\times$2（卡 1、2）\\版本：M* 2c8ffde8f · vLLM-Omni main a36de509e（vLLM 0.34.0+xpu）与 v0.32.0}

% 总览：结论与各自最优配置的结果（每个数字都标出配置）
\setpart{总览}

\begin{frame}[t]{结论：图像生成 M* 更快，图生文 vLLM-Omni 更快}
\small
\begin{itemize}\setlength\itemsep{2pt}
  \item \kw{测试内容}：卡 1、2 上的 BAGEL-7B-MoT，各框架取最优配置，等工作量比较端到端延迟。
  \item \kw{接入工作量}（第 1、2 部分）：M* 约 8 个文件、5.7 千行；vLLM-Omni 约 20 个文件、8.5 千行。
  \item \kw{并行拓扑}：TP2 每次文生图约 5000 次跨卡 all-reduce；改用 CFG 并行后两框架文生图均快约 41\%。
  \item \kw{compile 与 graph}：compile 使扩散快 7--9\%、解码慢约 13\%。graph 只作用于图生文解码：三卡 TP2 下 64 token 快约 50\%，扩散步不受影响。
  \item \kw{前缀缓存}：M* 默认开，vLLM-Omni 默认关，因此缓存关、开分别比较。
\end{itemize}
\takeaway{文生图 M* 快 \pct{mFastT2I}，图像编辑 M* 快 \pct{mFastI2I}（各自最优配置）。图生文 vLLM-Omni 更快：缓存关时 1 token 快 \pct{oFastI1off}、64 token 快 \pct{oFastI64off}；缓存开时分别快 \pct{oFastI1on}、\pct{oFastI64on}。绝对值只适用于卡 1+2 与本次软件版本。}
\end{frame}

\begin{frame}[t]{各自最优配置的结果（稳态中位数，秒）}
\scriptsize\def\codesize{\relax}
\renewcommand{\arraystretch}{1.2}\setlength{\tabcolsep}{3pt}
\begin{tabularx}{\textwidth}{>{\bfseries}p{2.4cm} Y Y Y Y r}
\toprule
\rowcolor{Mist}
测试项 & \textcolor{Teal}{M*} 结果 & \textcolor{Teal}{M*} 配置 & \textcolor{Orange}{vLLM-Omni} 结果 & \textcolor{Orange}{vLLM-Omni} 配置 & Omni 相对 M* \\
\midrule
文生图 512$^2$，50 步 & \textbf{\val{final/mstar-cfgpar/t2i}} & CFG 并行，卡 1、2；compile 开；缓存开；无 graph & \val{final/omni-single-cfgpar2/t2i} & 单 stage CFG 并行，卡 1、2；无 compile & \pct{sumT2I} \\
图像编辑 512$^2$，50 步 & \textbf{\val{final/mstar-cfgpar/i2i}} & CFG 并行，卡 1、2；compile 开；无 graph & \val{final/omni-single-tp2/i2i} & 单 stage TP2，卡 1、2 & \pct{sumI2I} \\
图生文 1 token，缓存关 & \raw{fact/mstar-cfgpar-nocache/i2t1} & CFG 并行，卡 1、2；compile 开；缓存关 & \raw{final/omni-two-stage/i2t1} & two-stage，Thinker 卡 1、DiT 卡 2；缓存关；eager & \pct{sumI1off} \\
图生文 1 token，缓存开 & \raw{final/mstar-cfgpar/i2t1} & CFG 并行，卡 1、2；compile 开；缓存开 & \raw{fact/omni-two-stage-prefixcache/i2t1} & two-stage，Thinker 卡 1、DiT 卡 2；缓存开；eager & \pct{sumI1on} \\
图生文 64 token，缓存关 & \raw{fact/mstar-cfgpar-nocache/i2t64} & 同上（缓存关） & \raw{final/omni-two-stage/i2t64} & 同上（缓存关） & \pct{sumI64off} \\
图生文 64 token，缓存开 & \raw{final/mstar-cfgpar/i2t64} & 同上（缓存开） & \raw{fact/omni-two-stage-prefixcache/i2t64} & 同上（缓存开） & \pct{sumI64on} \\
\bottomrule
\end{tabularx}
\vspace{4pt}
\begin{itemize}\scriptsize\setlength\itemsep{1pt}
  \item “Omni 相对 M*”：正值表示 vLLM-Omni 更慢，负值表示更快。
  \item 图像编辑两边指令相同。图生文输入规模一致（M* 1369 patch，Omni 1386 token）。
  \item 配置的完整定义见第 3 部分；三卡的结果单列。
\end{itemize}
\end{frame}

\begin{frame}[t]{阅读说明与标记}
\small
\begin{itemize}\setlength\itemsep{3pt}
  \item \kw{第 1 部分：支持范围}：两个框架各支持哪些模型、特性与硬件。
  \item \kw{第 2 部分：模型接入工作量}：新模型接入需要做什么，两个框架的投入，BAGEL 的实际规模。
  \item \kw{第 3 部分：BAGEL 实测}：测试环境、配置定义、统一工作量，两卡的全部结果，各因素的影响，三卡的 XPU graph 研究，附录。
\end{itemize}
\vspace{2pt}
\scriptsize\renewcommand{\arraystretch}{1.12}
\begin{tabularx}{\textwidth}{>{\bfseries}p{3.0cm} Y}
\toprule
\rowcolor{Mist}
标记 & 含义 \\
\midrule
\cellcolor{GoldBg}\textbf{\textcolor{Teal}{7.68}}\,$\star$ & 同一框架内该 case 的最佳配置（1\% 以内并列） \\
\textcolor{Muted}{30.76\textsuperscript{†}} & 输出图像异常（灰图）：只保留耗时作参考，不参与比较 \\
\textcolor{Muted}{描述不正确\textsuperscript{†}} & 图生文的描述与输入图不符：不显示耗时 \\
\textcolor{Red}{失败\textsuperscript{‡}} & 该配置在 XPU 上运行失败（NIXL 传输） \\
\textcolor{Muted}{2.55\textsuperscript{§}} & 输入或缓存设置与对方不等，仅供参考 \\
-- & 未测或不支持 \\
\textcolor{Teal}{M*} / \textcolor{Orange}{Omni} & 框架；全文颜色一致（蓝 = M*，橙 = vLLM-Omni） \\
\bottomrule
\end{tabularx}
\takeaway{延迟 = 单个请求从发出到收到完整响应的端到端秒数。2 次 warmup 后做 5 次稳态测量，取中位数，串行发送；首请求（含 JIT 与编译）单列。这是延迟对比，不是吞吐对比。}
\end{frame}

% 支持范围：硬件、Intel GPU、模型（2026-11-13 核查，只读源码、文档与上游 PR；2026-11-22 增补 M* 的 Intel 实现、按模型缺口与工作量；仓库内未在 XPU 上运行，B70 结果为用户报告）
\setpart{支持范围}

\begin{frame}[t]{硬件支持：两个框架的“支持”含义不同}
\scriptsize\def\codesize{\relax}
\renewcommand{\arraystretch}{1.12}
\setlength{\tabcolsep}{3pt}
\begin{tabularx}{\textwidth}{>{\bfseries}p{2.1cm} Y Y}
\toprule
\rowcolor{Mist}
后端 & \textcolor{Teal}{M*（2c8ffde）} & vLLM-Omni（main a36de509e） \\
\midrule
NVIDIA CUDA & 代码支持；仓库 CI 只跑 \code{ubuntu-latest}，没有 GPU & 代码；CI 有 CUDA 任务 \\
AMD ROCm & 代码，部分（AITER 注意力，\code{attn/dense.py} 等）；文档有 AITER 说明（\code{docs/models.rst}）；没有 ROCm CI & 代码、文档（gfx942）、CI 在 MI300 \\
Intel XPU & 代码，只有 BAGEL 的 TP2 路径；文档声明；没有 XPU CI & 代码、文档；CI 有 1 卡与 2 卡任务（队列 \code{intel-gpu-omni}），仓库内标记 B60 与 B70，结果不可见 \\
CPU & CPU 测试套件在 \code{ubuntu-latest} 上用 CPU 版 PyTorch 运行，没有真实模型推理 & 只有 CPU 标记的测试 \\
Ascend NPU & 没有 NPU 平台 & 代码；CI 在 A2 / A3 / A5 / 310P \\
\bottomrule
\end{tabularx}
\takeaway{“支持”至少分四层：代码、文档、CI、公开的实测结果。两家在最后一层上都看不到 Intel 的公开数据；本报告第 3 部分在本机补上了 BAGEL 的实测。}
\end{frame}

\begin{frame}[t]{Intel GPU（一）：代码、文档与 CI}
\scriptsize\def\codesize{\relax}
\renewcommand{\arraystretch}{1.12}
\setlength{\tabcolsep}{3pt}
\begin{tabularx}{\textwidth}{>{\bfseries}p{2.1cm} Y Y}
\toprule
\rowcolor{Mist}
项目 & \textcolor{Teal}{M*} & vLLM-Omni \\
\midrule
平台代码 & \code{XpuPlatform}；XPU 图后端；\code{xpu\_paged} 注意力；RoPE、RMSNorm 融合核；XPU 分布式组 & 独立的 XPU 平台类 \\
文档声明 & BAGEL：“Intel GPUs with XPU PyTorch for BAGEL”（\code{installation.rst:9--10}） & 模型表的 Intel 列多数打勾；XPU 安装页先写“Validated on Intel Arc B-Series”，后面又写“runtime validation … remains pending” \\
CI & 没有 XPU 任务；CI 全部在 \code{ubuntu-latest} 上跑 & Buildkite Intel 队列（\code{intel-gpu-omni}），PR 有 1 卡与 2 卡任务；仓库内标记 B60 与 B70；结果不可见 \\
依赖栈 & \code{Dockerfile.xpu}：torch 2.13.0+xpu，克隆 \code{main}，不是固定 commit；\code{vllm-xpu-kernels==0.1.15.4} 见 \code{pyproject.toml}（约第 167 行） & XPU 文档：保留上游 runtime（\code{Dockerfile.xpu:41} 为 triton 3.9.0+xpu），不换成 triton-xpu；运行验证仍待定 \\
\bottomrule
\end{tabularx}
\takeaway{M* 有 XPU 代码，但没有 XPU 上的 CI；vLLM-Omni 的模型声明最多，有 Intel 队列，但文档前后矛盾。}
\end{frame}

\begin{frame}[t]{Intel GPU（二）：证据与上游动态}
\scriptsize\def\codesize{\relax}
\renewcommand{\arraystretch}{1.12}
\setlength{\tabcolsep}{3pt}
\begin{tabularx}{\textwidth}{>{\bfseries}p{2.1cm} Y Y}
\toprule
\rowcolor{Mist}
项目 & \textcolor{Teal}{M*} & vLLM-Omni \\
\midrule
公开证据 & 没有日志或基准。PR \#221 的标题写 “validated on Intel GPU(XPU)”，只是标题 & 文档有 B70 测量数字（\code{denoise\_graph.md:100} 等）；XPU 端到端测试只有 Qwen3-TTS 与 Cosmos3；仓库内无日志 \\
近期合并 & 上游 pin 之后：\#1124（分布式上下文并行，11-12）改了 \code{attn/xpu.py}、\code{communication.py} 等 & pin 前已有 \#10124、\#9613；之后：\#10278（11-20，HunyuanImage-3 在 XPU 上运行）、\#10780（11-20，1 卡与 2 卡 PR 任务） \\
开放 PR & 没有标题含 XPU 的开放 PR；\#1413 标题是 Rust 前端，但改了 \code{docker/Dockerfile.xpu}（+4 / -1） & 11 个 XPU / Intel 标题的开放 PR，例如 \#10293（Intel CI 升级）、\#9314（Wan VAE） \\
\bottomrule
\end{tabularx}
\memo{开放 PR 数取 11-13 的上游状态，未重核；合并状态以各 checkout 为准。B70 上 BAGEL 的本机实测见第 3 部分（日志入库于 \code{results/}）。}
\takeaway{Intel 相关的活跃工作主要在 vLLM-Omni；M* 的 XPU 工作已经合并，目前没有开放的 XPU 功能 PR。}
\end{frame}

\begin{frame}[t]{Intel GPU（三）：M* 已有实现（一）：执行路径}
\scriptsize\def\codesize{\relax}
\renewcommand{\arraystretch}{1.12}
\setlength{\tabcolsep}{3pt}
\begin{tabularx}{\textwidth}{>{\bfseries}p{2.2cm} Y Y p{3.6cm}}
\toprule
\rowcolor{Mist}
层 & 做了什么 & 怎么实现 & 代码位置 \\
\midrule
平台识别 & 进程判断自己跑在 NVIDIA、AMD、Intel 或 CPU 上 & 先查 CUDA，再查 XPU；导入模块时不占用设备 & \code{platforms/\_\_init\_\_.py:40, 68}；\code{interface.py:184} \\
注意力（分页 KV） & 后端名 \code{xpu\_paged}，是 M* 唯一的 XPU 注意力 & 调用 vllm-xpu-kernels 的变长 FlashAttention；块表来自 KV 资源；图捕获时预分配计划缓冲区 & \code{attn/xpu.py:19, 232}；\code{attn/base.py:352--373} \\
注意力默认值 & \code{auto} 在 XPU 上也解析为 FlashInfer，代码里没有 XPU 分支 & 其他模型必须在配置里显式写 \code{xpu\_paged}，否则落到 FlashInfer & \code{attn/base.py:65--73} \\
图捕获 & \code{XPUGraph} 后端 & 与 CUDA 图共用接口，按设备类型选后端；BAGEL 的 XPU 配置默认关闭 & \code{accelerator\_graph\allowbreak\_backend.py:195--213}；\code{configs/bagel\allowbreak\_xpu\_tp2.yaml:3} \\
RoPE 与 RMSNorm & XPU 版内核 & 注册为 XPU 内核；残差加归一化的融合版只在 CUDA 上，XPU 上是非融合写法 & \code{position/rope.py:77}；\code{rms\_norm/rms\_norm.py:63, 131} \\
采样 & 批量采样，每个请求用设备端独立随机数 & 调用 vllm-xpu-kernels 的采样核（需 0.1.15 及以上） & \code{sampler/utils.py:1712, 1881} \\
\bottomrule
\end{tabularx}
\memo{以 M* 2c8ffde（2026-11-15，当前 main）为准。XPU 相关的合并提交：\#220（BAGEL TP2）、\#221（CFG TP2 与共享内存 KV 交接）、\#262（XPU 图）、\#347（采样）、\#704（镜像）、\#784（RoPE 表类型拒绝）。}
\takeaway{M* 的 XPU 实现集中在 BAGEL 这一条路径；其余模型的注意力、Triton 核和图捕获都没有在 XPU 上走过。}
\end{frame}

\begin{frame}[t]{Intel GPU（三）：M* 已有实现（二）：通信、配置与交付}
\scriptsize\def\codesize{\relax}
\renewcommand{\arraystretch}{1.12}
\setlength{\tabcolsep}{3pt}
\begin{tabularx}{\textwidth}{>{\bfseries}p{2.2cm} Y Y p{3.8cm}}
\toprule
\rowcolor{Mist}
层 & 做了什么 & 怎么实现 & 代码位置 \\
\midrule
通信组 & 张量并行组在 XPU 上可用 & 组内同时建 gloo（CPU）与 xpu 两个后端；对称内存全归约只在 CUDA 上启用 & \code{communication.py:586--587}；\code{all\_reduce.py:59} \\
同机传输 & 共享内存池有 XPU 分支 & 没有拷贝流 & \code{arena.py:204} \\
BAGEL 配置 & TP2；CFG 并行 TP2 & TP2：LLM 在 rank 1、2；CFG 版本共 7 张卡（0--6） & \code{bagel\_xpu\_tp2.yaml}；\code{bagel\_xpu\_cfg\_tp2.yaml} \\
依赖与镜像 & \code{bagel\_xpu} 额外依赖；Intel 镜像 & 钉住 vllm-xpu-kernels 0.1.15.4；镜像钉住驱动与 IGC，torch 2.13.0+xpu，但克隆 main、不锁 commit & \code{pyproject.toml:173--186}；\code{Dockerfile.xpu:5, 18--20, 81, 90} \\
测试 & XPU 单测与 BAGEL 测试 & 没有 XPU 设备时跳过；M* 的 CI 只在 ubuntu-latest 上跑 & \code{test\_xpu\_attention.py:349}；\code{ci.yml:16} \\
\bottomrule
\end{tabularx}
\memo{以 M* 2c8ffde（2026-11-15）为准。较 9c6f1ad 新增：共享内存池的 XPU 分支（\#1132，11-06）；DCP（\#1124，11-12）不含 XPU；\#1413 改了 \code{Dockerfile.xpu}。}
\takeaway{交付物齐全；缺的是 BAGEL 以外模型的配置，以及 XPU 硬件 CI。}
\end{frame}

\begin{frame}[t]{Intel GPU（四）：M* 上 7 个模型还差什么}
\scriptsize\def\codesize{\relax}
\renewcommand{\arraystretch}{1.12}
\setlength{\tabcolsep}{3pt}
\begin{tabularx}{\textwidth}{>{\bfseries}p{2.6cm} Y Y p{2.2cm}}
\toprule
\rowcolor{Mist}
模型 & 代码上已有 & 模型自身还缺什么 & 工作量（人周） \\
\midrule
BAGEL（T2I、I2I、I2T） & TP2、CFG-TP2 配置；本机实测 B70 上 M* 与 vLLM-Omni 都能跑，对比与日志见第 3 部分 & 图捕获：TP2 下捕获失败、TP1 下显存不足（本机实测）；2 卡 CFG 并行配置需自己写 & 1--2 \\
Qwen3-Omni（文本、语音） & 没有 XPU 配置；MoE 与音频解码器的滑窗核都是 Triton，代码无 XPU 分支 & 写配置（主干注意力切 \code{xpu\_paged}，确认无滑窗层）；验证 Triton 核 & 1--2 \\
Qwen3-TTS & 深度帧 Gluon 核只在 Hopper 启用，XPU 走 PyTorch 回退（\code{depth\_frame/\_\_init\_\_.py:60}） & 写配置；验证滑窗编解码器与 Triton 核 & 1 \\
Pi0.5（VLA） & 注意力默认 FlashInfer（\code{flashinfer\_backend: fa2}），XPU 上不可用 & 切到 \code{xpu\_paged}；核对前缀 KV 与动作专家的非因果注意力 & 1--2 \\
Cosmos3（T2I、T2V、I2V、V2V、机器人动作） & 配置普遍开图捕获，部分用 SP；稠密与块因果注意力只有 FlashInfer 或 FA3，都是 CUDA 专用 & SP 配置改用 TP（文档写明 XPU 不支持 SP）；另需下页的稠密与块因果注意力 & 2--4 \\
V-JEPA 2-AC（世界模型） & 主要用 PyTorch SDPA，依赖少 & 跑通、比对结果、图捕获 & 0.5--1 \\
Orpheus（TTS） & SDPA 与编译，无 CUDA 图；有 TP2 配置 & 跑通；TP2 的 XPU 通信与结果比对 & 1--2 \\
\bottomrule
\end{tabularx}
\memo{工作量为单人估计，未在 XPU 上试跑；不含下页的共性项。}
\takeaway{BAGEL 已通；V-JEPA 2-AC 与 Orpheus 依赖最少；Cosmos3 是最大的单项。}
\end{frame}

\begin{frame}[t]{Intel GPU（五）：共性缺口}
\scriptsize\def\codesize{\relax}
\renewcommand{\arraystretch}{1.12}
\setlength{\tabcolsep}{3pt}
\begin{tabularx}{\textwidth}{>{\bfseries}p{2.6cm} Y Y p{2.2cm}}
\toprule
\rowcolor{Mist}
缺口 & 现状与代码依据 & 要做什么 & 工作量（人周） \\
\midrule
分页注意力功能不全 & 滑窗、\code{sm\_scale}、块因果、稠密、MLA、上下文并行被拒绝（\code{attn/base.py:193--206, 261}）；封装接口有 \code{window\_size}、\code{softmax\_scale}（\code{vllm/\_xpu\_ops.py:1020}） & 接入滑窗与 scale；块因果与稠密另做；上下文并行不做 & 滑窗与 scale 1--2；块因果与稠密 3--5 \\
Triton 核 & 约 40 个文件；Gluon 核为 NVIDIA 专用；调优表只有 NVIDIA（\code{mstar/kernels/tables/}） & 逐个做正确性检查，补 XPU 调优或回退 & 2--4 \\
图捕获 & XPU 图后端已接（\code{accelerator\_graph\_backend.py:195}）；BAGEL 默认关 & 逐模型开启、比对、测性能 & 2--3 \\
观测与剖析 & 实验框架用 nvidia-smi、nsys；NVTX 在 XPU 禁用（\code{docs/serving.rst:87}） & 接入 xpu-smi 与 XPU 剖析器 & 1--2 \\
CI 与版本锁定 & M* CI 不跑 XPU；\code{Dockerfile.xpu:81} 克隆 main & 自托管 B70 runner；锁定版本 & 1--2 \\
\bottomrule
\end{tabularx}
\memo{工作量为单人估计，未在 XPU 上试跑。\code{window\_size}、\code{softmax\_scale} 来自上游封装，M* 锁定的 0.1.15.4 需在 B70 核对。}
\takeaway{共性里最大的两项是注意力功能不全和 Triton 核验证。}
\end{frame}

\begin{frame}[t]{Intel GPU（六）：工作量合计与下一步}
\scriptsize\def\codesize{\relax}
\renewcommand{\arraystretch}{1.12}
\setlength{\tabcolsep}{3pt}
\begin{tabularx}{\textwidth}{>{\bfseries}p{4.2cm} p{2.6cm} Y}
\toprule
\rowcolor{Mist}
项 & 工作量（人周） & 说明 \\
\midrule
7 个模型自身 & 7.5--14 & 上两页模型表的合计 \\
共性项 & 10--18 & 注意力功能、Triton 核、图捕获、观测、CI \\
\textbf{合计} & \textbf{约 18--32（约 4--7 人月）} & 单人估计；按模型重叠扣减后取这一区间 \\
其中 Cosmos3 相关 & 5--9 & 模型自身 2--4，加上稠密与块因果注意力 3--5 \\
不含 Cosmos3 & 约 13--23 & 可作为先行范围 \\
\bottomrule
\end{tabularx}
\memo{这些数字是估计。最快的收敛办法：用 B70 把 7 个模型的 YAML 各改一行注意力后端（\code{xpu\_paged}）跑冒烟测试，记下第一个失败点，约 2--3 天即可把估计落到具体问题。}
\takeaway{建议顺序：先做依赖少的 V-JEPA 2-AC、Orpheus、Pi0.5，再做 Qwen3-Omni 与 Qwen3-TTS，Cosmos3 放最后。}
\end{frame}

\begin{frame}[t]{模型支持：本报告涉及的模型}
\scriptsize\def\codesize{\relax}
\renewcommand{\arraystretch}{1.12}
\setlength{\tabcolsep}{3pt}
\begin{tabularx}{\textwidth}{>{\bfseries}p{3.2cm} Y Y}
\toprule
\rowcolor{Mist}
模型 & \textcolor{Teal}{M*} & vLLM-Omni \\
\midrule
Qwen3-Omni（文本与语音） & 支持 & 支持；文档标 Intel \\
Qwen3-TTS & 支持，三种变体 & 支持；CustomVoice 标 Intel \\
BAGEL（T2I、I2I、I2T） & 支持；Intel 只有文档声明（TP2 配置） & 支持；文档标 Intel \\
Pi0.5（VLA） & 支持，最多 3 路相机 & 支持；文档标 Intel \\
Cosmos3（视频 / 世界模型） & 支持：T2I、T2V、I2V、V2V，机器人动作 & 支持；文档列四种硬件 \\
V-JEPA 2-AC（世界模型） & 支持 & 没有 \\
Orpheus（TTS） & 支持；采样重归一化，与参考实现不完全一致 & 没有 \\
\bottomrule
\end{tabularx}
\takeaway{模型覆盖上 M* 和 vLLM-Omni 最接近；BAGEL、Pi0.5、V-JEPA 2-AC、Orpheus 只有 M* 或只有 vLLM-Omni 有。}
\end{frame}

\begin{frame}[t]{架构规模：注册名数量不能直接比较}
\scriptsize\def\codesize{\relax}
\renewcommand{\arraystretch}{1.12}
\setlength{\tabcolsep}{3pt}
\begin{tabularx}{\textwidth}{>{\bfseries}p{2.1cm} p{3.5cm} Y}
\toprule
\rowcolor{Mist}
仓库 & 注册规模 & 主要家族 \\
\midrule
M* & 60 个注册名（54 个字面量，加 Qwen3.5 系列生成的 5 个和 Qwen3.8-27B 1 个），38 个类；模型都是自己维护 & LLM：qwen3\_5、glm52；扩散：wan22、ltx2\_5、flux2\_klein；TTS：orpheus、kokoro、cosyvoice3；ASR：whisper、parakeet；世界模型：vjepa2、waypoint \\
vLLM-Omni & \code{\_OMNI\_MODELS} 登记 107 项（自有 105，与上游重叠 2 个）；\code{\_DIFFUSION\_MODELS} 登记 96 项，另有 1 个别名（89 个不同的类） & Omni：qwen2\_5\_omni、qwen3\_omni；扩散：QwenImage、Wan、LTX2、Flux；TTS：voxtral\_tts；世界模型：LingBot-World \\
\bottomrule
\end{tabularx}
\takeaway{vLLM-Omni 的登记数只计自己的登记表，上游 vLLM 的模型不计入。以模型表和本报告的实测为准。}
\end{frame}

% 新模型接入（正文）：讲清楚接入是什么、两个框架各要做什么、投入差在哪里。
% 文件与行号以各页的代码位置为准；指南与代码的差异见“常见问题”一页。
% 依据：各 checkout 的 docs/（官网指南的构建源）与 external_latest 最新 main 源码；历史投入来自近期接入提交的 git show --numstat。
\setpart{模型接入工作量}

\begin{frame}[t]{“接入一个新模型”到底是什么？}
\small
“能加载权重”不等于“接入完成”。一个模型要能在服务里被用户正常调用，至少要做到下面五件事：
\begin{itemize}
  \item \kw{框架认得它}：注册表里有它的模型类，HF 配置能解析，默认部署配置存在。缺了，服务可能能启动，但用户找不到入口。
  \item \kw{请求能进出}：把 OpenAI 形式的请求变成模型输入，再把输出变回音频、图像或文本。
  \item \kw{执行顺序对}：编码器、LLM、解码器、声码器等组件，按照图或 stage 的顺序执行；循环与停止条件要明确。
  \item \kw{状态不串、可取消}：每个请求自己的 KV、采样状态、音频缓冲互不干扰；取消、超时、报错后资源能回收。
  \item \kw{有验收}：CPU 单测、GPU 端到端，以及真实服务请求都通过。
\end{itemize}
\takeaway{接入的工作量来自上面五件事里“框架原本没有的部分”，而不是模型文件的数量。}
\end{frame}

\begin{frame}[t]{同一件事，两个框架叫法不同}
\scriptsize\def\codesize{\relax}
\renewcommand{\arraystretch}{1.15}
\setlength{\tabcolsep}{3pt}
\begin{tabularx}{\textwidth}{>{\bfseries}p{2.1cm} Y Y}
\toprule
\rowcolor{Mist}
要做的事 & \textcolor{Teal}{M*} & \textcolor{Orange}{vLLM-Omni} \\
\midrule
让框架认出模型 & \code{MODEL\_REGISTRY}（模型类路径）；需要 HF 下载时加 \code{HF\_MODELS} & \code{\_OMNI\_MODELS} 登记类；\code{OMNI\_PIPELINES} 登记拓扑 \\
描述执行结构 & 模型的“图”（graph）与 walk：一次请求走哪条路 & \code{PipelineConfig} 列出 stage，声明输入来源与终态 \\
每一步的计算单元 & \code{NodeSubmodule}，每个图节点一个 & 每个 stage 的模型类 \\
请求间的状态 & \code{declare\_state}：声明 Array / Sampler / SeqShard & stage input processor：full-payload、async-chunk 或 token-only 三种交接 \\
部署配置 & \code{configs/\allowbreak{}*.yaml} 的 \code{node\_groups}（节点放在哪些 GPU） & deploy YAML（GPU、显存、并行、connector、采样默认值） \\
对外 API 适配 & Python 适配器 + Rust frontend（按类名查找） & \code{tts\_adapters}（TTS 路由）；其余走 serving 层 \\
\bottomrule
\end{tabularx}
\vspace{2pt}
\memo{名字是各框架自己的，但都在回答同一个问题：框架怎样找到模型、怎样执行、怎样管状态、怎样部署、怎样对外。对照时只看“对应的职责”。}
\takeaway{理解了这张表，两家的接入步骤就能对上号。后面的页面都用这里的名字。}
\end{frame}

\begin{frame}[t]{框架提供的模板：接入时要填的契约（一）}
\scriptsize\def\codesize{\relax}
\renewcommand{\arraystretch}{1.14}
\setlength{\tabcolsep}{3pt}
\begin{tabularx}{\textwidth}{>{\bfseries}p{1.9cm} Y Y}
\toprule
\rowcolor{Mist}
契约 & \textcolor{Teal}{M*} & \textcolor{Orange}{vLLM-Omni} \\
\midrule
注册与发现 & \code{MODEL\_REGISTRY} 指向模型类；\code{HF\_MODELS} 给出权重来源；\code{DEFAULT\_CONFIGS} 给出默认配置 & \code{\_OMNI\_MODELS} 登记模型类；\code{OMNI\_PIPELINES} 登记拓扑（值可以是配置，也可以是按 HF config 解析的函数）；\code{\_DIFFUSION\_MODELS} 登记扩散管线 \\
请求进入 & \code{process\_prompt} 把文本与媒体变成张量包 \code{NameToTensorList}；输入顺序由 \code{input\_modalities} 决定 & 请求是 \code{OmniTokensPrompt}：\code{prompt\_token\_ids}、\code{prompt\_embeds}、\code{additional\_information}、多模态数据；\code{prompt\_transform\_func} 与 \code{prompt\_expand\_func} 可改写或展开请求 \\
执行结构 & \code{get\_graph\_walk\_graphs} 按 walk 返回图（\code{GraphNode}、\code{Sequential}、\code{Parallel}、\code{Loop}）；\code{walk\_plan} 决定每个分区走哪些 walk & \code{PipelineConfig} 列出 stage：\code{model\_stage}、\code{execution\_type}、\code{input\_sources}、\code{final\_output}、\code{owns\_tokenizer} \\
阶段间数据 & 边（\code{GraphEdge}）带名字与标志：\code{persist}、\code{is\_streaming}、\code{output\_modality} & stage 的 processor 有三种 hook：\code{full\_payload}、\code{async\_chunk}、\code{token\_only}；内容为 \code{OmniPayload}（键如 \code{hidden\_states}、\code{embed}、\code{codes}、\code{meta}） \\
数据传输 & 工作进程间的张量协议 \code{CommProtocol}：IPC、TCP、RDMA、SHM；KV 可经 Mooncake TransferEngine 传输 & \code{OmniConnectorBase}，方法为 \code{put}、\code{get}、\code{cleanup}；每条边在 YAML 中指定 connector（如 SharedMemory、Mooncake、NIXL） \\
\bottomrule
\end{tabularx}
\vspace{2pt}
\memo{依据：两家 checkout 的接入指南（\code{docs/}）与源码中的类定义。}
\takeaway{两家都提供成套的契约。M* 把执行写成图与张量边，vLLM-Omni 把执行写成 stage 流水线与 payload 交接。}
\end{frame}

\begin{frame}[t]{框架提供的模板：接入时要填的契约（二）}
\scriptsize\def\codesize{\relax}
\renewcommand{\arraystretch}{1.14}
\setlength{\tabcolsep}{3pt}
\begin{tabularx}{\textwidth}{>{\bfseries}p{1.9cm} Y Y}
\toprule
\rowcolor{Mist}
契约 & \textcolor{Teal}{M*} & \textcolor{Orange}{vLLM-Omni} \\
\midrule
输出 & \code{postprocess} 把输出转成客户端字节；\code{EMIT\_TO\_CLIENT} 边标明 \code{output\_modality} & \code{forward} 返回 \code{OmniOutput}；stage 结果为 \code{OmniRequestOutput}，\code{final\_output\_type} 标明模态 \\
请求状态 & \code{declare\_state} 声明 \code{Array}、\code{Sampler}、\code{SeqShard}；每请求的 \code{Layout} 由 \code{process\_prompt} 给出；\code{check\_stop} 判停 & 沿用 vLLM 的调度与 KV 块管理；\code{StagePipelineConfig} 声明 \code{retains\_state\_across\_chunks} 等 \\
对外接口 & \code{OpenAIAdapter} 的功能开关（chat、speech、images 等）决定暴露哪些路由；注册到 \code{ADAPTER\_REGISTRY} 并移植到 Rust frontend；只走 \code{/generate} 的模型不需要适配器 & \code{OmniServingCapability} 与 \code{endpoint\_restrictions} 控制路由；\code{chat\_template} 写在 \code{PipelineConfig}；TTS 用 \code{ARTTSAdapter} 或 \code{DiffusionTTSAdapter} \\
部署配置 & \code{node\_groups} 把节点映射到 GPU rank，可限定 \code{graph\_walks}；\code{resources} 设置 KV 与注意力后端 & \code{StageDeployConfig}：\code{devices}、\code{num\_replicas}、\code{input\_connectors}、\code{output\_connectors}、\code{async\_chunk}、\code{gpu\_memory\_utilization} \\
并行 & 每个 node group 的分片轴：\code{tp}、\code{cfg}、\code{sp}（MoE 另有 \code{ep\_config}） & stage 的 \code{tensor\_parallel\_size}、\code{cfg\_parallel\_size}、\code{sequence\_parallel\_size}，及 Ulysses、Ring 度数 \\
\bottomrule
\end{tabularx}
\vspace{2pt}
\memo{字段名以两家 external\_latest 的 main 为准。版本变化时，需要按同一清单重新核对。}
\takeaway{契约的名字不同，但每一项都必须有对应的实现。漏掉任何一项，模型要么找不到入口，要么在运行时出错。}
\end{frame}

\begin{frame}[t]{框架提供的模板：接入时要填的契约（三）}
\scriptsize\def\codesize{\relax}
\renewcommand{\arraystretch}{1.14}
\setlength{\tabcolsep}{3pt}
\begin{tabularx}{\textwidth}{>{\bfseries}p{1.9cm} Y Y}
\toprule
\rowcolor{Mist}
契约 & \textcolor{Teal}{M*} & \textcolor{Orange}{vLLM-Omni} \\
\midrule
错误与失败 & \code{ValueError} 转为 400（输入不可服务）；模型计划错误为 500；工作进程退出时，等待中的请求返回 503 & 错误语义见 \code{docs/design/module/error\_contracts.md}；connector 接收可设截止时间（\code{get\_with\_deadline}） \\
取消与回收 & 请求移除时引擎调用 \code{cleanup\_request}，删除该请求在节点上的状态 & connector 的 \code{cleanup(request\_id)} 回收数据；未完成的接收用 \code{abandon\_get} 放弃 \\
LoRA & \code{lora\_roots()} 声明适配器；\code{prompt\_lora\_adapters()} 列出可用于某条 prompt 的适配器；客户端指定的适配器经 \code{lora\_adapter} 传入 & 实现位于 \code{vllm\_omni/lora}，设计见 \code{docs/design/feature/lora.md} \\
观测 & \code{mstar/observability/} 记录边的传输与事件（如 \code{EdgeTransfers}） & 指标定义见 \code{docs/design/metrics.md} \\
\bottomrule
\end{tabularx}
\vspace{2pt}
\memo{并发与取消时的状态隔离，两家都需要单独的测试（见“接入完成”的验收门）。}
\takeaway{这些项目不影响单次输出，却决定服务在失败、取消和多适配器场景下是否正确。}
\end{frame}

\begin{frame}[t]{阶段之间的数据：两边都要回答“传什么、谁来搬、何时结束”}
\small
\begin{columns}[T,onlytextwidth]
  \column{.48\textwidth}
  \textbf{M*：边就是数据}
  \begin{itemize}\setlength\itemsep{3pt}
    \item 边是具名的张量包（\code{NameToTensorList}），边名即键。
    \item 边的标志定生命周期：\code{persist} 留在 conductor；\code{is\_streaming} 让 token 在目标缓冲中累积。
    \item 图内数据流由图声明：\code{img\_emb} 发给 \code{LLM}；去噪 latents 在步间传递。
    \item 工作进程之间直接路由张量；conductor 推进 walk 并发送输入。
    \item 传输协议由部署决定（IPC、TCP、RDMA、SHM）。
  \end{itemize}
  \column{.48\textwidth}
  \textbf{vLLM-Omni：processor 与 connector}
  \begin{itemize}\setlength\itemsep{3pt}
    \item 阶段间内容是 \code{OmniPayload}（TypedDict）；async chunk 用 \code{OmniPayloadStruct}。
    \item processor 决定发什么：\code{full\_payload} 在工作进程侧打包，\code{async\_chunk} 按块发出，\code{token\_only} 只分配下游槽位。
    \item connector 决定怎么运：\code{put} 与 \code{get}，默认经 \code{OmniSerde}（msgpack）序列化；每条边可选不同 connector。
    \item KV 交接由 \code{omni\_kv\_config} 声明，由 \code{OmniKVTransferManager} 执行。
  \end{itemize}
\end{columns}
\takeaway{两边都要回答同样的问题：边上传什么，谁负责搬运，何时结束。M* 的答案写在图的边里，vLLM-Omni 的答案写在 processor 与 connector 里。}
\end{frame}

\begin{frame}[t]{先判断新模型属于哪一层}
\small
按下面三个问题往下走，就能判断这次接入大概是哪一层：
\begin{enumerate}
  \item \kw{架构和任务已经能跑吗？}（同样的架构、同样的任务、同样的输出形式，只是换一个 checkpoint）\\
  是 → \kw{L0}：只补配置、权重版本和文档。
  \item \kw{框架的 core 里已有底层模型类吗？}（vLLM 的 AR 模型类）\\
  是 → \kw{L1}：写模型类和接入层的代码；否 → \kw{L2}：还要补底层模型本身。
  \item \kw{是多阶段、流式、带循环，或者是 diffusion 吗？}\\
  是 → 在 L1 或 L2 之上再加 \kw{L3} 的工作：stage 拓扑、chunk 交接、循环与停止、取消与资源回收。
\end{enumerate}
\vspace{4pt}
\memo{同一个模型可以跨层。例：TTS 模型的 LLM 部分已有 AR 实现（L1），但声码器和流式交接是新的（L3）。}
\takeaway{先判断层级，再查框架是否有现成的同类路径。层级比框架名更能决定投入。}
\end{frame}

\begin{frame}[t]{四个层级：额外要做的事与为什么贵}
\scriptsize\def\codesize{\relax}
\renewcommand{\arraystretch}{1.16}
\setlength{\tabcolsep}{3pt}
\begin{tabularx}{\textwidth}{>{\bfseries}p{1.6cm} Y Y}
\toprule
\rowcolor{Mist}
层级 & 额外要做的事 & 为什么贵（或不贵） \\
\midrule
L0 \newline 只换 checkpoint & 权重版本、部署配置、采样默认值、文档、smoke test & 没有新执行代码。风险主要在配置和 prompt 模板是否对得上。 \\
L1 \newline 新模型类，底层 AR 已有 & 模型包、stage 或 pipeline 配置、请求构造、payload、YAML & 框架已提供调度、KV、connector，主要是把新模型“接进”现成流程。 \\
L2 \newline 新 AR 架构 & 在 L1 的基础上，还要补底层模型：权重映射、core 模型、多模态接口 & 框架原本不认识这个结构，底层执行语义要自己实现。 \\
L3 \newline 多阶段 / 流式 / 循环 / diffusion & stage 拓扑或图、chunk 协议、循环步数与停止、每请求状态、取消回收、可能的共享基类 & 新增的是执行语义，而不是文件数。问题在于“框架原本没有这种形状”。 \\
\bottomrule
\end{tabularx}
\vspace{2pt}
\memo{规模参考见下一页。规模是整个 PR 的行数，含测试和文档，不是人天。}
\takeaway{最贵的通常不是“写模型”，而是“框架原本没有这个执行形状”。}
\end{frame}

\begin{frame}[t]{近期接入提交的规模（代理指标）}
\scriptsize\def\codesize{\relax}
\renewcommand{\arraystretch}{1.12}
\setlength{\tabcolsep}{3pt}
\begin{tabularx}{\textwidth}{>{\bfseries}p{1.2cm} p{1.2cm} Y Y Y}
\toprule
\rowcolor{Mist}
层级 & 框架 & 例子 & 文件（新增 / 修改） & 行数（+/−） \\
\midrule
L0 & vLLM-Omni & Fish S2 Pro，仅部署 YAML & 2（1/1） & +99 \\
L1 & M* & Ripple-65B，复用 MultimodalDecoderModel & 模型包 5 & +407（模型包） \\
L1 & vLLM-Omni & Haru-4B，已有 vLLM 类 & 28（28/5） & +3815/0 \\
L2 & M* & Fish Audio S2-Pro，无 scaffold，双 AR TTS & 26（13/13） & +4028/−1 \\
\bottomrule
\end{tabularx}
\vspace{6pt}
\memo{表格续下页。}
\end{frame}
\begin{frame}[t]{近期接入提交的规模（代理指标）（续）}
\scriptsize\def\codesize{\relax}
\renewcommand{\arraystretch}{1.12}
\setlength{\tabcolsep}{3pt}
\begin{tabularx}{\textwidth}{>{\bfseries}p{1.2cm} p{1.2cm} Y Y Y}
\toprule
\rowcolor{Mist}
层级 & 框架 & 例子 & 文件（新增 / 修改） & 行数（+/−） \\
\midrule
L2 & vLLM-Omni & YuE2-3B，新模型类与音频接口 & 27 & +3294/-6 \\
L3 & M* & Fun-CosyVoice3，流式 TTS 多 stage & 27（15/12） & +5349/−11 \\
L3 & vLLM-Omni & ABot-World，实时 AR-Diffusion & 21 & +4812/-6 \\
\bottomrule
\end{tabularx}
\vspace{1pt}
\memo{M* 的 Qwen3.8-27B（同一代码的新 checkpoint，9 个文件）未列入本页。}
\takeaway{同一层级里，行数差异主要来自测试、变体和共享层改动；不能直接当作人力。}
\end{frame}

\begin{frame}[t]{怎么读这些行数}
\small
\begin{itemize}
  \item \kw{行数含测试与文档。}两家的样本中，测试通常占 20%–35\%。M* 的 MiniMax-H3 有 21 个测试文件。
  \item \kw{有的提交顺便改了共享层。}例如 Ming3-Audio 改了 vLLM-Omni 自身的核心 24 个文件。这些改动不是模型自己的工作，但会计入 PR。
  \item \kw{M* 的历史提交用的是旧契约。}这些模型引入时用的是 \code{get\_node\_resources}，之后迁移到 \code{declare\_state} 的成本没有计入。
  \item \kw{同样的行数，难度不同。}一个 TTS 的流式交接和一个纯 AR 模型的 KV 配置，行数可能差不多，难度却差很多。
\end{itemize}
\takeaway{行数只能说明“这类改动大概有多大”，不能用来比较两个框架谁更省事。}
\end{frame}

\begin{frame}[t]{M*：接入的七个步骤}
\small
\begin{enumerate}
  \item \kw{登记模型。}在 \code{MODEL\_REGISTRY} 里把名字指向模型类；要从 Hugging Face 下载时加 \code{HF\_MODELS}（YAML 的 \code{model\_path\_hf} 可以覆盖）。
  \item \kw{写 Model 类。}四个必需方法：\code{get\_graph\_walk\_graphs}（画出模型的图）、\code{walk\_plan}（决定一次请求走哪条路）、\code{process\_prompt}（请求变输入）、\code{get\_submodule}（为每个节点造执行单元）。
  \item \kw{为每个节点写 NodeSubmodule。}\code{prepare\_inputs} 和 \code{forward} 必需；有 KV 或采样状态的节点要声明 \code{declare\_state}。
  \item \kw{写配置。}\code{DEFAULT\_CONFIGS} 加名字；\code{configs/\allowbreak{}<name>.yaml} 用 \code{node\_groups} 决定节点放在哪些 GPU；\code{docs/\allowbreak{}models.rst} 加一行。
  \item \kw{需要 OpenAI 接口时：}Python 适配器登记；Rust frontend 按 Python 类名查找适配器。
  \item \kw{过门禁。}包要加入 \code{test\_state\_lint.py} 的 \code{PORTED}；必须覆写 \code{request\_kwargs}（CI 会检查）。
  \item \kw{验证。}CPU 测试自动运行；GPU 集成测试和真实的 \code{mstar-serve} 请求通过。
\end{enumerate}
\takeaway{M* 最需要设计的是第 2、3 步：把模型描述成图，再把每个请求的状态声明清楚。}
\end{frame}

\begin{frame}[t]{M*：几个核心概念}
\scriptsize\def\codesize{\relax}
\renewcommand{\arraystretch}{1.16}
\setlength{\tabcolsep}{3pt}
\begin{tabularx}{\textwidth}{>{\bfseries}p{1.9cm} Y Y}
\toprule
\rowcolor{Mist}
概念 & 通俗解释 & 你要写什么 \\
\midrule
图（Graph） & 模型是一张有向图：节点是计算单元，边是数据流向。 & \code{get\_graph\_walk\_graphs} 返回图的描述。 \\
Walk & 一次请求沿图走的一条路径。同一模型，不同请求可以走不同路径（例如文生图与图生文）。 & \code{walk\_plan} 按请求选择 walk。 \\
节点（NodeSubmodule） & 图里一个计算节点的执行单元：准备输入、执行 forward、判断停止。 & 每个节点一个类，实现 \code{prepare\_inputs}、\code{forward}。 \\
状态（declare\_state） & 请求期间需要保留的数据：KV、采样器、按序列分片的数组。框架靠它做复用、回收和停止。 & 节点声明自己的状态；不要自己管理 stream。 \\
Loop 与 stop & 图里可以有循环（如扩散去噪、AR 解码）。循环上限和停止条件决定何时结束。 & \code{get\_loop\_caps} 定步数；\code{check\_stop} 判停止。 \\
\bottomrule
\end{tabularx}
\vspace{2pt}
\memo{代码位置：Model 基类 \code{mstar/\allowbreak{}model/\allowbreak{}base.py}；节点基类 \code{mstar/\allowbreak{}model/\allowbreak{}submodule\_base.py}。}
\takeaway{M* 的写法是“先把模型画成图，再填每个节点”。框架因此能知道每个请求的状态与停止点。}
\end{frame}

\begin{frame}[t]{M*：什么时候能套用现成 scaffold，什么时候要改 runtime}
\scriptsize\def\codesize{\relax}
\renewcommand{\arraystretch}{1.16}
\setlength{\tabcolsep}{3pt}
\begin{tabularx}{\textwidth}{>{\bfseries}p{2.1cm} Y Y}
\toprule
\rowcolor{Mist}
模型形状 & 可套用的 scaffold & 仍需手写 \\
\midrule
解码器 LLM + 视觉 / 音频编码器 & \code{MultimodalDecoderModel}（6 个包在用） & 模型类、编码器权重映射、prompt 模板、media 预处理 \\
flow-matching 去噪循环（DiT） & \code{DenoiseLoopSubmodule}（5 个包在用） & transformer 与 VAE 移植、各钩子（\code{seed\_latents}、\code{denoise} 等）、质量验证 \\
其它（如 Qwen3-Omni、BAGEL、LingBot-World、Cosmos3） & 不在上面两个 scaffold 上 & 按家族从头写，或使用各自的执行路径 \\
需要新 kernel、新 transport 或新资源类型 & 无 & 要改 M* runtime 本身，投入不再局限于模型包 \\
\bottomrule
\end{tabularx}
\vspace{2pt}
\memo{scaffold 的包数来自仓库内的 grep 统计；“在用”指当前 checkout 中有包继承它。}
\takeaway{形状与现有 scaffold 匹配时，M* 的接入明显更快；形状不匹配时，投入会落到 runtime 中。}
\end{frame}

\begin{frame}[t]{vLLM-Omni：接入的八个步骤}
\small
\begin{enumerate}
  \item \kw{HF 配置。}transformers 认识该 \code{model\_type} 吗？不认识，就在 \code{\_register\_omni\_hf\_configs()} 里加一项。
  \item \kw{模型类。}复用 vLLM 已有的层与 weight loader。Omni 统一入口按 \code{model\_stage} 分派到各 stage 类（参考 Qwen3-Omni）。
  \item \kw{登记类。}每个类在 \code{\_OMNI\_MODELS} 里登记一条。安装到 vLLM 的动作由 \code{register\_omni\_models\_to\_vllm()} 完成。
  \item \kw{pipeline。}\code{PipelineConfig} 列出各 stage：用哪个架构、怎样执行、输入来自哪里、哪个是终态。在 \code{OMNI\_PIPELINES} 登记。
\end{enumerate}
\vspace{6pt}
\memo{续下页。}
\end{frame}

\begin{frame}[t]{vLLM-Omni：接入的八个步骤（续）}
\small
\setcounter{enumi}{4}
\begin{enumerate}
  \item \kw{stage input processor。}定义 stage 之间的数据交接：full-payload、async-chunk 或 token-only。
  \item \kw{deploy YAML。}GPU 放置、显存、并行、connector、采样默认值。顶层 \code{pipeline:} 键可以指定 pipeline。
  \item \kw{对外接口（如有 TTS）。}新增 adapter 文件，在 \code{tts\_adapters/\allowbreak{}\_\_init\_\_.py} 加一行 import。不改 \code{serving\_speech.py}。
  \item \kw{测试与文档。}模型、pipeline、processor、online 测试；\code{recipes/} 与 supported\_models 各加一行。
\end{enumerate}
\takeaway{vLLM-Omni 的工作重心在 pipeline 和 stage contract；底层模型类只有 vLLM core 没有时才是主项。}
\end{frame}

\begin{frame}[t]{vLLM-Omni：一个模型怎样被选中 pipeline}
\small
部署时，框架按下面的顺序找 pipeline，找到就停：
\begin{enumerate}
  \item deploy YAML 顶层的 \code{pipeline:} 键（用户或 recipe 指定的，优先级最高）。
  \item checkpoint 的 HF \code{model\_type} 在 \code{OMNI\_PIPELINES} 里有映射：值是配置，或一个按 HF config 解析的函数（resolver）。
  \item 按 HF \code{architectures} 匹配 \code{PipelineConfig} 的 predicate。
\end{enumerate}
\vspace{4pt}
resolver 的作用是“同一个架构，不同变体走不同 pipeline”。例：Qwen3-Omni 按 \code{enable\_audio\_output} 决定是否带 talker；只输出文本时走 thinker-only。
\takeaway{新模型如果有多个变体，要想清楚每个变体该选哪条 pipeline，并在 resolver 里写出来。}
\end{frame}

\begin{frame}[t]{vLLM-Omni：TTS 是“两 stage 或单 stage”加一个 adapter}
\small
\begin{itemize}
  \item \kw{两种形态}：两个 stage（AR codec predictor → decoder），或单 stage AR streaming。
  \item \kw{对外接口}：新增 \code{ARTTSAdapter} 子类，实现 \code{validate} 与 \code{build}，声明 \code{stage\_keys}。
  \item \kw{在哪里改}：一个 adapter 文件，加 \code{tts\_adapters/\_\_init\_\_.py} 的一行 import；不改 \code{serving\_speech.py}。
  \item \kw{纯 diffusion 的 TTS}也走 adapter：OmniVoice 继承 \code{DiffusionTTSAdapter}。
  \item \kw{流式}：full-payload 与 async-chunk 两条路径都要测。
\end{itemize}
\vspace{4pt}
\memo{历史样例：Breeze-TTS-2（两 stage）。单 stage 的 Kokoro-82M 已于 11-16 移除，提交仍在历史中。}
\takeaway{TTS 的 endpoint 是小增量，但不替代模型自身的 stage、decoder 与流式实现。}
\end{frame}

\begin{frame}[t]{vLLM-Omni：diffusion 是完整的模型后端工作}
\small
\begin{itemize}
  \item \kw{要移植的}：transformer 与 pipeline；去掉 Diffusers mixin，换成 Omni 的 attention 与 config。
  \item \kw{注册}：在 \code{\_DIFFUSION\_MODELS} 里登记 (folder, file, class)。
  \item \kw{pre / post process}：可选。指南写成必需，代码里是可选的。
  \item \kw{必测}：L4 functionality，并加 \code{test\_alignment.py} 的对齐门；与 Diffusers 的对照是手动 PR 项。
  \item \kw{独立扩展}：TP / CFG / SP、step execution、streaming、cache、offload，不算“基础跑通”。
\end{itemize}
\vspace{4pt}
\memo{所列三个 diffusion 提交（Krea 2、Flux3 base、Stable Audio 3）都没有 stage processor 和 deploy YAML；并非所有 diffusion 提交都如此。}
\takeaway{新 diffusion port 通常是完整模型后端的工作，不能按 TTS 的小增量估计。}
\end{frame}

\begin{frame}[t]{节点、walk、stage：两边的执行单位粒度不同}
\scriptsize\def\codesize{\relax}
\renewcommand{\arraystretch}{1.16}
\setlength{\tabcolsep}{3pt}
\begin{tabularx}{\textwidth}{>{\bfseries}p{2.0cm} Y Y}
\toprule
\rowcolor{Mist}
层次 & \textcolor{Teal}{M*} & \textcolor{Orange}{vLLM-Omni} \\
\midrule
最小执行单元 & 节点（\code{NodeSubmodule}）：一个 \code{torch.nn.Module}，对应图中的一个节点，在同一运行时内执行 & 模型类的 \code{forward}，由 vLLM 引擎在某个 stage 内调度 \\
运行单位 & walk：一次请求沿图走的一条路径，如 \code{prefill\_text}、\code{prefill\_vit}、去噪循环 & stage：独立的引擎实例，有自己的设备、并行度与显存预算 \\
循环 & \code{Loop} 由图声明；循环上限由 \code{get\_loop\_caps} 给出，停止由 \code{check\_stop} 判断 & 去噪或解码的循环在 stage 内部完成；stage 之间由 \code{input\_sources} 连接（BAGEL 为 0 → 1） \\
跨单元交接 & 同一图内是张量边；跨工作进程由通信协议搬运 & stage 之间经 processor 与 connector 交接 \\
部署归属 & \code{node\_groups} 把节点组放到 rank 上，节点组可跨多个 rank（分片） & \code{devices} 把 stage 放到卡上；一个 stage 也可以占多张卡（TP、CFG 并行） \\
\bottomrule
\end{tabularx}
\vspace{2pt}
\memo{两边没有一对一的对应。M* 的节点与 walk 位于同一运行时内；vLLM-Omni 的 stage 是独立引擎，其内部计算不以节点表示。比较时对齐职责，不对齐名字。}
\takeaway{M* 的图在节点层面可见，vLLM-Omni 的流水线在 stage 层面可见。两边的 BAGEL 实现分别落在这两层上。}
\end{frame}

\begin{frame}[t]{同一个模型在两边的落地：BAGEL}
\scriptsize\def\codesize{\relax}
\renewcommand{\arraystretch}{1.14}
\setlength{\tabcolsep}{3pt}
\begin{tabularx}{\textwidth}{>{\bfseries}p{1.9cm} Y Y}
\toprule
\rowcolor{Mist}
 & \textcolor{Teal}{M*} & \textcolor{Orange}{vLLM-Omni（two-stage）} \\
\midrule
结构 & 7 个节点，按 walk 组织：\code{prefill\_text}（\code{LLM}）、\code{prefill\_vit}（\code{vit\_encoder} → \code{LLM}）、\code{prefill\_vae}（\code{vae\_encoder} → \code{LLM}）、去噪循环与 \code{vae\_decoder} & 2 个 stage：Thinker（\code{LLM\_AR}，拥有 tokenizer）→ DiT（\code{DIFFUSION}，\code{input\_sources = (0,)}） \\
视觉与 VAE 的位置 & ViT、VAE 节点与 rank 0 上的 LLM 同卡 & ViT 在 Thinker 内；VAE 在 DiT stage 内 \\
CFG 的表达 & 图节点 \code{LLM\_cfg\_text}、\code{LLM\_cfg\_img}，由 \code{combine\_cfg} 合并；CFG 并行时两个分支放在 rank 1 & 附加请求：\code{expand\_cfg\_prompts} 生成 cfg\_text 的请求；DiT 侧 \code{collect\_cfg\_kv\_caches} 收齐各路 KV \\
跨卡交接 & 张量边；CFG 并行每步同步一次（一张图约 50 次） & KV 经 connector 从 Thinker 发给 DiT（\code{need\_send\_cache}、\code{need\_recv\_cache}） \\
配置文件 & \code{configs/mstar/bagel\_cfgpar\_2gpu.yaml}：\code{node\_groups} 把节点放到 rank 0、1 & \code{configs/vllm\_omni/bagel\_two\_stage\_2gpu.yaml}：\code{stages} 与 \code{connectors} \\
\bottomrule
\end{tabularx}
\vspace{2pt}
\memo{M* 的 walk 与节点见 \code{mstar/model/bagel/bagel\_model.py}；vLLM-Omni 的拓扑见 \code{vllm\_omni/model\_executor/models/bagel/pipeline.py}。}
\takeaway{同一个模型，两边都要解决 ViT、VAE、CFG 的放置与交接。差别在于 M* 把它们写进图，vLLM-Omni 把它们拆成 stage 与 processor。}
\end{frame}

\begin{frame}[t]{同一接入范围下，两个框架的投入是条件性的}
\scriptsize\def\codesize{\relax}
\renewcommand{\arraystretch}{1.14}
\setlength{\tabcolsep}{2.8pt}
\begin{tabularx}{\textwidth}{>{\bfseries}p{2.0cm} Y Y}
\toprule
\rowcolor{Mist}
接入层级 & \textcolor{Teal}{M*} & \textcolor{Orange}{vLLM-Omni} \\
\midrule
L0 新 checkpoint & 低：注册、YAML、文档、门禁。 & 低：deploy YAML、pipeline 别名、recipe。 \\
L1 新模型类，底层 AR 已有 & 中：Model 与 NodeSubmodule；scaffold 可套用时偏低。 & 中：pipeline、stage processor、deploy；复用 vLLM 类。 \\
L2 新 AR 架构 & 高：模型实现；缺资源或 kernel 时进入 runtime。 & 很高：core 模型、Omni 注册、HF config 注册、pipeline。 \\
L3 多阶段 / 流式 TTS & 高：partition、stream edge、chunk、stop。 & 高：stage 与 chunk 契约、adapter、流式验收。 \\
L3 diffusion / 循环 & 中到高（scaffold 可套用时）；否则高。 & 高到很高：transformer 与 pipeline 移植，L4 与对齐门。 \\
\bottomrule
\end{tabularx}
\vspace{2pt}
\memo{“低 / 中 / 高”只表示实现面、跨层注册和验证范围，不是人天。两家没有同一模型的开发记录，所以不做工时对比，也不做总排名。}
\takeaway{没有无条件“最省”的框架：底层 core 的复用度、执行形状与验收范围共同决定总投入。}
\end{frame}

\begin{frame}[t]{每家最容易花时间的地方}
\small
\begin{columns}[T,onlytextwidth]
  \column{.48\textwidth}
  \textbf{M*}
  \begin{itemize}
    \item 把模型画成图，并声明每个请求的状态。
    \item loop、流式 partition 与 stop 的组合。
    \item 形状不匹配 scaffold 时，要改 runtime。
  \end{itemize}
  \column{.48\textwidth}
  \textbf{vLLM-Omni}
  \begin{itemize}
    \item stage 之间的数据交接（processor 与 connector）。
    \item 新 AR 架构时，core 模型要自己实现。
    \item diffusion 的 transformer 移植与对齐测试。
  \end{itemize}
\end{columns}
\takeaway{每个框架的难点不同。计划接入时，先用这一页确认你的模型会落在哪个难点上。}
\end{frame}

\begin{frame}[t]{“接入完成”的验收门，按框架对应的测试}
\scriptsize\def\codesize{\relax}
\renewcommand{\arraystretch}{1.16}
\setlength{\tabcolsep}{3pt}
\begin{tabularx}{\textwidth}{>{\bfseries}p{1.6cm} Y Y}
\toprule
\rowcolor{Mist}
验收门 & \textcolor{Teal}{M*} & \textcolor{Orange}{vLLM-Omni} \\
\midrule
发现与配置 & \code{test/\allowbreak{}modular/\allowbreak{}test\_model\_registration.py}：注册、默认配置、docs 行、HF\_MODELS & 模型注册与 pipeline 解析测试 \\
组件正确性 & \code{test/modular/} CPU 测试：graph、worker、资源生命周期 & 模型加载、权重映射、processor 序列化 \\
执行路径 & \code{test\_model\_contracts.py}：\code{request\_kwargs} 覆写 & 每个 stage 的输入输出与 modality \\
服务与流式 & GPU 集成测试；真实 \code{mstar-serve} 请求 & L4 functionality；TTS adapter 测试；diffusion 对齐门 \\
\bottomrule
\end{tabularx}
\vspace{2pt}
\memo{路径以各 checkout 为准。并发与取消的状态隔离，两家都需要单独的测试；指南中的覆盖并不完整。}
\takeaway{“服务能启动”或单次输出正常都不算通过。每一道门都要有可复现的测试入口。}
\end{frame}

\begin{frame}[t]{接入时的五类常见问题}
\small
\begin{enumerate}
  \item \kw{只加权重，不加注册和文档。}服务能起，用户找不到。M* 的 \code{docs/\allowbreak{}mintlify/\allowbreak{}} 副本没有测试检查。
  \item \kw{照抄旧名。}M* 指南第 5 步还在用 CUDA graph 旧名，当前是 accelerator 名。
  \item \kw{按指南改旧文件。}vLLM-Omni 指南末尾仍写要改 \code{serving\_speech.py}，当前应改 TTS adapter。
  \item \kw{把 diffusion 当成 TTS 的变体。}diffusion 的 pre / post、L4 测试与对齐门都不同，指南对此写得不完整。
  \item \kw{用行数估人力。}行数含测试、文档和共享层改动，不能直接换算成工时。
\end{enumerate}
\takeaway{每一类都可以对照两家 checkout 中的指南与源码核对。}
\end{frame}

\begin{frame}[t]{来源、版本与边界}
\scriptsize\def\codesize{\relax}
\renewcommand{\arraystretch}{1.12}
\setlength{\tabcolsep}{3pt}
\begin{tabularx}{\textwidth}{>{\bfseries}p{1.8cm} p{3.4cm} Y}
\toprule
\rowcolor{Mist}
框架 & \code{external\_latest} checkout & 指南与核对说明 \\
\midrule
M* & \code{main} / 2c8ffde8 / 11-15 & Adding a New Model，更新于 11-15；契约与 checkout 一致，第 5 步有过期名称。 \\
vLLM-Omni & \code{main} / a36de509e / 11-21 & Omni 指南最后修改 615f97b9c（2026-09-06）；Diffusion 11-11、TTS 11-04 沿用旧记录，未在新 checkout 重核。 \\
\bottomrule
\end{tabularx}
\vspace{2pt}
\memo{官网页面本次未能直接抓取（网络策略拦截）。指南内容取自各 checkout 的 docs/，即官网页面的构建源。历史投入来自 squash 提交的 \code{git show --numstat}。未运行测试、未安装依赖、未下载权重。}
\takeaway{这些是最新 main 快照，用于接入流程的核对；性能结果仍以 pinned 版本与 run manifest 为准。}
\end{frame}

% BAGEL 实测：测试环境、配置定义、统一工作量
\setpart{BAGEL 实测 · 环境与配置}

\begin{frame}[t]{测试环境：Intel Arc Pro B70}
\scriptsize\def\codesize{\relax}
\renewcommand{\arraystretch}{1.15}
\begin{tabularx}{\textwidth}{>{\bfseries}p{2.6cm} Y Y}
\toprule
\rowcolor{Mist}
 & \textcolor{Teal}{M*} & \textcolor{Orange}{vLLM-Omni} \\
\midrule
源码 & \code{external/mstar} main 2c8ffde8f（最新提交） & \code{external/vllm-omni} main a36de509e（2026-11-21）；对照版本 v0.32.0 \\
PyTorch / 内核 & torch 2.13.0+xpu；vllm-xpu-kernels 0.1.15.4（M* 的 \code{[bagel\_xpu]} 固定） & torch 2.15.0+xpu；vllm-xpu-kernels 0.1.17（vLLM 0.34.0+xpu 固定） \\
环境 & \code{mstar\_xpu/.venv} & \code{\textasciitilde/workspaces/.venvs/vllm-omni-main}（v0.32 环境保留为 \code{vllm-omni-v0.32.0}） \\
模型 & \multicolumn{2}{Y}{BAGEL-7B-MoT（同一份 BF16 权重），约 27 GiB；两个框架都不做量化、缓存加速或减步} \\
硬件 & \multicolumn{2}{Y}{Arc Pro B70 $\times$3（每张 32 GiB），单 NUMA。两卡结果使用卡 1、2；三卡研究使用卡 0、1、2。运行前检查卡占用，只在空闲卡上运行} \\
\bottomrule
\end{tabularx}
\vspace{4pt}
\begin{itemize}\scriptsize\setlength\itemsep{1pt}
  \item 两个框架各用各自源码要求的 torch 与内核版本，没有强行统一。两边都使用 eager 执行，即不启用 XPU graph；M* 默认开启 \code{torch.compile}，vLLM-Omni 无 compile。
  \item 并行拓扑、compile、缓存的每种组合都是一个命名配置，下两页给出定义。
\end{itemize}
\takeaway{两卡结论适用于卡 1+2 与本次软件版本；三卡结论适用于卡 0、1、2。卡对之间的通信速度不同，其他卡对未测。}
\end{frame}

\begin{frame}[t]{配置定义：两卡（卡 1、2）}
\scriptsize\def\codesize{\relax}
\renewcommand{\arraystretch}{1.1}\setlength{\tabcolsep}{2.5pt}
\begin{tabularx}{\textwidth}{>{\bfseries}p{3.2cm} p{1.0cm} p{3.4cm} Y c c c}
\toprule
\rowcolor{Mist}
配置 & 框架 & 卡与分工 & 说明 & compile & 前缀缓存 & graph \\
\midrule
M* CFG 并行 & \textcolor{Teal}{M*} & 卡 1、2 各一份完整 LLM & 卡 1 算主分支，卡 2 算引导分支；无 TP 通信 & 开 & 开 & 无 \\
M* CFG 并行，缓存关 & \textcolor{Teal}{M*} & 同上 & 仅关闭 KV 前缀缓存 & 开 & 关 & 无 \\
M* CFG 并行，compile 关 & \textcolor{Teal}{M*} & 同上 & hook 关闭 \code{torch.compile}，非官方，仅诊断 & 关 & 开 & 无 \\
M* TP2 & \textcolor{Teal}{M*} & 卡 1、2 切分 LLM；ViT、VAE 在卡 1 & 每层 all-reduce；官方 \code{bagel\_xpu\_tp2} 的两卡版 & 开 & 开 & 无 \\
Omni 单 stage CFG 并行 & \textcolor{Orange}{Omni} & 整个模型一个 stage & 两路 CFG 分支各占一卡；I2I 需要 3 路 & 无 & — & 无 \\
Omni 单 stage TP2 & \textcolor{Orange}{Omni} & 整个模型一个 stage，TP2 & 图像编辑的两卡最优方案 & 无 & — & 无 \\
Omni two-stage & \textcolor{Orange}{Omni} & Thinker 在卡 1，DiT 在卡 2 & KV 经 connector 传给 DiT；Thinker 的 \code{max\_model\_len} 为 6144 & 无 & 关 & 无 \\
Omni two-stage，缓存开 & \textcolor{Orange}{Omni} & 同上 & 仅开启 Thinker 前缀缓存 & 无 & 开 & 无 \\
Omni two-stage，graph & \textcolor{Orange}{Omni} & 同上 & Thinker 启用 graph；扩散仍为 eager & 无 & 关 & 开 \\
\bottomrule
\end{tabularx}
\vspace{3pt}
\memo{两卡上 M* 的最优配置是 CFG 并行（compile 开，缓存开）；vLLM-Omni 文生图用单 stage CFG 并行，图像编辑用单 stage TP2，图生文用 two-stage。}
\end{frame}


\begin{frame}[t]{配置定义：三卡（卡 0、1、2）}
\scriptsize\def\codesize{\relax}
\renewcommand{\arraystretch}{1.1}\setlength{\tabcolsep}{2.5pt}
\begin{tabularx}{\textwidth}{>{\bfseries}p{3.4cm} p{1.0cm} p{3.4cm} Y c c c}
\toprule
\rowcolor{Mist}
配置 & 框架 & 卡与分工 & 说明 & compile & 前缀缓存 & graph \\
\midrule
M* CFG 并行（三卡） & \textcolor{Teal}{M*} & 卡 0 主分支；卡 1、2 各一路引导分支 & 官方三 rank 的 CFG 方案，三卡都是完整模型 & 开 & 开 & 无 \\
M* CFG 并行（三卡），缓存关 & \textcolor{Teal}{M*} & 同上 & 仅关闭 KV 前缀缓存 & 开 & 关 & 无 \\
M* CFG 并行（三卡），graph & \textcolor{Teal}{M*} & 同上 & graph 开启；未打补丁 & 开 & 开 & 开 \\
M* TP2（三卡） & \textcolor{Teal}{M*} & 卡 0 为 ViT、VAE；卡 1、2 切分 LLM & LLM 切分与两卡 TP2 相同 & 开 & 开 & 无 \\
M* TP2（三卡），graph & \textcolor{Teal}{M*} & 同上 & 需对 M* 打补丁，未合入上游 & 开 & 开 & 开 \\
Omni two-stage（三卡） & \textcolor{Orange}{Omni} & Thinker TP2 在卡 0、1；DiT 在卡 2 & 图生文不使用 DiT & 无 & 关 & 无 \\
Omni two-stage（三卡），graph & \textcolor{Orange}{Omni} & 同上 & Thinker 启用 graph；扩散仍为 eager & 无 & 关 & 开 \\
Omni 单 stage CFG 并行（三卡） & \textcolor{Orange}{Omni} & 卡 0、1、2 各一路 CFG 分支 & I2I 的三路 CFG 方案 & 无 & — & 无 \\
\bottomrule
\end{tabularx}
\vspace{3pt}
\memo{三卡配置用于 XPU graph 研究与 I2I 的三路 CFG 比较；两卡结论不依赖三卡数据。}
\end{frame}


\begin{frame}[t]{统一工作量：同一个客户端、同一张输入图、相同的步数与 CFG}
\scriptsize\def\codesize{\relax}
\renewcommand{\arraystretch}{1.15}
\begin{tabularx}{\textwidth}{>{\bfseries}p{2.0cm} Y}
\toprule
\rowcolor{Mist}
case & 定义（\code{scripts/bench\_bagel.py}，两个框架共用） \\
\midrule
T2I 文生图 & prompt \code{A red ceramic cup on a white table}，512$\times$512，50 个时间步，shift 3.0，CFG text 4.0 / img 1.0，作用区间 [0.4, 1.0]，renorm \code{global}（区间内 2 路分支，区间外 1 路） \\
I2T 1 token & 512$\times$512 输入图 + 一句描述问题，贪心解码，只生成 1 个 token（即图像编码 + 预填充，不含解码） \\
I2T 64 token & 同一输入图 + 详细描述问题，贪心解码，\code{ignore\_eos} 生成精确 64 个 token（用 usage 与 tokenizer 校验） \\
I2I 图像编辑 & 同一输入图，指令 \code{Replace the bagel with a croissant.}，输出 512$\times$512，50 步，CFG text 4.0 / img 2.0，区间 [0, 1]，renorm \code{text\_channel}（每步 3 路分支） \\
\bottomrule
\end{tabularx}
\vspace{6pt}
\memo{T2I 与 I2T、I2I 的参数在两个框架上一致。两个框架各自的执行方式（并行拓扑、compile、缓存）见上一页的配置定义。}
\takeaway{四类 case 的请求定义在两个框架上相同，定义全部来自 \code{scripts/bench\_bagel.py}。}
\end{frame}
\begin{frame}[t]{统一工作量（续）：输入 token、缓存与测量方法}
\small
\begin{itemize}\setlength\itemsep{3pt}
  \item \kw{输入图像 token}（以 vLLM-Omni 返回的 \code{usage.prompt\_tokens} 核对）：I2T 为 \textbf{1386 token}（1369 个图像 patch + 17 个文本 token），M* 默认预处理同为 1369 patch。I2I 两边均为 1369 个图像 patch，vLLM-Omni 由源码确认（usage 不记录图像 token）。I2I 指令为“Replace the bagel with a croissant.”。
  \item \kw{前缀缓存}：M* 默认开，vLLM-Omni 的 Thinker 默认关。基准每次发送相同的图和问题，缓存会命中，因此 I2T 分缓存关、开两组比较。
  \item \kw{测量与步数}：每个 case 做 2 次 warmup 与 5 次稳态测量，串行 HTTP，端到端计时。步数与 shift 固定为 50 / 3.0；M* 不暴露随机种子，只比较延迟。
  \item \kw{只做不改变计算量的优化}：并行拓扑、compile、graph、connector、前缀缓存。不使用 TeaCache、Cache-DiT、量化或减少步数。
\end{itemize}
\takeaway{四类 case 两边的计算量一致：图像 token 均为 1369 patch（I2I 另有 1024 个 VAE latent），}
\end{frame}

% BAGEL 两卡实测结果：全部结果矩阵、最优配置对比、等量比较、首请求
\setpart{BAGEL 实测 · 两卡结果}

\begin{frame}[t]{两卡全部实测结果：卡 1、2，稳态中位数，秒}
\scriptsize\def\codesize{\relax}
\renewcommand{\arraystretch}{1.12}\setlength{\tabcolsep}{3pt}
\begin{tabularx}{\textwidth}{>{\bfseries}p{5.9cm} p{0.9cm} Y Y Y Y}
\toprule
\rowcolor{Mist}
配置 & 框架 & T2I 512²·50 步 & I2T 1 token & I2T 64 token & I2I 512²·50 步 \\
\midrule
M* CFG 并行 · compile 开 · 缓存开 · 无 graph & \textcolor{Teal}{M*} & \cellcolor{GoldBg}\textbf{\textcolor{Teal}{7.68}}\,$\star$ & 0.085 & 2.46 & \cellcolor{GoldBg}\textbf{\textcolor{Teal}{16.79}}\,$\star$ \\
M* CFG 并行 · compile 关（诊断）· 缓存开 & \textcolor{Teal}{M*} & 8.48 & \cellcolor{GoldBg}\textbf{\textcolor{Teal}{0.083}}\,$\star$ & \cellcolor{GoldBg}\textbf{\textcolor{Teal}{2.17}}\,$\star$ & 18.25 \\
M* TP2 · compile 开 · 缓存开 · 无 graph & \textcolor{Teal}{M*} & 12.96 & 0.330 & 3.16 & 21.21 \\
M* CFG 并行 · compile 开 · 缓存关 & \textcolor{Teal}{M*} & -- & 0.293 & 2.66 & \cellcolor{GoldBg}\textbf{\textcolor{Teal}{16.82}}\,$\star$ \\
\midrule
Omni 单 stage CFG 并行 · 无 compile & \textcolor{Orange}{Omni} & \cellcolor{GoldBg}\textbf{\textcolor{Teal}{9.98}}\,$\star$ & -- & -- & -- \\
Omni 单 stage TP2 · 无 graph & \textcolor{Orange}{Omni} & 17.04 & -- & -- & \cellcolor{GoldBg}\textbf{\textcolor{Teal}{28.08}}\,$\star$ \\
Omni two-stage · 缓存关 · eager & \textcolor{Orange}{Omni} & 17.99 & 0.193 & 1.88 & \textcolor{Muted}{30.76\textsuperscript{†}} \\
Omni two-stage · 缓存关 · Thinker graph & \textcolor{Orange}{Omni} & -- & 0.191 & 1.86 & -- \\
Omni two-stage · 缓存开 · eager & \textcolor{Orange}{Omni} & -- & \cellcolor{GoldBg}\textbf{\textcolor{Teal}{0.053}}\,$\star$ & \cellcolor{GoldBg}\textbf{\textcolor{Teal}{1.74}}\,$\star$ & -- \\
\bottomrule
\end{tabularx}

\vspace{3pt}
\begin{itemize}\scriptsize\setlength\itemsep{0pt}
  \item 配置列即该行的完整定义，见第 3 部分。$\star$ = 同一框架内该 case 的最佳（1\% 以内并列）。n = 5，min--max 见附录 A。
  \item 图生文两边输入规模一致（M* 1369 patch，Omni 1386 token）；图像编辑两边指令相同。--：未测或不支持；\textsuperscript{†}灰图，只保留耗时。
\end{itemize}
\takeaway{M* 最优：文生图 \val{final/mstar-cfgpar/t2i} s，图像编辑 \val{final/mstar-cfgpar/i2i} s。vLLM-Omni 最优：文生图 \val{final/omni-single-cfgpar2/t2i} s，图生文 \val{fact/omni-two-stage-prefixcache/i2t1} s（1 token）/ \val{fact/omni-two-stage-prefixcache/i2t64} s（64 token，缓存开）。}
\end{frame}

\newcommand{\barpair}[6]{% y, label, mstar value, omni value, mstar tag, omni tag
    \pgfmathsetmacro{\mx}{max(#3,#4)}
    \pgfmathsetmacro{\wm}{#3/\mx*6.6}
    \pgfmathsetmacro{\wo}{#4/\mx*6.6}
    \node[anchor=east,text width=3.6cm,align=right,font=\scriptsize\bfseries] at (3.7,#1-0.34) {#2};
    \fill[Teal] (3.9,#1-0.04) rectangle ++(\wm,-0.28);
    \node[anchor=west,font=\tiny] at (3.95+\wm,#1-0.18) {#3\,s\ \textcolor{Muted}{#5}};
    \fill[Orange] (3.9,#1-0.36) rectangle ++(\wo,-0.28);
    \node[anchor=west,font=\tiny] at (3.95+\wo,#1-0.50) {#4\,s\ \textcolor{Muted}{#6}};}

\begin{frame}[t]{各 case 的最优配置对比（柱长按 case 内较大值归一）}
\begin{tikzpicture}[x=1cm,y=.72cm,font=\scriptsize]
  \fill[Teal] (3.9,.85) rectangle ++(.3,-.25); \node[anchor=west,font=\scriptsize] at (4.3,.72) {M*（CFG 并行，compile 开）};
  \fill[Orange] (8.2,.85) rectangle ++(.3,-.25); \node[anchor=west,font=\scriptsize] at (8.6,.72) {vLLM-Omni};
  \barpair{0}{文生图 T2I}{\raw{final/mstar-cfgpar/t2i}}{\raw{final/omni-single-cfgpar2/t2i}}{CFG 并行}{单 stage CFG 并行}
  \barpair{-1.0}{图像编辑 I2I}{\raw{final/mstar-cfgpar/i2i}}{\raw{final/omni-single-tp2/i2i}}{CFG 并行}{单 stage TP2}
  \barpair{-2.0}{图生文 1 token，缓存关}{\raw{fact/mstar-cfgpar-nocache/i2t1}}{\raw{final/omni-two-stage/i2t1}}{CFG 并行}{two-stage，eager}
  \barpair{-3.0}{图生文 1 token，缓存开}{\raw{final/mstar-cfgpar/i2t1}}{\raw{fact/omni-two-stage-prefixcache/i2t1}}{}{two-stage + 前缀缓存}
  \barpair{-4.0}{图生文 64 token，缓存关}{\raw{fact/mstar-cfgpar-nocache/i2t64}}{\raw{final/omni-two-stage/i2t64}}{CFG 并行}{two-stage，eager}
  \barpair{-5.0}{图生文 64 token，缓存开}{\raw{final/mstar-cfgpar/i2t64}}{\raw{fact/omni-two-stage-prefixcache/i2t64}}{}{two-stage + 前缀缓存}
\end{tikzpicture}
\takeaway{vLLM-Omni 相对 M* 的耗时：文生图为 \pct{bestTti}，图像编辑为 \pct{bestIti}。图生文 1 token 为 \pct{i2tOff1}（缓存关）、\pct{i2tOn1}（缓存开）；64 token 为 \pct{i2tOff64}、\pct{i2tOn64}。}
\end{frame}

\begin{frame}[t]{图生文：输入与缓存设置对齐后的结果}
\scriptsize\def\codesize{\relax}
\renewcommand{\arraystretch}{1.12}\setlength{\tabcolsep}{3pt}
\begin{tabularx}{\textwidth}{>{\bfseries}p{4.6cm} Y Y Y Y}
\toprule
\rowcolor{Mist}
M* 配置 & 1 token，4900 patch（早期） & 1 token，1369 patch（本次） & 64 token，4900 patch（早期） & 64 token，1369 patch（本次） \\
\midrule
CFG 并行（compile 开） & \textcolor{Muted}{0.193} & \textbf{\textcolor{Teal}{0.085}} & \textcolor{Muted}{2.55} & \textbf{\textcolor{Teal}{2.46}} \\
CFG 并行（compile 关） & \textcolor{Muted}{0.200} & \textbf{\textcolor{Teal}{0.083}} & \textcolor{Muted}{2.29} & \textbf{\textcolor{Teal}{2.17}} \\
TP2（compile 开） & \textcolor{Muted}{0.959} & \textbf{\textcolor{Teal}{0.330}} & \textcolor{Muted}{3.76} & \textbf{\textcolor{Teal}{3.16}} \\
CFG 并行，KV 前缀缓存关 & \textcolor{Muted}{0.826} & \textbf{\textcolor{Teal}{0.293}} & \textcolor{Muted}{3.18} & \textbf{\textcolor{Teal}{2.66}} \\
\midrule
\textcolor{Orange}{Omni two-stage（1386 token 提示），缓存关} & -- & \textbf{\textcolor{Orange}{0.193}} & -- & \textbf{\textcolor{Orange}{1.882}} \\
\textcolor{Orange}{Omni two-stage，Thinker 前缀缓存开} & -- & \textbf{\textcolor{Orange}{0.053}} & -- & \textbf{\textcolor{Orange}{1.744}} \\
\bottomrule
\end{tabularx}

\vspace{4pt}
\begin{itemize}\scriptsize\setlength\itemsep{1pt}
  \item vLLM-Omni two-stage 的图生文请求返回 \code{usage.prompt\_tokens = 1386}（1369 个图像 patch + 17 个文本 token）。M* 的 \code{vit\_patches(512, 512, "default")} 同为 1369。
  \item 早期测试的 M* 图生文使用了 4900 patch（\code{image\_preprocess:"vllm"}），与 vLLM-Omni main 的 1386 token 不等量，标 \textsuperscript{§}，仅供参考。
  \item 灰色 = 早期的 4900 patch 数据；粗体 = 输入对齐后的数据（\code{results/20261122\_i2t\_parity}）。1 token 延迟降低 56--66\%，64 token 降低 4--17\%。
\end{itemize}
\takeaway{对齐后，M* 的 1 token 延迟为 0.085 s（缓存开）、0.293 s（缓存关）；vLLM-Omni 为 0.053 s、0.193 s。两种设置下 vLLM-Omni 都更快。}
\end{frame}

\begin{frame}[t]{图像编辑：同一指令、同一图像 token 量}
\scriptsize\def\codesize{\relax}
\renewcommand{\arraystretch}{1.08}\setlength{\tabcolsep}{3pt}
\begin{tabularx}{\textwidth}{>{\bfseries}p{4.6cm} p{2.5cm} Y}
\toprule
\rowcolor{Mist}
配置 & 图像 token（ViT / VAE） & I2I 中位数（min--max，秒） \\
\midrule
\textcolor{Teal}{M* CFG 并行（compile 开）} & 1369 / 1024 & \textbf{16.79\,\textcolor{Muted}{(16.75--16.81)}} \\
\textcolor{Teal}{M* CFG 并行（compile 关）} & 1369 / 1024 & 18.25\,\textcolor{Muted}{(18.23--18.26)} \\
\textcolor{Teal}{M* TP2（compile 开）} & 1369 / 1024 & 21.21\,\textcolor{Muted}{(21.20--21.21)} \\
\textcolor{Teal}{M* CFG 并行，KV 前缀缓存关} & 1369 / 1024 & 16.82\,\textcolor{Muted}{(16.80--16.83)} \\
\textcolor{Orange}{vLLM-Omni 单 stage TP2（512$^2$ 原生输入）} & 1369 / 1024 & \textbf{28.08\,\textcolor{Muted}{(28.07--28.08)}} \\
\midrule
\multicolumn{3}{l}{\textcolor{Muted}{早期指令（“Add a small cat sitting beside the bagel on the plate.”）的运行}} \\
\textcolor{Muted}{M* CFG 并行，\code{vllm}（4900 patch）} & 4900 / 1024 & 18.98\,\textcolor{Muted}{(18.96--19.00)} \\
\textcolor{Muted}{M* CFG 并行，\code{default}（输出未加入猫）} & 1369 / 1024 & 16.83\,\textcolor{Muted}{(16.81--16.83)} \\
\textcolor{Muted}{vLLM-Omni 单 stage TP2（512$^2$ 原生）} & 1369 / 1024 & 28.07\,\textcolor{Muted}{(28.06--28.07)} \\
\textcolor{Muted}{vLLM-Omni 单 stage TP2（输入放大到 980$^2$）} & 4900 / 3721 & 31.24\,\textcolor{Muted}{(31.24--31.25)} \\
\bottomrule
\end{tabularx}

\vspace{2pt}
\begin{itemize}\scriptsize\setlength\itemsep{0pt}
  \item 上半部分：两边同一指令 \code{Replace the bagel with a croissant.}，ViT 1369 patch，VAE 1024 latent，输出均为可颂。
  \item 下半部分：早期“加猫”指令，仅作记录。M* 默认预处理下未生效，故改用上述指令。
  \item Omni 放大到 980$^2$：VAE 为 3721 latent，不等量，不用于比较。
\end{itemize}
\takeaway{等量后 M* CFG 并行仍快约 40\%：\val{final/mstar-cfgpar/i2i} s，vLLM-Omni 为 \val{final/omni-single-tp2/i2i} s。}
\end{frame}

\begin{frame}[t]{首个请求（含 JIT、编译、图捕获）的延迟，秒}
\scriptsize\def\codesize{\relax}
\renewcommand{\arraystretch}{1.12}\setlength{\tabcolsep}{3pt}
\begin{tabularx}{\textwidth}{>{\bfseries}p{5.2cm} Y Y Y Y}
\toprule
\rowcolor{Mist}
配置 & T2I & I2T 1 token & I2T 64 token & I2I \\
\midrule
M* CFG 并行 · compile 开 · 缓存开 · 无 graph & 17.5 & 8.4 & 3.5 & 34.9 \\
M* CFG 并行 · compile 关（诊断）· 缓存开 & 10.9 & 1.6 & 2.2 & 22.0 \\
M* TP2 · compile 开 · 缓存开 · 无 graph & 24.6 & 16.7 & 4.3 & 39.7 \\
M* CFG 并行 · compile 开 · 缓存关 & -- & 6.7 & 3.6 & 32.3 \\
Omni 单 stage CFG 并行 · 无 compile & 10.9 & -- & -- & -- \\
Omni 单 stage TP2 · 无 graph & 17.7 & -- & -- & 29.5 \\
Omni two-stage · 缓存关 · eager & 18.7 & 0.6 & 1.9 & -- \\
Omni two-stage · 缓存关 · Thinker graph & -- & 0.4 & 1.9 & -- \\
Omni two-stage · 缓存开 · eager & -- & 0.6 & 1.7 & -- \\
\bottomrule
\end{tabularx}

\vspace{4pt}
\begin{itemize}\scriptsize\setlength\itemsep{1pt}
  \item 首请求指服务启动后，该 case 的第一次请求（warmup 1）。第二次 warmup 与稳态已无明显差异。M* 的图生文行来自单独的等量运行；vLLM-Omni 的两行同理。
  \item M* 默认开启 \code{torch.compile}：T2I、I2I、I2T 的首请求均多出约 6--7 s。关闭后首请求更短，但扩散稳态慢 7--9\%。
  \item vLLM-Omni 无 compile，首请求只多出内核懒加载的时间。单 stage CFG 并行的首个 T2I 比稳态多约 1 s。
\end{itemize}
\takeaway{需要冷启动即时响应的场景中，M* 的 compile 会使第一个请求多等约 6--7 s。稳态服务中该开销可以忽略。}
\end{frame}

% 各因素对性能的影响
\setpart{BAGEL 实测 · 各因素的影响}

\begin{frame}[t]{各因素对延迟的影响一览（A $\rightarrow$ B 的变化，负值为更快）}
\scriptsize\def\codesize{\relax}
\renewcommand{\arraystretch}{1.12}\setlength{\tabcolsep}{3pt}
\begin{tabularx}{\textwidth}{>{\bfseries}p{3.1cm} p{3.6cm} r r r r Y}
\toprule
\rowcolor{Mist}
因素 & A $\rightarrow$ B & T2I & I2I & I2T 1 & I2T 64 & 备注 \\
\midrule
\textcolor{Teal}{M*} 并行拓扑 & TP2 $\rightarrow$ CFG 并行 & \pct{topoMstarTti} & \pct{topoMstarIti} & \pct{topoMstarPre} & \pct{topoMstarDec} & 影响最大 \\
\textcolor{Orange}{Omni} 并行拓扑 & 单 stage TP2 $\rightarrow$ CFG 并行 & \pct{topoOmniTti} & 不支持 & -- & -- & I2I 需要 3 卡 \\
\textcolor{Teal}{M*} torch.compile & 关 $\rightarrow$ 开（CFG 并行） & \pct{cmpCfgTti} & \pct{cmpCfgIti} & \pct{cmpCfgPre} & \pct{cmpCfgDec} & 解码变慢 \\
\textcolor{Orange}{Omni} two-stage connector & SHM $\rightarrow$ Mooncake-TCP & \pct{connMoonTti} & \pct{connMoonIti} & -- & -- & 见下页 \\
\textcolor{Orange}{Omni} two-stage connector & SHM $\rightarrow$ NIXL & \pct{connNixlTti} & \pct{connNixlIti} & -- & -- & NIXL 无法使用 \\
\textcolor{Orange}{Omni} Thinker 前缀缓存 & 关 $\rightarrow$ 开 & -- & -- & \pct{cacheOmni1} & \pct{cacheOmni64} & 同图重复提问时命中 \\
\textcolor{Teal}{M*} KV 前缀缓存 & 关 $\rightarrow$ 开（默认开） & -- & \pct{cacheMstarIti} & \pct{cacheMstar1} & \pct{cacheMstar64} & 同上 \\
\textcolor{Orange}{Omni} 软件版本 & v0.32 $\rightarrow$ main（单 stage TP2） & \pct{verOmniTti} & \pct{verOmniIti} & -- & -- & 两阶段解码 \pct{verOmni2Dec} \\
SDPA 修复（OpenCL） & MATH $\rightarrow$ 默认 oneDNN & $\approx$0 & -- & $-$33\%\textsuperscript{§} & -- & 仅 M* 受影响 \\
\bottomrule
\end{tabularx}
\vspace{2pt}
\memo{百分比 = B 相对 A 的稳态中位数变化。“--” 表示该组合不适用或无有效数据。\textsuperscript{§}两侧输入相同，仅作 A/B。}
\takeaway{并行拓扑的影响（$-$17\% 到 $-$74\%）远大于 compile 与缓存（个位数到约 10\%）。两框架的差距主要来自拓扑选择。}
\end{frame}

\begin{frame}[t]{并行拓扑：CFG 并行与 TP2 有什么区别}
\scriptsize\def\codesize{\relax}
\renewcommand{\arraystretch}{1.15}\setlength{\tabcolsep}{3pt}
\begin{tabularx}{\textwidth}{>{\bfseries}p{2.4cm} Y Y}
\toprule
\rowcolor{Mist}
 & TP2（张量并行） & CFG 并行 \\
\midrule
切分方式 & 每一层的权重矩阵切成两半，两张卡各算一半 & 权重不切分。两路 CFG 分支各占一张卡 \\
每张卡的权重 & 约一半（约 13--14 GiB） & 完整一份（约 27 GiB） \\
每张卡的计算 & 每路分支的一半 & 一整路分支 \\
卡间同步 & 每一层算完都要 all-reduce，两张卡交换部分结果并相加 & 每个去噪步交换一次引导结果 \\
文生图的同步次数 & 约 5000 次（50 步，每步约 100 次） & 约 50 次（每步 1 次） \\
图像编辑（3 路分支） & 两张卡各算一半的全部分支，同步次数更多 & 主分支在卡 0，两路引导分支串行在卡 1，卡 1 是关键路径 \\
\bottomrule
\end{tabularx}
\vspace{4pt}
\begin{itemize}\scriptsize\setlength\itemsep{2pt}
  \item \textbf{一个去噪步的过程}。TP2：两张卡同时算第 1 层的一半，all-reduce；算第 2 层的一半，all-reduce；……，每步约 100 次。CFG 并行：卡 1 完整算条件分支，同时卡 2 完整算无条件分支；交换一次，合并引导，进入下一步。
  \item 两种方式在理论上每张卡的计算量相同，区别在于 \textbf{卡间同步的次数与粒度}。TP2 的同步在每层之间，粒度小、次数多；CFG 并行的同步只在每步之间。
  \item 代价：CFG 并行每张卡都要放一份完整模型，显存占用更高（每卡约 31 GiB，几乎占满）。
\end{itemize}
\takeaway{TP2 的代价是每步约 100 次跨卡同步；CFG 并行每步只有 1 次。在本机的 PCIe 上，一次同步约 1.3 ms，所以 TP2 每张图多花约 6.5 s。}
\end{frame}

\begin{frame}[t]{并行拓扑：数值与通信开销}
\scriptsize\def\codesize{\relax}
\begin{minipage}[t]{.52\textwidth}
\begin{itemize}\scriptsize\setlength\itemsep{2pt}
  \item \textbf{通信量}：一次文生图（50 步，两路分支）约 \textbf{5000 次} all-reduce（实测计数）。
  \item \textbf{单次开销}：卡 (1,2) 上一次 7 MB 的 all-reduce 约 \textbf{1.3 ms}（约 5.6 GB/s，经 PCIe）。5000 $\times$ 1.3 ms $\approx$ 6.5 s。
  \item \textbf{与实测对照}：TP2 比 CFG 并行慢 5.3 s（M*）、7.1 s（vLLM-Omni），与通信估算基本吻合。两框架各自的计算量相同，差距基本来自通信。
  \item \textbf{I2T}：CFG 并行下也快很多（M* 1 token 从 \raw{final/mstar-tp2/i2t1} s 降到 \raw{final/mstar-cfgpar/i2t1} s）。TP2 的预填充与解码每层都有 all-reduce。
\end{itemize}
\end{minipage}\hfill
\begin{minipage}[t]{.45\textwidth}
\renewcommand{\arraystretch}{1.15}\setlength{\tabcolsep}{3pt}
\begin{tabularx}{\linewidth}{>{\bfseries}p{2.4cm} Y Y}
\toprule
\rowcolor{Mist}
秒 & TP2 & CFG 并行 \\
\midrule
\textcolor{Teal}{M*} T2I & \val{final/mstar-tp2/t2i} & \textbf{\val{final/mstar-cfgpar/t2i}} \\
\textcolor{Teal}{M*} I2I & \val{final/mstar-tp2/i2i} & \textbf{\val{final/mstar-cfgpar/i2i}} \\
\textcolor{Teal}{M*} I2T 1 token & \val{final/mstar-tp2/i2t1} & \textbf{\val{final/mstar-cfgpar/i2t1}} \\
\textcolor{Teal}{M*} I2T 64 token & \val{final/mstar-tp2/i2t64} & \textbf{\val{final/mstar-cfgpar/i2t64}} \\
\textcolor{Orange}{Omni} T2I & \val{final/omni-single-tp2/t2i} & \textbf{\val{final/omni-single-cfgpar2/t2i}} \\
\bottomrule
\end{tabularx}
\end{minipage}
\vspace{4pt}
\begin{itemize}\scriptsize\setlength\itemsep{1pt}
  \item 通信速度取决于卡对。卡 (0,1) 的数据未测，因此在通信更快的卡对上，两种拓扑的相对优劣可能不同，未验证。
\end{itemize}
\takeaway{在这对卡上，PCIe all-reduce 是 TP2 的主要开销。使用不需要层内通信的 CFG 并行后，两个框架都快 40\% 以上。}
\end{frame}

\begin{frame}[t]{torch.compile：在 Intel GPU 上收益有限}
\scriptsize\def\codesize{\relax}
\renewcommand{\arraystretch}{1.15}\setlength{\tabcolsep}{3pt}
\begin{tabularx}{\textwidth}{>{\bfseries}p{3.4cm} r r r r r r r r}
\toprule
\rowcolor{Mist}
M* 配置 & T2I 关 & T2I 开 & I2I 关 & I2I 开 & I2T 64 关 & I2T 64 开 & 首请求 T2I 关 & 首请求 T2I 开 \\
\midrule
CFG 并行（秒） & \val{final/mstar-cfgpar-nocompile/t2i} & \textbf{\val{final/mstar-cfgpar/t2i}} & \val{final/mstar-cfgpar-nocompile/i2i} & \textbf{\val{final/mstar-cfgpar/i2i}} & \textbf{\val{final/mstar-cfgpar-nocompile/i2t64}} & \val{final/mstar-cfgpar/i2t64} & \warm{final/mstar-cfgpar-nocompile/t2i} & \warm{final/mstar-cfgpar/t2i} \\
TP2（秒，早期批次） & \val{v032/mstar-tp2-nocompile/t2i} & \textbf{\val{v032/mstar-tp2/t2i}} & \val{v032/mstar-tp2-nocompile/i2i} & \textbf{\val{v032/mstar-tp2/i2i}} & \textbf{\anyval{v032/mstar-tp2-nocompile/i2t64}}\textsuperscript{§} & \anyval{v032/mstar-tp2/i2t64}\textsuperscript{§} & \warm{v032/mstar-tp2-nocompile/t2i} & \warm{v032/mstar-tp2/t2i} \\
\bottomrule
\end{tabularx}
\vspace{4pt}
\begin{itemize}\scriptsize\setlength\itemsep{2pt}
  \item 同一拓扑下，compile 开与关的对比。\textsuperscript{§}早期批次的图生文为缓存默认开的设置，仅作 A/B。compile 在 XPU 上主要做 Inductor / Triton-XPU 的算子融合：扩散加速 \textbf{7--9\%}（CFG 并行：T2I \pct{cmpCfgTti}，I2I \pct{cmpCfgIti}）；小 batch 的解码反而\textbf{变慢}（\pct{cmpCfgDec}，TP2 下更明显）。
  \item 首请求多出约 6--7 s 的编译时间。M* 没有官方开关，“关”是通过 hook 让 \code{torch.compile} 失效实现的，仅作诊断。
  \item CUDA 上 compile 的主要收益来自 CUDA graph，减少了内核启动开销。B70 上 GPU 利用率采样的最高值只有 14--22\%，推测 eager 路径受内核启动限制。XPU graph 在两卡上无法用于 BAGEL；三卡结果见第 6 部分。
\end{itemize}
\takeaway{Intel GPU 上 compile 只带来个位数百分比的扩散加速，并使解码变慢约 10\%。CUDA 上的大幅提速主要来自 CUDA graph；B70 上两卡无法获得；三卡结果见第 6 部分。}
\end{frame}


\begin{frame}[t]{两阶段 connector：SharedMemory、Mooncake、NIXL}
\scriptsize\def\codesize{\relax}
\renewcommand{\arraystretch}{1.12}\setlength{\tabcolsep}{3pt}
\begin{tabularx}{\textwidth}{>{\bfseries}p{5.9cm} p{0.9cm} Y Y}
\toprule
\rowcolor{Mist}
配置 & 框架 & T2I 512²·50 步 & I2I 512²·50 步 \\
\midrule
SharedMemory（默认） & \textcolor{Orange}{Omni} & \cellcolor{GoldBg}\textbf{\textcolor{Teal}{17.962\,\textcolor{Muted}{(17.823--17.979)}}}\,$\star$ & \textcolor{Muted}{30.74\textsuperscript{†}} \\
Mooncake TransferEngine（TCP） & \textcolor{Orange}{Omni} & \cellcolor{GoldBg}\textbf{\textcolor{Teal}{17.983\,\textcolor{Muted}{(17.976--17.994)}}}\,$\star$ & \textcolor{Muted}{29.92\textsuperscript{†}} \\
NIXL（UCX） & \textcolor{Orange}{Omni} & \textcolor{Red}{失败\textsuperscript{‡}} & \textcolor{Red}{失败\textsuperscript{‡}} \\
\bottomrule
\end{tabularx}

\vspace{3pt}
\begin{itemize}\scriptsize\setlength\itemsep{1.5pt}
  \item 同一 two-stage 拓扑（Thinker 在卡 1，DiT 在卡 2，eager），仅改变两阶段之间 KV 的传输方式。n = 5，括号为 min--max。
  \item \textbf{SharedMemory 与 Mooncake（TCP）}：文生图相同（\pct{connMoonTti}）。KV 很小，传输不是瓶颈。图像编辑中 Mooncake 快约 0.8 s（\pct{connMoonIti}），KV 约 6 k token。
  \item \textbf{NIXL（UCX）无法在 B70 上使用}\textsuperscript{‡}。UCX 把 XPU 显存识别为主机内存，显存注册失败。KV 传输重试 3 次后失败，DiT 等待 30 s 超时后继续。因此 T2I 的 48.3 s = 18 s + 30 s，\textbf{这不是传输耗时}。\code{nixl\_connector.py} 按张量所在设备选择内存类型，没有 CPU 暂存选项。
  \item \textsuperscript{†}two-stage 的 I2I 输出为灰图（main 版本），耗时仅供参考。NIXL 失败时（DiT 未收到 KV），I2I 输出正常，说明灰图与 KV 传输路径有关。原因未查明。
\end{itemize}
\memo{对照：NVIDIA RTX 5880 上，Mooncake-TCP 的 I2I（50 步）节省 2.25 s（\textminus 8\%）。B70 上同一设置只节省约 0.8 s（\textminus 3\%），原因未单独测量。}
\takeaway{在 B70 上，connector 的选择对 I2I 的影响约为 3\%。NIXL 目前无法在 XPU 上使用。Mooncake-TCP 可用，但收益有限。}
\end{frame}

\begin{frame}[t]{前缀缓存：相同输入命中缓存后，I2T 预填充明显缩短}
\scriptsize\def\codesize{\relax}
\renewcommand{\arraystretch}{1.12}\setlength{\tabcolsep}{3pt}
\begin{tabularx}{\textwidth}{>{\bfseries}p{3.6cm} p{3.2cm} r r r r}
\toprule
\rowcolor{Mist}
框架 & 设置 & I2T 1 token & I2T 64 token & 解码每 token（ms） & 说明 \\
\midrule
\textcolor{Teal}{M*}（CFG 并行） & KV 前缀缓存关 & \raw{fact/mstar-cfgpar-nocache/i2t1} & \raw{fact/mstar-cfgpar-nocache/i2t64} & \pct{msTokMoff} & 无缓存命中 \\
\textcolor{Teal}{M*}（CFG 并行） & KV 前缀缓存开（默认） & \raw{final/mstar-cfgpar/i2t1} & \raw{final/mstar-cfgpar/i2t64} & \pct{msTokM} & 缓存命中 \\
\textcolor{Orange}{Omni} two-stage & Thinker 缓存关（默认，eager） & \raw{final/omni-two-stage/i2t1} & \raw{final/omni-two-stage/i2t64} & \pct{msTokOg} & 无缓存命中 \\
\textcolor{Orange}{Omni} two-stage & Thinker 缓存开 & \raw{fact/omni-two-stage-prefixcache/i2t1} & \raw{fact/omni-two-stage-prefixcache/i2t64} & \pct{msTokO} & 缓存命中 \\
\bottomrule
\end{tabularx}
\vspace{4pt}
\begin{itemize}\scriptsize\setlength\itemsep{1.5pt}
  \item 缓存对 1 token（即预填充）影响最大：M* \pct{cacheMstar1}，vLLM-Omni \pct{cacheOmni1}。对 64 token 只省去预填充部分（M* \pct{cacheMstar64}，vLLM-Omni \pct{cacheOmni64}）。
  \item \textbf{每个 token 的解码}：vLLM-Omni 约 26.8 ms，M* 约 37.7 ms（compile 关时为 33.1 ms）。vLLM-Omni 的解码快约 29\%，与缓存设置无关。
  \item 缓存只影响图生文。T2I 没有共享前缀；M* 的 I2I 缓存关、开分别为 \raw{fact/mstar-cfgpar-nocache/i2i} s、\raw{final/mstar-cfgpar/i2i} s，无明显差别。
  \item 实际服务中，同一张图被连续提问时才会命中缓存。单次请求的延迟应以“缓存关”为准。两个框架的默认设置不同，只比较默认值会把设置差异误认为框架差距。
\end{itemize}
\takeaway{缓存关对缓存关、开对开时，vLLM-Omni 的图生文都更快（1 token 快 34--38\%，64 token 快 29\%）。缓存可以把 1 token 延迟降到约三分之一，但只在重复输入时有效。}
\end{frame}

% XPU graph 研究（三卡）：图生文解码、扩散步、TP2 图捕获、去噪步 kernel 分布
\setpart{BAGEL 实测 · XPU graph 与 kernel 分布（三卡）}

\begin{frame}[t]{XPU graph（三卡，卡 0、1、2）：图生文解码，中位数，秒}
\scriptsize\def\codesize{\relax}
\renewcommand{\arraystretch}{1.12}\setlength{\tabcolsep}{3pt}
% generated by scripts/make_tex_tables.py -- do not edit
\scriptsize
\begin{tabularx}{\textwidth}{>{\bfseries}p{4.6cm} c r r r r r}
\toprule
\rowcolor{Mist}
配置 & 缓存 & I2T 1 token eager & graph & I2T 64 token eager & graph & 64 token 变化 \\
\midrule
M* CFG 并行（三卡） & 开 & 0.089 & 0.085 & 2.42 & 1.95 & $-$20\% \\
M* CFG 并行（三卡） & 关 & 0.304 & 0.300 & 2.68 & 2.15 & $-$20\% \\
M* TP2（三卡，LLM 卡 1、2） & 开 & 0.258 & 0.226 & 3.05 & 1.52 & $-$50\% \\
vLLM-Omni two-stage（三卡，Thinker TP2 卡 0、1） & 关 & 0.127 & 0.125 & 1.36 & 1.10 & $-$19\% \\
vLLM-Omni two-stage（三卡，Thinker TP2 卡 0、1） & 开 & 0.054 & 0.051 & 1.25 & 1.02 & $-$19\% \\
\bottomrule
\end{tabularx}

\vspace{4pt}
\begin{itemize}\scriptsize\setlength\itemsep{1pt}
  \item 配置：M* CFG 并行，三个 rank 分别为卡 0 的主分支、卡 1 与卡 2 的两路引导分支；M* TP2，卡 0 为 ViT 与 VAE，卡 1、2 切分 LLM；Omni two-stage，Thinker TP2 在卡 0、1，DiT 在卡 2（图生文不使用 DiT）。
  \item graph 只作用于 LLM 解码步。M* CFG 并行的 graph 无需补丁；M* TP2 的 graph 需要补丁（见下页）。
  \item 两边图像输入规模一致：M* 为默认预处理的 1369 个图像 patch，Omni 提示为 1386 token（含 17 个文本 token）。
\end{itemize}
\takeaway{三卡下 graph 使 M* 的 64 token 解码在 CFG 并行中快 \pct{gCfgOn64}，在 TP2 中快 \pct{gTp2G64}。M* TP2 graph 后为 \val{tp2g/mstar-tp2-graph-fixed/i2t64} s，Omni two-stage graph 为 \val{g3/omni-tp2-graph-chat/i2t64} s，M* 仍慢 \pct{gTp2Gap}。}
\end{frame}

\begin{frame}[t]{XPU graph（三卡）：扩散步不变，TP2 图捕获的限制}
\scriptsize\def\codesize{\relax}
\renewcommand{\arraystretch}{1.12}\setlength{\tabcolsep}{3pt}
% generated by scripts/make_tex_tables.py -- do not edit
\scriptsize
\begin{tabularx}{\textwidth}{>{\bfseries}p{5.6cm} r r r r r}
\toprule
\rowcolor{Mist}
配置 & T2I eager & T2I graph & I2I eager & I2I graph & I2I 变化 \\
\midrule
M* CFG 并行（三卡） & 7.73 & 7.75 & 9.28 & 9.26 & $-$0\% \\
M* TP2（三卡，LLM 卡 1、2） & 7.68 & 7.64 & -- & -- & n/a \\
vLLM-Omni 单 stage CFG 并行（三卡） & -- & -- & 11.19 & -- & n/a \\
vLLM-Omni two-stage（三卡，Thinker TP2 卡 0、1，DiT 卡 2） & 17.91 & 18.01 & -- & -- & n/a \\
\bottomrule
\end{tabularx}

\vspace{4pt}
\begin{itemize}\scriptsize\setlength\itemsep{1pt}
  \item 扩散步两边都是 eager，graph 不作用于扩散。M* CFG 并行的 T2I 与 I2I 在 graph 前后基本持平（T2I \pct{gDiffT}，I2I \pct{gDiffI}）。vLLM-Omni 的 DiT 不支持 denoise graph。
  \item M* TP2 图捕获失败的原因有三：(1) XCCL 的 \code{all\_reduce} 消息超过约 4 MiB 时无法进入 graph；(2) inductor 生成的 functional \code{all\_gather}（lm\_head 汇聚 logits）在 XPU graph 中不可用，普通 \code{all\_gather} 在约 4 行 logits 以内可用；(3) 最大的 prefill 捕获失败后，捕获状态残留，decode 随之失败。
  \item 补丁（\code{scripts/patches/mstar\_xpu\_tp2\_graph.patch}）作用于已安装的 M* 包，未合入上游：\code{all\_gather} 不参与 compile；prefill 捕获 token 上限 512；decode 与 prefill 的捕获 batch size 限为 [1, 2, 4]。
  \item 两卡 CFG 并行的 graph 因显存不足无法启动（每卡约 27 GiB 的完整模型）。
\end{itemize}
\takeaway{graph 只加速图生文的解码：64 token 快 19--50\%；1 token（预填充为主）只快 1--12\%。文生图与图像编辑的扩散步不经过 graph，前后差异在 1\% 以内。M* TP2 的 graph 需要补丁，CFG 并行则不需要。}
\end{frame}

\begin{frame}[t]{去噪步的 kernel 分布（单步，rank 0）}
\scriptsize\def\codesize{\relax}
\renewcommand{\arraystretch}{1.08}\setlength{\tabcolsep}{3pt}
% generated by scripts/make_kernel_table.py -- do not edit
\scriptsize\setlength{\tabcolsep}{2.5pt}
\begin{tabularx}{\textwidth}{>{\bfseries}p{5.0cm} r r r r r}
\toprule
\rowcolor{Mist}
kernel 类别 & Omni ms & Omni 占比 & Omni 次数 & M* 占比 & M* 次数 \\
\midrule
gemm（矩阵乘） & 119.7 & 61.8\% & 229 & 84.1\% & 78 \\
attention（FMHA） & 5.9 & 3.1\% & 28 & 4.5\% & 18 \\
融合 kernel（Triton / vLLM） & 8.8 & 4.5\% & 56 & 10.5\% & 351 \\
index、gather、scatter & 27.1 & 14.0\% & 1017 & 0.8\% & 18 \\
copy、cat & 15.4 & 8.0\% & 1021 & 0.0\% & 0 \\
RMSNorm 拆散部分（pow、mean、reduce） & 5.7 & 2.9\% & 452 & 0.0\% & 0 \\
其他 elementwise & 11.2 & 5.8\% & 1396 & 0.1\% & 18 \\
\midrule
合计（kernel 忙碌时间） & 193.7 & 100\% & 4199 & 100\% & 483 \\
窗口时长 & 201.7 & & & 57.6 & \\
空闲（窗口 − 忙碌） & 8.0 & 4.0\% & & 11.2\% & \\
gemm 平均耗时（ms） & 0.52 & & & 0.55 & \\
\bottomrule
\end{tabularx}

\vspace{4pt}
\begin{itemize}\scriptsize\setlength\itemsep{1pt}
  \item Omni：单 stage CFG 并行，卡 0、1（完整一步）。M*：CFG 并行，卡 0、1、2（58 ms 片段，只比较占比）。
\end{itemize}
\takeaway{Omni 的非 gemm、非 attention 小 kernel 约 68 ms（占 35\%）；单个 gemm 两边平均耗时相同（0.52 与 0.55 ms）。M* 整步未取得，整步差距未完全分配。}
\end{frame}

% 附录：两卡最终矩阵的范围、排除的运行、早期测试、复现
\setpart{附录}

\begin{frame}[t]{附录 A：两卡最终矩阵的稳态中位数与 min--max（秒，n = 5）}
\tiny\def\codesize{\relax}
\renewcommand{\arraystretch}{1.1}\setlength{\tabcolsep}{2.5pt}
\begin{tabularx}{\textwidth}{>{\bfseries}p{5.9cm} p{0.9cm} Y Y Y Y}
\toprule
\rowcolor{Mist}
配置 & 框架 & T2I 512²·50 步 & I2T 1 token & I2T 64 token & I2I 512²·50 步 \\
\midrule
M* CFG 并行 · compile 开 · 缓存开 · 无 graph & \textcolor{Teal}{M*} & \cellcolor{GoldBg}\textbf{\textcolor{Teal}{7.678\,\textcolor{Muted}{(7.655--7.706)}}}\,$\star$ & 0.085\,\textcolor{Muted}{(0.085--0.086)} & 2.458\,\textcolor{Muted}{(2.454--2.461)} & \cellcolor{GoldBg}\textbf{\textcolor{Teal}{16.792\,\textcolor{Muted}{(16.755--16.811)}}}\,$\star$ \\
M* CFG 并行 · compile 关（诊断）· 缓存开 & \textcolor{Teal}{M*} & 8.476\,\textcolor{Muted}{(8.463--8.486)} & \cellcolor{GoldBg}\textbf{\textcolor{Teal}{0.083\,\textcolor{Muted}{(0.082--0.092)}}}\,$\star$ & \cellcolor{GoldBg}\textbf{\textcolor{Teal}{2.169\,\textcolor{Muted}{(2.168--2.172)}}}\,$\star$ & 18.251\,\textcolor{Muted}{(18.235--18.258)} \\
M* TP2 · compile 开 · 缓存开 · 无 graph & \textcolor{Teal}{M*} & 12.955\,\textcolor{Muted}{(12.951--12.961)} & 0.330\,\textcolor{Muted}{(0.329--0.333)} & 3.156\,\textcolor{Muted}{(3.151--3.165)} & 21.208\,\textcolor{Muted}{(21.198--21.212)} \\
M* CFG 并行 · compile 开 · 缓存关 & \textcolor{Teal}{M*} & -- & 0.293\,\textcolor{Muted}{(0.292--0.293)} & 2.656\,\textcolor{Muted}{(2.647--2.658)} & \cellcolor{GoldBg}\textbf{\textcolor{Teal}{16.815\,\textcolor{Muted}{(16.801--16.826)}}}\,$\star$ \\
\midrule
Omni 单 stage CFG 并行 · 无 compile & \textcolor{Orange}{Omni} & \cellcolor{GoldBg}\textbf{\textcolor{Teal}{9.978\,\textcolor{Muted}{(9.973--9.992)}}}\,$\star$ & -- & -- & -- \\
Omni 单 stage TP2 · 无 graph & \textcolor{Orange}{Omni} & 17.040\,\textcolor{Muted}{(17.038--17.043)} & -- & -- & \cellcolor{GoldBg}\textbf{\textcolor{Teal}{28.076\,\textcolor{Muted}{(28.074--28.082)}}}\,$\star$ \\
Omni two-stage · 缓存关 · eager & \textcolor{Orange}{Omni} & 17.987\,\textcolor{Muted}{(17.971--17.992)} & 0.193\,\textcolor{Muted}{(0.192--0.193)} & 1.882\,\textcolor{Muted}{(1.881--1.891)} & \textcolor{Muted}{30.76\textsuperscript{†}} \\
Omni two-stage · 缓存关 · Thinker graph & \textcolor{Orange}{Omni} & -- & 0.191\,\textcolor{Muted}{(0.191--0.192)} & 1.861\,\textcolor{Muted}{(1.856--1.862)} & -- \\
Omni two-stage · 缓存开 · eager & \textcolor{Orange}{Omni} & -- & \cellcolor{GoldBg}\textbf{\textcolor{Teal}{0.053\,\textcolor{Muted}{(0.052--0.053)}}}\,$\star$ & \cellcolor{GoldBg}\textbf{\textcolor{Teal}{1.744\,\textcolor{Muted}{(1.743--1.749)}}}\,$\star$ & -- \\
\bottomrule
\end{tabularx}

\vspace{3pt}
\memo{单元格格式：中位数（min--max）。同一配置内的重复测量波动小于 1\%（M* CFG 并行文生图 7.655--7.706 s；vLLM-Omni 单 stage CFG 并行文生图 9.973--9.992 s）。框架间的差距远大于测量波动。原始数据：\code{results/20261122\_best\_cfg/final/<label>/timings.csv}。}
\takeaway{所有框架间的差距都远大于重复测量的波动（波动小于 1\%）。}
\end{frame}

\begin{frame}[t]{附录 B：排除的运行与原因（一）}
\scriptsize\def\codesize{\relax}
\renewcommand{\arraystretch}{1.12}
\setlength{\tabcolsep}{3pt}
\begin{tabularx}{\textwidth}{>{\bfseries}p{3.6cm} Y}
\toprule
\rowcolor{Mist}
运行 & 现象与处理 \\
\midrule
vLLM-Omni two-stage，Thinker 显存利用率 0.90 / 0.95 & 单卡 KV 空间为零或为负，无法启动。利用率 0.98 且 \code{max\_model\_len} 为 6144 时，KV 约 0.33 GiB，可以运行 \\
vLLM-Omni two-stage，\code{max\_num\_batched\_tokens} = 6144 & 曾怀疑分块导致 I2I 灰图。恢复为 8192 后仍为灰图，排除该原因 \\
vLLM-Omni two-stage，手写 \code{<|image\_pad|>} 提示词 & Thinker 将图像描述为披萨，描述与输入不符。改用服务端 chat 模板后描述正确。该批数据不纳入 \\
vLLM-Omni two-stage I2I（main） & 输出纯灰图（像素极差约 39）。SharedMemory、Mooncake 均为灰图；NIXL 传输失败时输出正常。v0.32.0 输出正常，判断为 main 版本的回归。耗时仅供参考 \\
vLLM-Omni two-stage + NIXL（UCX） & UCX 无法注册 XPU 显存，KV 传输重试 3 次后失败，DiT 等待 30 s 超时后继续。T2I 的 48.3 s、I2I 的 121.6 s 为失败路径耗时，不代表传输速度 \\
\bottomrule
\end{tabularx}
\takeaway{输出内容不正确（灰图、描述错误）的运行，其耗时不进入性能对比。耗时只在明确标注时作参考。}
\end{frame}

\begin{frame}[t]{附录 B：排除的运行与原因（二）}
\scriptsize\def\codesize{\relax}
\renewcommand{\arraystretch}{1.12}
\setlength{\tabcolsep}{3pt}
\begin{tabularx}{\textwidth}{>{\bfseries}p{3.6cm} Y}
\toprule
\rowcolor{Mist}
运行 & 现象与处理 \\
\midrule
M* 图生文使用 \code{image\_preprocess:"vllm"}（早期批次） & 4900 patch，vLLM-Omni main 为 1386 token，输入不等量。已用默认预处理重测；早期数据以 \textsuperscript{§} 标注 \\
vLLM-Omni 单 stage 图生文 & 额外进行了 VAE 编码，描述与输入图不符，且不支持 \code{ignore\_eos}。不作为图生文结果 \\
M* 图像编辑默认预处理，指令“Add a small cat ...”（早期） & 编辑未生效，输出与输入几乎相同（未加入猫）；改用可颂指令后生效。早期运行仅作记录 \\
vLLM-Omni 图像编辑输入放大到 980$^2$ & ViT 4900 patch 与 M* 一致，但 VAE 变为 3721 latent（976$^2$），与 M* 的 1024 不等，不用于比较 \\
两卡 M* TP2 + XPU graph；两卡 M* CFG 并行 + XPU graph & 前者捕获失败并回退 eager；后者显存不足，无法启动 \\
vLLM 0.34 安装 & \code{vllm-xpu-kernels 0.1.17} 不在默认索引中，需要额外的 index（\code{install\_omni\_env.sh} 已处理） \\
\bottomrule
\end{tabularx}
\takeaway{输出内容不正确（灰图、描述错误）的运行，其耗时不进入性能对比。耗时只在明确标注时作参考。}
\end{frame}

\begin{frame}[t]{附录 C：早期测试与本次的测量差异}
\scriptsize\def\codesize{\relax}
\renewcommand{\arraystretch}{1.1}
\setlength{\tabcolsep}{3pt}
\begin{tabularx}{\textwidth}{>{\bfseries}p{3.0cm} p{4.2cm} Y}
\toprule
\rowcolor{Mist}
项目 & 早期测试 & 本次 \\
\midrule
vLLM-Omni 的图像 token 数 & v0.32 固定为 4900 patch（I2T、I2I 均如此） & main 按输入尺寸计算，512$^2$ 为 1369 patch \\
M* 的图像预处理 & I2T、I2I 均用 \code{vllm}（4900 patch） & I2T 改用默认预处理（1369 patch）；I2T、I2I 均改用默认预处理（1369 patch）；I2I 指令改为 M* 能跟随的指令（第 4 部分） \\
前缀缓存 & M* 开（默认），Omni 关（默认） & 两框架分别报告缓存关、开 \\
workaround & 强制 MATH SDPA，关闭 \code{torch.compile}，数字非默认配置 & 去掉 workaround；关闭 compile 只作诊断 \\
输出尺寸 & M* 的 I2I 输出取输入图尺寸（512$\times$464） & 统一为 512$\times$512，客户端校验 \\
vLLM-Omni 单 stage I2T & 描述与输入图不符，且做了多余的 VAE 编码 & 改用 two-stage 的 Thinker，描述正确，64 token 精确 \\
输出检查 & 只检查 PNG 文件头 & 检查图像内容与 caption 文本 \\
\bottomrule
\end{tabularx}
\vspace{2pt}
\memo{早期 T2I 的 CFG 设置、512$^2$、50 步、warmup 与稳态分开、串行计时，本次沿用。}
\takeaway{早期的 I2T 对比在缓存设置上不对等，I2I 在图像 token 数上不对等。I2T 与 I2I 本次均已等量，输出正确。}
\end{frame}

\begin{frame}[t]{附录 C：M* 默认 SDPA 的 OpenCL 报错（不影响性能）}
\scriptsize\def\codesize{\relax}
\begin{minipage}[t]{.47\textwidth}
\textbf{现象}：M* 默认 worker 的 SDPA 报 \code{OpenCL device not found}，随后 \code{could not create a primitive}。\\[3pt]
\textbf{原因}：用 micromamba 安装 FFmpeg 时带入了 conda 的 \code{libOpenCL.so.1}，\code{.bashrc} 把 \code{ffmpeg/lib} 加入了 \code{LD\_LIBRARY\_PATH}。它遮蔽了系统的 OpenCL ICD loader，而该 loader 只读取一个空的 vendors 目录，因此 \code{clinfo} 看到 0 个平台。系统的 \code{intel-opencl-icd}、L0 driver 与 IGC 均正常。\\[3pt]
\textbf{修复（已验证，无需安装额外组件）}：
\begin{itemize}\scriptsize\setlength\itemsep{1pt}
\item \code{export OCL\_ICD\_VENDORS=/etc/OpenCL/vendors}（已写入 \code{scripts/env.sh}）
\item 或移除 \code{ffmpeg/lib/libOpenCL.so*}，或从 \code{LD\_LIBRARY\_PATH} 中去掉 ffmpeg 目录
\end{itemize}
\textbf{最小复现}：\code{scripts/diagnostics/ocl\_sdpa\_smoke.py}。修复前 ViT 与 VAE 形状的 SDPA 均失败，修复后全部通过。
\end{minipage}\hfill
\begin{minipage}[t]{.5\textwidth}
\textbf{对性能的影响（卡 1+2，早期测试批次，中位数秒）}\\[2pt]
\renewcommand{\arraystretch}{1.1}\setlength{\tabcolsep}{2.5pt}
\begin{tabularx}{\linewidth}{>{\bfseries}p{3.9cm} Y Y}
\toprule
\rowcolor{Mist}
配置 & T2I & I2T 1 token \\
\midrule
M* 强制 MATH SDPA & \val{v032/diag-mstar-mathsdpa/t2i} & \anyval{v032/diag-mstar-mathsdpa/i2t1}\textsuperscript{§} \\
M* 默认 SDPA（已修复） & \val{v032/mstar-tp2/t2i} & \textbf{\anyval{v032/mstar-tp2/i2t1}}\textsuperscript{§} \\
vLLM-Omni 强制 MATH SDPA & \val{v032/diag-omni-single-mathsdpa/t2i} & -- \\
vLLM-Omni 不修复（可运行） & \val{v032/diag-omni-single-no-ocl-fix/t2i} & -- \\
vLLM-Omni 修复后 & \val{v032/omni-single-tp2/t2i} & -- \\
\bottomrule
\end{tabularx}
\vspace{2pt}
{\color{Muted}\scriptsize 带 ViT 的路径上，修复后的默认 SDPA 比 MATH 快约 1.5$\times$；T2I 与 SDPA 后端无关。\textsuperscript{§}两行 M* 的 I2T 均为 M* 默认缓存开，只用于 A/B 比较。vLLM-Omni 不经过受影响的 oneDNN 路径，不修复也能运行，因此该问题只影响 M* 的默认配置。}
\end{minipage}
\takeaway{OpenCL 报错的原因是 conda 的 FFmpeg 自带的 libOpenCL 遮蔽了系统 ICD，设置 \code{OCL\_ICD\_VENDORS} 即可解决。修复本身没有性能代价，与强制 MATH 相比，ViT 路径快约 1.5$\times$。}
\end{frame}

\begin{frame}[t]{附录 C：候选配置：两卡上每个框架的摆法}
\scriptsize\def\codesize{\relax}
\renewcommand{\arraystretch}{1.12}\setlength{\tabcolsep}{3pt}
\begin{tabularx}{\textwidth}{>{\bfseries}p{3.4cm} p{1.0cm} p{3.9cm} Y}
\toprule
\rowcolor{Mist}
配置 & 框架 & 拓扑 & 说明 \\
\midrule
CFG 并行 & \textcolor{Teal}{M*} & LLM 为 TP1，两卡各一份：rank 0 算主分支，rank 1 算引导分支 & \textbf{无 TP 通信}；每卡约 31--32 GiB，显存几乎占满 \\
TP2 & \textcolor{Teal}{M*} & LLM 为 TP2（卡 1、2），ViT 与 VAE 与 rank 0 同卡 & 官方 \code{bagel\_xpu\_tp2} 的两卡版本；每层 all-reduce \\
compile 关 & \textcolor{Teal}{M*} & 同 CFG 并行 & 通过 hook 让 \code{torch.compile} 失效。非官方设置，仅用于诊断 \\
+ XPU graph & \textcolor{Teal}{M*} & 两卡 TP2 / CFG 并行；三卡 TP2 & 两卡 TP2 捕获失败，回退 eager；两卡 CFG 并行显存不足。三卡 TP2 捕获成功（需补丁） \\
\midrule
单 stage，CFG 并行 & \textcolor{Orange}{Omni} & 整个模型放在 DiT stage，\code{cfg\_parallel\_size: 2} & BAGEL 只有 2 路 CFG 分支（T2I）；I2I 需要 3 路，即 3 张卡 \\
单 stage，TP2 & \textcolor{Orange}{Omni} & 整个模型做 TP2 & 两卡下 I2I 的唯一可选方案 \\
two-stage & \textcolor{Orange}{Omni} & Thinker（AR）在卡 1，DiT 在卡 2，KV 经 connector 传输 & I2T 的有效路径。Thinker 单卡 KV 空间小（\code{max\_model\_len} 为 6144） \\
two-stage + Thinker graph & \textcolor{Orange}{Omni} & 同上，Thinker 启用 graph & 捕获成功，收益约 1\% \\
\bottomrule
\end{tabularx}
\takeaway{两个框架在两张卡上各有 2--4 种摆法。拓扑决定通信量，也决定某些 case 能否运行。}
\end{frame}

\begin{frame}[t]{附录 C：候选配置筛选（3 次稳态，T2I 与 64 token）}
\scriptsize\def\codesize{\relax}
\renewcommand{\arraystretch}{1.12}\setlength{\tabcolsep}{3pt}
\begin{tabularx}{\textwidth}{>{\bfseries}p{5.9cm} p{0.9cm} Y Y}
\toprule
\rowcolor{Mist}
配置 & 框架 & T2I 512²·50 步 & I2T 64 token \\
\midrule
M* TP2 · compile 开 · 缓存开 · 无 graph & \textcolor{Teal}{M*} & 12.946\,\textcolor{Muted}{(12.943--12.953)} & \textcolor{Muted}{3.75\textsuperscript{§}} \\
TP2 + XPU graph（捕获失败→回退 eager） & \textcolor{Teal}{M*} & 12.964\,\textcolor{Muted}{(12.960--13.017)} & \textcolor{Muted}{3.79\textsuperscript{§}} \\
M* CFG 并行 · compile 开 · 缓存开 · 无 graph & \textcolor{Teal}{M*} & \cellcolor{GoldBg}\textbf{\textcolor{Teal}{7.702\,\textcolor{Muted}{(7.688--7.707)}}}\,$\star$ & \textcolor{Muted}{2.58\textsuperscript{§}} \\
CFG 并行 + XPU graph（显存耗尽，起不来） & \textcolor{Teal}{M*} & -- & -- \\
\midrule
Omni 单 stage TP2 · 无 graph & \textcolor{Orange}{Omni} & 17.060\,\textcolor{Muted}{(17.048--17.060)} & -- \\
Omni 单 stage CFG 并行 · 无 compile & \textcolor{Orange}{Omni} & \cellcolor{GoldBg}\textbf{\textcolor{Teal}{9.970\,\textcolor{Muted}{(9.949--9.974)}}}\,$\star$ & -- \\
Omni two-stage · 缓存关 · eager & \textcolor{Orange}{Omni} & 17.941\,\textcolor{Muted}{(17.862--17.965)} & 1.881\,\textcolor{Muted}{(1.876--1.881)} \\
Omni two-stage · 缓存关 · Thinker graph & \textcolor{Orange}{Omni} & 17.970\,\textcolor{Muted}{(17.939--17.971)} & \cellcolor{GoldBg}\textbf{\textcolor{Teal}{1.860\,\textcolor{Muted}{(1.856--1.860)}}}\,$\star$ \\
\bottomrule
\end{tabularx}

\vspace{4pt}
\begin{itemize}\scriptsize\setlength\itemsep{1pt}
  \item 中位数（min--max）。$\star$ = 本轮同框架最佳。\textsuperscript{§}本轮 M* 的图生文使用了 4900 patch（与 vLLM-Omni main 不等量），已由等量数据取代。
  \item \textbf{M* TP2 + graph}：捕获失败（\code{Graph nodes cannot depend on events from outside the graph}，LLM 前向含 oneCCL 通信），回退 eager，数字与 TP2 相同。\textbf{M* CFG 并行 + graph}：捕获时显存耗尽，25 分钟内无法启动。
  \item vLLM-Omni two-stage 的 T2I 慢于单 stage：DiT 阶段需要等待 Thinker 经 connector 送来 KV。
\end{itemize}
\takeaway{筛选决定了最终矩阵：M* 取 CFG 并行；vLLM-Omni 文生图取单 stage CFG 并行，图像编辑取单 stage TP2，图生文取 two-stage。}
\end{frame}

\begin{frame}[t]{附录 C：早期结果：两框架均为 TP2，vLLM-Omni v0.32.0}
\scriptsize\def\codesize{\relax}
\renewcommand{\arraystretch}{1.12}\setlength{\tabcolsep}{3pt}
\begin{tabularx}{\textwidth}{>{\bfseries}p{5.9cm} p{0.9cm} Y Y Y Y}
\toprule
\rowcolor{Mist}
配置 & 框架 & T2I 512²·50 步 & I2T 1 token & I2T 64 token & I2I 512²·50 步 \\
\midrule
TP2，compile 开（默认） & \textcolor{Teal}{M*} & \cellcolor{GoldBg}\textbf{\textcolor{Teal}{12.954\,\textcolor{Muted}{(12.944--12.959)}}}\,$\star$ & \textcolor{Muted}{0.96\textsuperscript{§}} & \textcolor{Muted}{3.75\textsuperscript{§}} & \cellcolor{GoldBg}\textbf{\textcolor{Teal}{22.817\,\textcolor{Muted}{(22.802--22.819)}}}\,$\star$ \\
TP2，compile 关（hook） & \textcolor{Teal}{M*} & 13.905\,\textcolor{Muted}{(13.901--13.915)} & \textcolor{Muted}{0.95\textsuperscript{§}} & \textcolor{Muted}{3.04\textsuperscript{§}} & 24.666\,\textcolor{Muted}{(24.662--24.670)} \\
单 stage TP2 & \textcolor{Orange}{Omni} & \cellcolor{GoldBg}\textbf{\textcolor{Teal}{17.656\,\textcolor{Muted}{(17.653--17.659)}}}\,$\star$ & \textcolor{Muted}{描述不正确\textsuperscript{†}} & \textcolor{Muted}{描述不正确\textsuperscript{†}} & \cellcolor{GoldBg}\textbf{\textcolor{Teal}{30.740\,\textcolor{Muted}{(30.736--30.742)}}}\,$\star$ \\
two-stage（Thinker 卡1 + DiT 卡2） & \textcolor{Orange}{Omni} & 18.528\,\textcolor{Muted}{(18.411--18.558)} & \cellcolor{GoldBg}\textbf{\textcolor{Teal}{1.523\,\textcolor{Muted}{(1.517--1.525)}}}\,$\star$ & \cellcolor{GoldBg}\textbf{\textcolor{Teal}{3.347\,\textcolor{Muted}{(3.338--3.349)}}}\,$\star$ & 40.125\,\textcolor{Muted}{(40.114--40.133)} \\
\bottomrule
\end{tabularx}

\vspace{4pt}
\begin{itemize}\scriptsize\setlength\itemsep{1pt}
  \item \textsuperscript{§}M* 的图生文数字为缓存默认开的设置，vLLM-Omni 为缓存默认关，两边不等（图像 token 数相同）。仅供参考。
  \item 这一批两个框架都用 TP2（M* 为默认配置，vLLM-Omni 为单 stage）。T2I 上 M* 快约 27\%。改用 CFG 并行后两边都快 40\% 以上，差距缩小为 23\%。
  \item vLLM-Omni 单 stage 的图生文描述不正确，不列。vLLM-Omni two-stage 在 v0.32.0 的 I2I 输出正常（40.1 s），在 main 上为灰图。
  \item 本批 vLLM-Omni 为 v0.32.0 + vLLM 0.32，与最终矩阵（main + vLLM 0.34）版本不同，不直接比较。版本的影响见因素部分。
\end{itemize}
\takeaway{早期结果只反映 TP2 这一种拓扑，已由最终矩阵取代，保留作为完整记录。}
\end{frame}

\begin{frame}[t]{附录 C：早期图生文结果为什么偏向 M*}
\scriptsize\def\codesize{\relax}
\renewcommand{\arraystretch}{1.15}\setlength{\tabcolsep}{3pt}
\begin{tabularx}{\textwidth}{>{\bfseries}p{3.2cm} Y Y}
\toprule
\rowcolor{Mist}
项目 & 早期测试（vLLM-Omni v0.32.0） & 本次（vLLM-Omni main） \\
\midrule
图像 token & M*：4900 patch（\code{vllm}）；vLLM-Omni：固定 4900 & M*：1369 patch（默认）；vLLM-Omni：按尺寸计算，1369 \\
前缀缓存 & M*：开（默认）；vLLM-Omni：关（默认） & 缓存关、开分别比较 \\
1 token 延迟，早期 & M* TP2：\anyval{v032/mstar-tp2/i2t1}\textsuperscript{§} s；vLLM-Omni two-stage：\raw{v032/omni-two-stage-2gpu/i2t1} s & -- \\
1 token 延迟，本次 & M* CFG 并行：缓存关 \raw{fact/mstar-cfgpar-nocache/i2t1} s，缓存开 \raw{final/mstar-cfgpar/i2t1} s & vLLM-Omni two-stage：缓存关 \raw{final/omni-two-stage/i2t1} s，缓存开 \raw{fact/omni-two-stage-prefixcache/i2t1} s \\
\bottomrule
\end{tabularx}
\vspace{4pt}
\begin{itemize}\scriptsize\setlength\itemsep{2pt}
  \item 早期 M* 的 1 token 延迟更低，原因是设置不一致：M* 的缓存默认开，同一请求重复发送时预填充命中缓存；vLLM-Omni（v0.32）默认关，每次都要重算。两边的图像 token 数相同（都是 4900），因此差异来自缓存。
  \item 切换到 main 后，vLLM-Omni 的图像 token 由 4900 降为 1369（源码确认）。1 token 延迟从 1.52 s 降到 0.19 s；token 数在其中的贡献是推断，未在 v0.32 上用 1369 token 单独测量。
  \item 早期的 M* 为 TP2、vLLM-Omni 为 two-stage，拓扑也不同。本次的 M* 为 CFG 并行。
\end{itemize}
\takeaway{早期 M* 的图生文优势来自缓存设置不一致。设置对齐后，图生文由 vLLM-Omni 更快。}
\end{frame}

\begin{frame}[t,fragile]{附录 D：复现命令与数据目录}
\scriptsize\def\codesize{\relax}
\begin{minipage}[t]{.57\textwidth}
\textbf{运行一个配置}（自动选择空闲卡；卡数不足时不运行）
\begin{Verbatim}[fontsize=\tiny]
cd ~/workspaces/mstar_xpu
scripts/run_one.sh results/<run> mstar-cfgpar mstar \
  configs/mstar/bagel_cfgpar_2gpu.yaml 2 \
  t2i,i2t_1tok,i2t_64tok,i2i
scripts/run_one.sh results/<run> omni-single-cfgpar2 omni \
  configs/vllm_omni/bagel_single_stage_cfgpar2.yaml 2 t2i
scripts/run_one.sh results/<run> omni-two-stage omni \
  configs/vllm_omni/bagel_two_stage_2gpu.yaml 2 \
  i2t_1tok,i2t_64tok --omni-ignore-eos \
  --omni-prompt-style chat
\end{Verbatim}
\textbf{表格与报告}
\begin{Verbatim}[fontsize=\tiny]
scripts/make_tex_tables.py   # 生成报告中的表格与数值
report/build.sh              # 生成表格并编译报告
\end{Verbatim}
\end{minipage}\hfill
\begin{minipage}[t]{.4\textwidth}
\textbf{数据目录}
\begin{itemize}\tiny\setlength\itemsep{0pt}
\item \code{results/20261122\_best\_cfg/}：候选配置与最终矩阵
\item \code{results/20261122\_rerun/}：vLLM-Omni v0.32.0，默认配置
\item \code{results/20261122\_factors/}：connector 与前缀缓存
\item \code{results/20261122\_i2t\_parity/}：M* 图生文输入对齐
\item \code{results/archive/20261118\_initial/}：早期测试
\end{itemize}
\end{minipage}\par
\vspace{3pt}
\memo{失败与被排除的运行保留在 results/ 下的 \texttt{\_failed\_*}、\texttt{\_invalid\_*}、\texttt{\_explore\_*} 目录。每个运行目录包含 timings.csv、manifest.json、样例输出、服务日志、GPU 时间线与生效配置。配置文件位于 configs/mstar 与 configs/vllm\_omni。}
\takeaway{复现时只需两张空闲卡。3 卡布局与卡对 (0,1) 尚未测量。}
\end{frame}


\end{document}
